\documentclass[11pt]{article}
\usepackage[margin=0.85in]{geometry}
\usepackage[T1]{fontenc}
\usepackage[utf8]{inputenc}
\usepackage{lmodern,microtype,amsmath,amssymb,amsthm,mathtools}
\usepackage{booktabs,graphicx,enumitem,natbib,xcolor,tabularx,array}
\usepackage[colorlinks=true,allcolors=blue!55!black]{hyperref}
\setlist{nosep}
\clubpenalty=10000
\widowpenalty=10000
\displaywidowpenalty=10000
\newcolumntype{Y}{>{\raggedright\arraybackslash}X}
\newcolumntype{P}[1]{>{\raggedright\arraybackslash}p{#1}}
\newtheorem{theorem}{Theorem}
\newtheorem{lemma}{Lemma}
\newtheorem{proposition}{Proposition}
\theoremstyle{remark}
\newtheorem{remark}{Remark}
\newcommand{\R}{\mathbb R}
\newcommand{\E}{\mathbb E}
\newcommand{\PP}{\mathbb P}
\newcommand{\ip}[2]{\langle #1,#2\rangle}
\newcommand{\norm}[1]{\lVert #1\rVert}
\newcommand{\rank}{\operatorname{rank}}
\newcommand{\soft}{\operatorname{soft}}
\newcommand{\TV}{\operatorname{TV}}
\newcommand{\Tu}{\mathrm{Tu}}
\newcommand{\TT}{\mathrm{TT}}
\newcommand{\NL}{\mathrm{NL}}
\newcommand{\calL}{\mathcal L}
\emergencystretch=1em
\providecommand{\tightlist}{\setlength{\itemsep}{0pt}\setlength{\parskip}{0pt}}
\title{Technical Supplement: Honest Tensor-Functional Inference}
\author{}\date{Revised September 13, 2026}
\begin{document}\maketitle
This supplement contains the complete supplied proofs supporting the accompanying revised main paper. The main paper states the final central results and empirical findings; this document records the auxiliary comparisons, calibrations, computational contracts and proofs. Fixed-budget intermediate results are retained here as proof ingredients, not as unresolved final rate claims. All formal results have internal status PROVED; this means a checked supplied derivation, not formal verification or external peer review.

The notation is local to each explicit experiment: $D_{\mathrm{amb}}$ is ambient dimension, $\mathsf d$ polynomial degree, $C_{\mathrm{diag}}$ a diagonal-envelope coefficient, $\varepsilon_A$ an approximation allowance and $\eta_k$ a pilot operator bound. The main paper's $q_K=2/K$ is denoted $q$ in scalar-threshold proofs; its $z$-quantile notation (or $z_{\mathrm{val}}$) is used for Gaussian critical values. Canonical complexity $J$ and support overlap $H$ are distinguished.

\tableofcontents
\section{Experiment and parameter classes}\label{sec:model}

Let $K\ge2$, $p_k\ge1$, and $D_{\mathrm{amb}}=\prod_{k=1}^K p_k$. We observe
\begin{equation}\label{eq:model}
 Z=A+\tau G\in\R^{p_1\times\cdots\times p_K},\qquad
 \operatorname{vec}(G)\sim N(0,I_{D_{\mathrm{amb}}}),\qquad \tau>0.
\end{equation}
The noise level is known in the baseline model. Section~\ref{sec:robust_adaptive} states different, explicitly calibrated noise and rank assumptions. Equivalently, $Z$ is the sufficient sample mean of independent full observations $W_j=A+\sigma G_j$, $j=1,\ldots,n$, with known $\sigma$ and $\tau=\sigma/\sqrt n$. This is an exact Gaussian mean equivalence, not an equivalence to scalar-response tensor regression. The observed variables are the tensors $W_j$, or $Z$; $A$ is a fixed unknown parameter. Factor and core representations are nonunique coordinates, not additional latent random variables.

Write $A_{(k)}$ for the mode-$k$ unfolding and $A_{\langle k\rangle}$ for the matrix joining modes $1,\ldots,k$ against $k+1,\ldots,K$. For supplied positive integer caps, define
\begin{align}
 \Theta_{\Tu}(M,\bar r)&=\{A:\norm A_F\le M,\ \rank(A_{(k)})\le\bar r_k,\ k=1,\ldots,K\},\label{eq:tucker}\\
 \Theta_{\TT}(M,\bar r)&=\{A:\norm A_F\le M,\ \rank(A_{\langle k\rangle})\le\bar r_k,\ k=1,\ldots,K-1\}.\label{eq:tt}
\end{align}
The known bound $M$ is positive. Actual ranks may be smaller, including zero; nonzero singular values may be arbitrarily small. No incoherence or coordinate sparsity condition is imposed. The main experiment assumes Gaussian isotropic noise and exact caps. Appendix~\ref{app:refinements} states separately calibrated noise and approximation extensions; missing observations remain untreated.

The prescribed loading has the form
\begin{equation}\label{eq:loading}
 L=\sum_{i=1}^m w_i V_i,\qquad
 V_i=\bigotimes_{k=1}^K v_{ki},\qquad
 v_{ki}^\top v_{kj}=\mathbf1\{i=j\},\qquad \norm w_2=1.
\end{equation}
Here $m\le\min_k p_k$, and each loading family is supplied independently of the data. Thus $\norm L_F=1$. Although this is restrictive, $L$ need not be sparse in the observation basis: known orthogonal mode rotations can make it dense. The estimand and relevant observations are
\begin{equation}\label{eq:diagonal}
 \psi(A)=\ip LA=\sum_iw_i a_i,\qquad
 a_i=\ip A{V_i},\qquad d_i=\ip Z{V_i}=a_i+\tau g_i,
\end{equation}
where the $g_i$ are independent standard normals.

An interval $C(Z)$ is honest at level $1-\alpha$ on $\Theta$ if
\begin{equation}\label{eq:honesty}
 \inf_{A\in\Theta}\PP_A\{\psi(A)\in C(Z)\}\ge1-\alpha.
\end{equation}
This is a finite-sample requirement, not a pointwise limit along a fixed-rank stratum. Length and expected length are the primary criteria. For example,
\begin{equation}\label{eq:minimaxcriteria}
 \calL_0(\Theta)=\inf_{C:\,\eqref{eq:honesty}}\E_0|C|,
 \qquad \calL_{\max}(\Theta)=\inf_{C:\,\eqref{eq:honesty}}\sup_{A\in\Theta}\E_A|C|.
\end{equation}
The infima range over all measurable intervals, allowing independent randomization. The usual squared Frobenius estimation loss has population risk $R_A(B)=\norm{B-A}_F^2+D_{\mathrm{amb}}\sigma^2$ and empirical risk $R_n(B)=n^{-1}\sum_j\norm{W_j-B}_F^2$. We do not minimize these risks or require a full-tensor estimator.

The fixed-budget benchmark comparison in experiment~\eqref{eq:model} uses $p_k=m=p$, fixed $K\ge3$, fixed rank caps, and $M=\tau\sqrt p$. A concrete sequence is $M=\sigma=1$, $n=p$, so $D_{\mathrm{amb}}=n^K$. Section~\ref{sec:norm_loading} removes the relation between $M$ and $\tau$ for the sharp full-data results and extends them to flat orthogonal loadings in unequal ambient dimensions.


The structural and probabilistic assumptions have distinct roles, summarized in Table~\ref{tab:assumptions}. A1--A4 define the inferential problem; A5 is used for the original asymptotic comparisons. The norm/noise relation in A5 is not required in Section~\ref{sec:norm_loading}.
\begin{table}[ht]
\centering\small
\begin{tabularx}{\textwidth}{@{}lYY@{}}
\toprule
 & Assumption & Role and restriction \\
\midrule
A1 & Full observation with independent Gaussian coordinates and known $\tau$ & Gives exact diagonal noise, direct calibration and Gaussian likelihood identities. \\
A2 & Supplied $M$ and Tucker or TT upper rank bounds & Defines the outer honesty class; actual rank and signal strength are unrestricted below it. \\
A3 & Prescribed loading in \eqref{eq:loading} & Gives orthogonal noise coordinates and the diagonal envelope; data-selected loadings are excluded; arbitrary prescribed loadings are treated separately in Section~\ref{sec:robust_adaptive}. \\
A4 & Finite-sample uniform coverage in \eqref{eq:honesty} & Coverage includes zero and local sequences approaching lower-rank strata. \\
A5 & Equal dimensions, fixed $K\ge3$, fixed caps and $M=\tau\sqrt p$ & Produces the stated vanishing length ratio; not needed for the finite-sample interval. \\
\bottomrule
\end{tabularx}
\caption{Assumptions, their purpose and the scope they exclude. Known upper bounds are not replaced by same-data estimates anywhere in the stated procedure.}
\label{tab:assumptions}
\end{table}

\section{Interval construction}\label{sec:method}

Define
\begin{equation}\label{eq:constants}
 q=\frac2K,\quad
 J_{\Tu}=\min_k\prod_{\ell\ne k}\bar r_\ell,\quad
 J_{\TT}=\prod_{k=1}^{K-1}\bar r_k,\quad B=J^{1-q/2}M^q,
\end{equation}
where $J$ is chosen for the applicable class. Set
\begin{equation}\label{eq:threshold}
 \lambda=\tau\sqrt{2\log(2m/\alpha)},\qquad
 \widehat\psi_{\NL}=\sum_iw_i\soft_\lambda(d_i),\qquad
 \soft_\lambda(x)=\operatorname{sign}(x)(|x|-\lambda)_+.
\end{equation}
The three candidate intervals have radii
\begin{align}
 h_{\NL}&=\norm w_\infty B(2\lambda)^{1-q},&
 I_{\NL}&=[\widehat\psi_{\NL}-h_{\NL},\widehat\psi_{\NL}+h_{\NL}],\label{eq:nl}\\
 h_G&=z_{1-\alpha/2}\tau,&
 I_G&=[\ip LZ-h_G,\ip LZ+h_G],\label{eq:direct}\\
 h_0&=M\left(\sum_{i=1}^s w_{(i)}^2\right)^{1/2},&
 I_0&=[-h_0,h_0],\label{eq:zero}
\end{align}
where $|w_{(1)}|\ge\cdots\ge|w_{(m)}|$ and $s$ is the smaller of $m$ and the smallest supplied rank cap. Choose
\begin{equation}\label{eq:select}
 C_*=I_{j_*},\qquad j_*\in\arg\min_{j\in\{\NL,G,0\}}h_j,
\end{equation}
with a fixed rule for ties. This choice depends only on known quantities. Choosing an interval after examining its center or an estimated rank is a different procedure and is not justified by the proof.

The constants in \eqref{eq:constants} calibrate the threshold procedure~\eqref{eq:threshold}; the direct alternative~\eqref{eq:direct} supplies a valid fallback. Only $m$ contractions are used. For coordinate loadings, computation reads $m$ tensor entries and takes $O(m)$ operations. Obtaining contractions for a dense rotated loading has additional cost, which can be large. The procedure does not estimate a Tucker core, TT cores, factors or actual ranks.

\section{Statistical theory}\label{sec:theory}

All numbered theorems and propositions are \textbf{PROVED}, with complete derivations in the appendix. The result-status table separates the proved statements from open computational and adaptive questions. The results use the experiment and classes of Section~\ref{sec:model}. Lower-bound constants are stated for $0<\alpha<1/4$ when this simplifies positivity.

\subsection{Diagonal geometry and honest coverage}

\begin{theorem}[Orthogonal diagonal compressibility]\label{thm:diagonal}
For either class in \eqref{eq:tucker}--\eqref{eq:tt}, under loading \eqref{eq:loading},
\begin{equation}\label{eq:compress}
 \sum_{i=1}^m|\ip A{V_i}|^{2/K}\le J^{1-1/K}M^{2/K}.
\end{equation}
No minimum singular value is required.
\end{theorem}

The proof expands a canonical tensor representation into at most $J$ product terms whose squared component weights sum to $\norm A_F^2$. For TT tensors, these terms need not be orthogonal. A prefix-weight identity establishes the sum exactly; assuming termwise orthogonality would be an incorrect shortcut. H\"older's inequality and Bessel's inequality then give \eqref{eq:compress}. The rank factor may be loose, particularly for growing caps.

For rank cap one the diagonal image is exactly the ball $\{a:\sum_i|a_i|^{2/K}\le M^{2/K}\}$, as shown in Lemma~\ref{lem:image}. Thus the diagonal-only statistical experiment is precisely a classical Gaussian sequence problem on an $\ell_q$ ball. The full tensor experiment contains additional off-diagonal information. This identity permits a sharp benchmark for the restricted observation experiment; it does not remove the additional information in the rest of the tensor.

\begin{theorem}[Finite-sample honest interval]\label{thm:coverage}
For every $0<\alpha<1$, each interval in \eqref{eq:nl}--\eqref{eq:zero} is honest on the applicable cap class. Consequently, $C_*$ is honest and has deterministic length
\begin{equation}\label{eq:length}
 |C_*|=2\min\{h_{\NL},h_G,h_0\}.
\end{equation}
\end{theorem}

The event $\max_i|\tau g_i|\le\lambda$ has probability at least $1-\alpha$. On this event soft thresholding satisfies
\[
 |\soft_\lambda(a_i+\tau g_i)-a_i|\le\min(|a_i|,2\lambda)
 \le |a_i|^q(2\lambda)^{1-q}.
\]
Summing proves the nonlinear interval bound. No central limit theorem, rank recovery or pilot remainder is needed. Known $M$ is essential to this deterministic calibration.

\subsection{Two lower-bound comparisons}

For the remainder of this subsection let all dimensions equal $p$, let $E_i=e_i^{\otimes K}$, and put
\begin{equation}\label{eq:equalL}
 L=p^{-1/2}\sum_{i=1}^p E_i,\qquad z=z_{1-\alpha/2}.
\end{equation}

\begin{theorem}[Fixed-radius affine intervals]\label{thm:affine}
Suppose the honesty class contains $0$ and all $\pm M E_i$. If $0<\alpha<1/2$ and
\[
 C(Z)=[\ip WZ+b-h,\ip WZ+b+h]
\]
is honest, with deterministic $W,b,h$, then
\begin{equation}\label{eq:affinelower}
 h\ge\frac{(M/\sqrt p)z\tau}{M/\sqrt p+z\tau}.
\end{equation}
On the rank-one tensor Frobenius ball, the best radius in this affine class is therefore of order $\min(M/\sqrt p,\tau)$, up to constants depending on $\alpha$.
\end{theorem}

The theorem allows $W$ to use the entire tensor. It does not cover affine centers with data-dependent radii. The matching upper order follows by choosing between $I_G$ and $I_0$ for rank cap one.

For a nonempty proper subset $S$ of modes, write $A_{[S]}$ for the corresponding bipartition unfolding and define
\begin{equation}\label{eq:relax}
 \Theta_S=\{A:\rank(A_{[S]})\le1,\ \norm A_F\le M\}.
\end{equation}

\begin{theorem}[Cost of a single unfolding honesty class]\label{thm:unfolding}
If $0<\alpha<1/4$, every interval honest on $\Theta_S$ satisfies
\begin{equation}\label{eq:unfoldlower}
 \E_0|C|\ge c_\alpha\min\left\{\frac M{\sqrt p},\frac\tau4\right\},\qquad
 c_\alpha=1-2\alpha-\frac12\sqrt{\sqrt{8/7}-1}>0.
\end{equation}
The lower order $\min(M/\sqrt p,\tau)$ is attained up to $\alpha$-dependent constants by the direct/structural choice on $\Theta_S$.
\end{theorem}

The alternatives in the proof are rank one in the chosen matrix unfolding. They generally violate the other tensor rank constraints. This is necessary for the comparison: the theorem measures the cost of requiring coverage on the larger class. For Tucker caps, a mode unfolding is a permissible relaxation; for TT caps, a contiguous cut is permissible. With all tensor rank caps one, every bipartition in \eqref{eq:relax} contains the original tensor class. Increasing an unfolding cap further preserves the lower bound.

\begin{theorem}[Strict length separations]\label{thm:separation}
Let $0<\alpha<1/4$, $K\ge3$ and the rank caps be fixed. Under \eqref{eq:equalL} and $M=\tau\sqrt p$, $C_*$ is honest on the Tucker or TT class and
\begin{equation}\label{eq:separation}
 |C_*|\le2\tau J^{1-1/K}
 \left\{\frac{8\log(2p/\alpha)}p\right\}^{(K-2)/(2K)}=o(\tau).
\end{equation}
In contrast, every fixed-radius affine interval honest on the same class has length at least $2z\tau/(1+z)$. Every interval honest on a class retaining only one constrained unfolding has expected length at zero at least $c_\alpha\tau/4$.
\end{theorem}

For growing caps, the same conclusion follows if $J^{1-1/K}(\log p/p)^{(K-2)/(2K)}\to0$. This condition may be restrictive. At $K=2$ the exponent in \eqref{eq:separation} is zero, so the argument gives no polynomial separation. Theorem~\ref{thm:separation} neither asserts that the displayed exponent is optimal nor compares favorably at every strong signal with an efficient oracle projection interval. The following benchmark concerns the diagonal observation experiment; the subsequent full-data theorem gives the final rate.

\subsection{An initial full-data lower bound}\label{sec:gap}

\begin{theorem}[Within-class lower bound]\label{thm:jointlower}
Under \eqref{eq:model} and \eqref{eq:equalL}, for $0<\alpha<1/4$, every interval honest on a class containing the rank-one tensor Frobenius ball of radius $M$ satisfies
\begin{equation}\label{eq:jointlower}
 \E_0|C|\ge(3/4-2\alpha)
 \frac{\min\{M,\tau\sqrt{\log(1+p/4)}\}}{\sqrt p}.
\end{equation}
\end{theorem}

This is a sparse-coordinate mixture bound. At $M=\tau\sqrt p$, for fixed caps and fixed $K\ge3$, the results imply
\begin{equation}\label{eq:gap}
 c\tau\sqrt{\frac{\log p}p}
 \le\calL_0(\Theta)\le\calL_{\max}(\Theta)
 \le C\tau\left(\frac{\log p}p\right)^{(K-2)/(2K)}
\end{equation}
for sufficiently large $p$, with constants depending on $K$, the caps and $\alpha$. The ratio of these initial scales grows as $(p/\log p)^{1/K}$. Section~\ref{sec:sharp_full} closes the gaps: the two criteria in \eqref{eq:minimaxcriteria} have the same rate, with an additional logarithmic factor in the cubic case.

\subsection{Optimality when only the diagonal is observed}\label{sec:diagonal_optimality}

Let $\Theta_1$ denote the Frobenius ball of tensors with rank at most one, including zero. Restrict the interval to be measurable in $d=(Z_{i,\ldots,i})_{i=1}^p$ and an independent random seed. Write $\mathcal C_{\alpha,d}(\Theta_1)$ for such intervals that are honest on $\Theta_1$, and define
\[
 \mathcal L_{0,d}(\Theta_1)=\inf_{C\in\mathcal C_{\alpha,d}(\Theta_1)}\E_0|C|,
 \qquad
 \mathcal L_{\max,d}(\Theta_1)
 =\inf_{C\in\mathcal C_{\alpha,d}(\Theta_1)}\sup_{A\in\Theta_1}\E_A|C|.
\]
The next result specializes classical sparse-sequence arguments \citep{cailownonconvex2004,collier2017} to honest intervals. It is a benchmark for the proposed use of diagonal information, not a new Gaussian-sequence minimax principle.

\begin{proposition}[Classical diagonal-information benchmark]\label{prop:diagonal_minimax}
Assume \eqref{eq:model}, \eqref{eq:equalL} and $0<\alpha<1/4$. For each integer $1\le s\le p$, put
\[
 b_s=\min\left\{M s^{-K/2},\
       \tau\sqrt{\log\left(1+\frac{p}{4s^2}\right)}\right\},
 \qquad
 c_{\alpha,d}=1-2\alpha-\frac12\sqrt{e^{1/4}-1}>0.
\]
Every $C\in\mathcal C_{\alpha,d}(\Theta_1)$ satisfies
\begin{equation}\label{eq:diagonal_lower}
 \E_0|C|\ge c_{\alpha,d}\frac{s b_s}{\sqrt p}.
\end{equation}
For fixed $K\ge3$ and $M=\tau\sqrt p$,
\begin{equation}\label{eq:diagonal_rate}
 \mathcal L_{0,d}(\Theta_1)\asymp
 \mathcal L_{\max,d}(\Theta_1)\asymp
 \tau\left(\frac{\log p}{p}\right)^{(K-2)/(2K)}.
\end{equation}
The comparison constants depend on $K$ and $\alpha$.
\end{proposition}

The rate in \eqref{eq:diagonal_rate} follows because the effective support size $s\asymp(p/\log p)^{1/K}$ makes \eqref{eq:diagonal_lower} match the threshold interval's upper order. The same asymptotic conclusion holds for any fixed Tucker or TT caps: the class contains $\Theta_1$ and Theorem~\ref{thm:coverage} changes the upper bound only by a fixed rank factor. Thus further improvement in the power of $p$ cannot come from recalibrating a procedure restricted to $d$.

For comparison with the full-data result, consider the preliminary bound in \eqref{eq:gap}. For the common-target tensors $A_S=b_s(\sum_{i\in S}e_i)^{\otimes K}$, two support sets with overlap $H$ have diagonal inner product $b_s^2H$ but full-tensor inner product $b_s^2H^K$. The likelihood second moment in the full experiment is consequently different. A valid diagonal lower bound cannot be transferred by ignoring the extra observations. Theorem~\ref{thm:scan} establishes that they do permit shorter expected length at zero.

Table~\ref{tab:theory_comparison} separates the procedure and parameter classes in the comparisons.
\begin{table}[ht]
\centering\small
\begin{tabularx}{\textwidth}{@{}P{0.24\textwidth}P{0.22\textwidth}Y@{}}
\toprule
Comparison & Information and honesty class & Length statement when $M=\tau\sqrt p$ \\
\midrule
Direct Gaussian interval & All means; prescribed $L$ & $2z\tau$, with exact coverage. \\
Selected interval & Joint cap class; uses $d$ & At most $2\tau J^{1-1/K}\{8\log(2p/\alpha)/p\}^{\beta}$; $\beta=(K-2)/(2K)$. \\
Any fixed-radius affine interval & Full tensor; same joint class & Length at least $2z\tau/(1+z)$. \\
Any interval on one unfolding class & Full tensor; enlarged honesty class & Expected length at zero at least $c_\alpha\tau/4$. \\
Any diagonal-based interval & Only $d$; rank-one or fixed-cap class & Optimal expected length at zero and in worst case has order $\tau(\log p/p)^\beta$. \\
Any full-data interval & Full tensor; joint class & Optimal expected length both at zero and in the worst case is $r_K(p)$ in \eqref{eq:full_rates}. \\
\bottomrule
\end{tabularx}
\caption{Precisely delimited comparisons for fixed $K\ge3$ and fixed rank caps. Lower bounds for relaxed honesty classes or restricted observations are not lower bounds for unrestricted full-data procedures on the original class.}
\label{tab:theory_comparison}
\end{table}




\subsection{Using the full tensor at the zero mean}\label{sec:full_scan}

The diagonal rate is a constraint on procedures using the diagonal. A finite scan over product directions of coordinate subtensors can give shorter expected length at zero while preserving the original uniform length bound. The scan is an existence construction with exponential complexity; no efficient implementation is asserted.

For a tensor $B$, define its injective norm
\[
 \rho(B)=\sup_{\norm{u_k}_2=1}
       |\ip B{u_1\otimes\cdots\otimes u_K}|.
\]
For every nonempty subset $S\subseteq\{1,\ldots,p\}$ of size $s$, choose fixed $1/(2K)$-nets of the unit sphere in $\R^s$, one per mode, each with at most $(1+4K)^s$ points. Let $N_S(Z)$ be the largest absolute contraction of $Z_{S,\ldots,S}$ over their product. Define
\begin{align}
 \lambda_s^{\rm scan}
 &=\tau\left[2\{s\log(ep/s)+Ks\log(1+4K)+\log(2p/\gamma)\}\right]^{1/2},
 \label{eq:scan_threshold}\\
 R_\gamma
 &=\frac{4J^{K/2}}{\sqrt p}
       \sum_{s=1}^p\lambda_s^{\rm scan}s^{-K/2}.
 \label{eq:scan_radius}
\end{align}
Accept the zero test if $N_S(Z)\le\lambda_{|S|}^{\rm scan}$ for every nonempty $S$. On acceptance, let $C_{\rm scan}=[-h_{\rm scan},h_{\rm scan}]$, where $h_{\rm scan}=\min(h_0,R_\gamma)$; on rejection let $C_{\rm scan}=[-h_0,h_0]$. Here $h_0=M\sqrt{s_{\rm cap}/p}$ and $s_{\rm cap}$ is the smallest cap, truncated at $p$. For $0<\alpha_0,\gamma<1$ with $\alpha_0+\gamma\le\alpha$, intersect this interval with $C_*(\alpha_0)$ and call the result $C_\dagger$; replace an empty intersection by $\{0\}$. Write $\ell_*=|C_*(\alpha_0)|$.

\begin{theorem}[A full-data interval with short expected length at zero]\label{thm:scan}
Under \eqref{eq:model}, \eqref{eq:equalL} and either cap class, $C_{\rm scan}$ is honest at level $1-\gamma$, and $C_\dagger$ is honest at level $1-\alpha$. Moreover,
\begin{equation}\label{eq:scan_expectation}
 \sup_{A\in\Theta}\E_A|C_\dagger|\le\ell_*,
 \qquad
 \E_0|C_\dagger|\le
 \min\{\ell_*,\,2R_\gamma+\ell_*\gamma\}.
\end{equation}
For fixed $0<\alpha<1/4$, fixed $K\ge3$, fixed positive caps and $M=\tau\sqrt p$, choose $\alpha_0=\alpha/2$ and $\gamma=\alpha/(2p)$. Then
\begin{equation}\label{eq:scan_rate}
 \E_0|C_\dagger|
 \le C_{K,J,\alpha}\frac{\tau\sqrt{\log p}}{\sqrt p}
       \sum_{s=1}^p s^{-(K-1)/2}
\end{equation}
for sufficiently large $p$. Consequently,
\begin{equation}\label{eq:zero_sharp}
 \calL_0(\Theta)\asymp\tau\sqrt{\frac{\log p}{p}},
 \qquad K\ge4.
\end{equation}
\end{theorem}

The matching lower bound in \eqref{eq:zero_sharp} is Theorem~\ref{thm:jointlower}, applied to the same full-data experiment and honesty class. The sum in \eqref{eq:scan_rate} is bounded for $K\ge4$ and harmonic for $K=3$. The number of scanned product directions is bounded by $(1+(1+4K)^K)^p-1$. Thus the result establishes attainable statistical precision, without a polynomial-time guarantee. The worst-case bound in \eqref{eq:scan_expectation} is the initial diagonal upper order. The residual projection in Theorem~\ref{thm:full_minimax} strengthens it to the sharp full-data rate.

\begin{theorem}[A harmonic obstruction in order three]\label{thm:harmonic}
Under \eqref{eq:model} and \eqref{eq:equalL}, let $K=3$, $M=\tau\sqrt p$ and $0<\alpha<1/4$. Every interval honest on a class containing the rank-one tensor Frobenius ball satisfies
\begin{equation}\label{eq:harmonic_lower}
 \E_0|C|\ge
 (3/4-2\alpha)\frac{\tau H_p}{200\sqrt p},
 \qquad H_p=\sum_{i=1}^p i^{-1}.
\end{equation}
\end{theorem}

Theorem~\ref{thm:harmonic} uses a full-data mixture of rank-one tensors with sign-coupled factors and a common harmonic diagonal. It refutes the earlier conjecture that $\calL_0$ always has order $\tau\sqrt{\log p/p}$ at rank one. Combining \eqref{eq:harmonic_lower} and \eqref{eq:scan_rate} gives the intermediate cubic bracket
\begin{equation}\label{eq:cubic_gap}
 c_\alpha^{\rm cub}\frac{\tau\log p}{\sqrt p}
 \le\calL_0(\Theta)
 \le C_{J,\alpha}\frac{\tau(\log p)^{3/2}}{\sqrt p}.
\end{equation}
Here $c_\alpha^{\rm cub}=(3/4-2\alpha)/200>0$. The intermediate bracket \eqref{eq:cubic_gap} has a factor $\sqrt{\log p}$ gap. Theorem~\ref{thm:sparse_harmonic} closes it from below. The constant $1/200$ is conservative and is not an optimal finite-sample constant.

\paragraph{Further verified refinements.}
The accompanying theory review proves a smaller rank coefficient in the diagonal bound, an exact one-dimensional characterization of the optimal fixed affine radius, and a noise-envelope version with an additive approximation term. It also records finite-variance and unknown-Gaussian-variance extensions requiring retained replicates. These supporting results are collected in Appendix~\ref{app:refinements}; the simulations retain the original Gaussian threshold procedure.


\subsection{Sharp full-data rates, including worst-case length}\label{sec:sharp_full}

A sparse harmonic profile strengthens the cubic lower bound. Inverting the residual scan then turns the zero-point construction into an interval with uniformly controlled length.

\begin{theorem}[Sparse harmonic lower bound]\label{thm:sparse_harmonic}
Under \eqref{eq:model} and \eqref{eq:equalL}, let $K=3$ and $0<\alpha<1/4$.
For every integer $1\le m\le p$, put
\[
 b_m=\min\left\{\frac{M}{\sqrt{27m}},
       \frac{\tau}{\sqrt{27}}\sqrt{\log(1+p/(4m^2))}\right\},
 \qquad c_{\alpha,d}=1-2\alpha-\tfrac12\sqrt{e^{1/4}-1}.
\]
Every interval honest on a class containing the rank-one Frobenius ball satisfies
\begin{equation}\label{eq:sparse_harmonic}
 \E_0|C|\ge c_{\alpha,d}\frac{b_mH_m}{\sqrt p}.
\end{equation}
If $M=\tau\sqrt p$, then for $p\ge256$ and $m=\lfloor p^{1/4}\rfloor$,
\[
 \E_0|C|\ge\frac{c_{\alpha,d}}{16\sqrt{27}}\,
                  \frac{\tau(\log p)^{3/2}}{\sqrt p}.
\]
\end{theorem}

The profile weights, rather than additional rank components, produce the harmonic sum. The alternatives are nonnegative rank-one tensors. Their full likelihood overlap is bounded linearly in the support intersection, permitting the classical hypergeometric calculation.

Write $\lambda_s(\gamma)=\lambda_s^{\rm scan}$ for the threshold in
\eqref{eq:scan_threshold} at error budget $\gamma$.
Use the product nets above with thresholds $\lambda_s(\alpha)$. Define
\[
 {\cal F}(Z)=\{B\in\Theta(M,\bar r):
        N_S(Z-B)\le\lambda_{|S|}(\alpha)\text{ for every nonempty }S\}.
\]
Let $C_F$ be the interval between the minimum and maximum of $\psi$ on
$\mathcal F(Z)$, and set $C_F=\{0\}$ when this set is empty.
Write $J_2=J(2\bar r)$, the expansion constant with doubled caps, and put
\[
 W_\alpha=\frac{4J_2^{K/2}}{\sqrt p}
           \sum_{s=1}^p\lambda_s(\alpha)s^{-K/2}.
\]

\begin{theorem}[Uniform diameter and full-data minimax length]\label{thm:full_minimax}
Under \eqref{eq:model}, \eqref{eq:equalL}, $0<\alpha<1$ and either positive cap class,
$C_F$ is measurable and honest at level $1-\alpha$, and for every observation
\begin{equation}\label{eq:diameter}
 |C_F|\le\min(2h_0,W_\alpha).
\end{equation}
For fixed $0<\alpha<1/4$, fixed $K\ge3$, fixed positive caps and $M=\tau\sqrt p$,
\begin{equation}\label{eq:full_rates}
 \calL_0(\Theta)\asymp\calL_{\max}(\Theta)\asymp r_K(p),
 \qquad
 r_K(p)=
 \begin{cases}
 \tau(\log p)^{3/2}/\sqrt p,&K=3,\\
 \tau\sqrt{\log p/p},&K\ge4.
 \end{cases}
\end{equation}
Constants may depend on $K$, the caps and $\alpha$.
\end{theorem}

Differences of feasible tensors obey doubled unfolding caps; the proof does not assume that a rank class is closed under subtraction. Computing the confidence-set extrema is not claimed to be efficient.

\begin{proposition}[Squared-error consequence]\label{prop:mse_sharp}
In the regime of \eqref{eq:full_rates},
\[
 \inf_{\widehat\psi}\sup_{A\in\Theta}
          \E_A(\widehat\psi-\psi(A))^2\asymp r_K(p)^2.
\]
The midpoint of $C_F$ with error level $\alpha_p=p^{-2}$ attains the upper order for $p\ge2$.
\end{proposition}

\subsection{Computation: a certified hierarchy and a restricted lower bound}\label{sec:computation}

The exact residual projection supplies the statistical benchmark. A truncated zero test is computable with an explicit search budget. It does not compute the extrema defining $C_F$.

\begin{proposition}[Truncated scans]\label{prop:truncated}
For $1\le t\le p$, test zero using only supports of sizes at most $t$, retaining thresholds calibrated simultaneously over all sizes. Write $\lambda_s=\lambda_s(\gamma)$ and put
\[
 q=2/K,\quad B_q=J^{1-1/K}M^q,\quad
 a_t=4J^{K/2}\lambda_t t^{-K/2},
\]
\[
 R_t=\frac{4J^{K/2}}{\sqrt p}\sum_{s=1}^t\lambda_s s^{-K/2}
       +\mathbf1_{\{t<p\}}\frac{B_q}{\sqrt p}a_t^{1-q}.
\]
On acceptance use radius $\min(h_0,R_t)$ centered at zero; on rejection use the structural interval.
Intersect with $C_*(\alpha_0)$, where $0<\alpha_0,\gamma<1$,
$\alpha_0+\gamma\le\alpha$, and $\gamma$ is the scan error budget, using the singleton convention.
The resulting $C_t$ is honest, has length at most $\ell_*=|C_*(\alpha_0)|$, and
\[
 \E_0|C_t|\le\min(\ell_*,2R_t+\ell_*\gamma).
\]
For fixed $K\ge3$, fixed caps and fixed $0<\alpha<1$, with $M=\tau\sqrt p$, $\alpha_0=\alpha/2$ and $\gamma=\alpha/(2p)$,
\begin{align}\label{eq:trunc_rate}
 R_t\lesssim{}&
 \frac{\tau\sqrt{\log p}}{\sqrt p}
       \sum_{s=1}^t s^{-(K-1)/2}\\
 &+\mathbf1_{\{t<p\}}\tau(\log p/p)^\beta
                t^{-(K-1)\beta},\qquad
 \beta=(K-2)/(2K).\nonumber
\end{align}
\end{proposition}

Explicit normalized grid nets have at most $(3+4K\sqrt s)^s$ directions per mode.
Using their cardinalities in the Gaussian thresholds preserves \eqref{eq:trunc_rate}.
The arithmetic cost through level $t$ is bounded by
\[
 O\left\{\sum_{s=1}^t\binom ps
             (3+4K\sqrt s)^{Ks}s^K\right\}.
\]
For fixed $t,K$ this is polynomial in $p$. Fixed $t$ does not close the power gap in the tail term of \eqref{eq:trunc_rate}. Completed levels share one noise event, so successive intersections remain honest even if the decision to compute another complete level uses the data.

A complementary obstruction holds for polynomial estimators under squared-error loss. This is the restricted low-degree framework used in sparse tensor PCA \citep{choo2021,fu2026}; it does not characterize all polynomial-time computation.

\begin{theorem}[A low-degree estimation separation]\label{thm:low_degree}
Let $K\ge3$, $1\le s\le\sqrt p$, and mix uniformly over
\[
 x_S=s^{-1/2}\mathbf1_S,\qquad
 A_S=\tau\lambda x_S^{\otimes K},\qquad |S|=s.
\]
Assume $0<\tau\lambda\le M$, and write
\[
 \delta=\frac{\tau\lambda s^{1-K/2}}{\sqrt p},\quad
 \mu=s^2/p,\quad B_{\mathsf d}=\frac{\lambda^2}{s^K}(K\mathsf d)^K,
\]
where $\mathsf d\ge1$ is the polynomial degree. If $B_{\mathsf d}<1$, the centered likelihood ratio projected onto polynomials of degree at most $\mathsf d$ has squared null norm at most
\[
 \eta_{\mathsf d}^2=\mu e^\mu B_{\mathsf d}/(1-B_{\mathsf d}).
\]
Every scalar polynomial estimator $T(Z)$ of degree at most $\mathsf d$ satisfies
\begin{equation}\label{eq:lowdegree_risk}
 \sup_{A\in\Theta}\E_A(T-\psi(A))^2
       \ge\frac{\delta^2}{(2+\eta_{\mathsf d})^2}
\end{equation}
on any parameter class containing zero and these alternatives.
In particular, let $M=\tau\sqrt p$, $\lambda=\sqrt p$,
$s=\lfloor p^a\rfloor$, and $\mathsf d=\lfloor p^b\rfloor$, where
\[
 \frac1K<a<\min\left(\frac12,\frac1{K-2}\right),
 \qquad 0<b<a-\frac1K.
\]
Then $\eta_{\mathsf d}\to0$ and the lower order in \eqref{eq:lowdegree_risk} is polynomially larger than the unrestricted rate in Proposition~\ref{prop:mse_sharp}.
\end{theorem}

This theorem includes $K=3$ and uses nonnegative sparse factors. It does not imply the same limitation for arbitrary interval endpoints or arbitrary thresholded algorithms.

\begin{proposition}[What an efficient short interval would imply]\label{prop:ci_detection}
Let $Q$ be supported on parameters with common target $\delta>0$, and let the null target be zero.
An interval honest at level $1-\alpha$ at the null and these alternatives defines
$\phi(Z)=\mathbf1\{\delta\in C(Z)\}$ with
\[
 Q(\phi=0)\le\alpha,\qquad
 P_0(\phi=1)\le\alpha+\E_0|C|/\delta.
\]
A computational hypothesis excluding polynomial-time tests with total error below $c>2\alpha$ therefore excludes polynomial-time honest intervals with
$\E_0|C|<(c-2\alpha)\delta$ for that prior.
\end{proposition}

The last implication is conditional on a testing-hardness hypothesis.
Theorem~\ref{thm:low_degree} alone does not establish that hypothesis.
Efficient minimax confidence intervals remain open, with a proved restricted estimation obstruction and an explicit computable hierarchy.

\section{Norm budgets and loading geometry}\label{sec:norm_loading}

The restriction $M=\tau\sqrt p$ can be removed for the equal diagonal contrast. The key is to impose the Frobenius budget and the residual-scan constraints on each ordered diagonal entry, before summing. Throughout this section the tensor order and positive rank caps are fixed; constants are not claimed uniform in growing caps.

Put $P_2=2J_2^{(K-1)/2}$ and $Q_2=4J_2^{K/2}$, where $J_2=J(2\bar r)$, and define
\begin{equation}\label{eq:norm_diameter}
 U_{K,p}(M;\alpha)=\frac1{\sqrt p}\sum_{s=1}^p
       \min\{P_2M,Q_2\lambda_s(\alpha)\}\,s^{-K/2}.
\end{equation}
Here $\lambda_s(\alpha)$ is the finite-net threshold in \eqref{eq:scan_threshold}, with error budget $\alpha$.

\begin{theorem}[Diameter using both norm and scan constraints]\label{thm:norm_diameter}
For every $p\ge1$, $K\ge3$, $M,\tau>0$ and $0<\alpha<1$, the residual-projection interval $C_F$ is honest and satisfies, for every observation,
\[
 |C_F|\le\min\{2h_0,U_{K,p}(M;\alpha)\}.
\]
The bound retains the supplied norm budget inside the ordered-coordinate sum; taking the minimum only after summation can lose the intermediate cubic regime.
\end{theorem}

Let $\ell_p=\log(ep)$, $x=M/(\tau\sqrt{\ell_p})$, and define
\begin{align}
 g_p(x)&=\min\{x,1+\log_+x,\ell_p\},\qquad
          \log_+x=\max(0,\log x),\label{eq:norm_g}\\
 \mathfrak r_K(p,M,\tau)&=\frac{\tau\sqrt{\ell_p}}{\sqrt p}
 \begin{cases}
 g_p(x),&K=3,\\
 \min(x,1),&K\ge4.
 \end{cases}\label{eq:norm_rate}
\end{align}
These definitions also cover $p=1$, when the rate is $\min(M,\tau)$.

\begin{theorem}[Full-data length for every positive norm budget]\label{thm:all_norm_rates}
Fix $K\ge3$, positive Tucker or TT caps, and $0<\alpha<1/4$. For the equal diagonal contrast there are constants $0<c<C<\infty$, depending only on these fixed quantities, such that for every $p\ge1$ and every $M,\tau>0$,
\[
 c\,\mathfrak r_K(p,M,\tau)
 \le\mathcal L_0(\Theta)
 \le\mathcal L_{\max}(\Theta)
 \le C\,\mathfrak r_K(p,M,\tau).
\]
The lower bounds already hold when the honesty class is restricted to rank-one tensors. The interval in Theorem~\ref{thm:norm_diameter} attains the upper order.
\end{theorem}

For $K\ge4$ the rate is $\min\{M,\tau\sqrt{\log(ep)}\}/\sqrt p$. For $K=3$, the norm-limited regime $x\le1$ has rate $M/\sqrt p$; above it the multiplier grows as $1+\log x$ until capped at $\ell_p$. These are rates up to fixed multiplicative constants, not exact transition constants. Substituting $M=\tau\sqrt p$ recovers Theorem~\ref{thm:full_minimax}.

For squared-error estimation, a fixed-confidence midpoint argument is insufficient when $M/\tau$ is arbitrarily large: its off-coverage term can grow with $M^2$. Instead use an approximate minimum-residual estimator. Fix $\beta=1/8$, and put
\[
 T(R)=\max\bigl(0,\max_{S\ne\varnothing}
                  \{N_S(R)-\lambda_{|S|}(\beta)\}\bigr).
\]
Choose a fixed countable dense sequence $(B_j)$ in the compact cap class. Define
\begin{equation}\label{eq:residual_fit}
 d(Z)=\inf_j T(Z-B_j),\quad
 j(Z)=\min\{j:T(Z-B_j)\le d(Z)+\tau\},\quad
 \widehat\psi=\psi(B_{j(Z)}).
\end{equation}
The tolerance $\tau$ ensures the displayed index exists. This defines a statistical estimator, not a claimed efficient optimization algorithm.

\begin{theorem}[Squared risk uniformly over norm budgets]\label{thm:all_norm_mse}
Under the fixed-order, fixed-cap conditions of Theorem~\ref{thm:all_norm_rates},
\[
 \inf_{\widetilde\psi}\sup_{A\in\Theta}
       \E_A(\widetilde\psi-\psi(A))^2
       \asymp\mathfrak r_K(p,M,\tau)^2
\]
uniformly over $p\ge1$ and $M,\tau>0$. The measurable estimator in \eqref{eq:residual_fit} attains the upper bound. It uses neither a failure probability depending on $M$ nor a minimum nonzero singular value.
\end{theorem}

\subsection{Nonuniform orthogonal loadings}

Return to the prescribed orthogonal product loading in \eqref{eq:loading}, with $m$ nonzero weights and possibly unequal ambient mode dimensions. Let $w_s^*$ denote decreasing absolute weights. Use orthogonal mode transformations and restrict to the $m^K$ coordinates spanned by the prescribed loading vectors. The resulting Gaussian tensor retains the rank and norm bounds. Write $\lambda_s^{(m)}$ for the scan threshold with $p=m$, and define
\begin{align}
 U_w&=\sum_{s=1}^m w_s^*
          \min\{P_2M,Q_2\lambda_s^{(m)}(\alpha)\}s^{-K/2},\label{eq:weighted_upper}\\
 L_w&=\max_{1\le d\le m} w_d^*\sqrt d\,
                  \mathfrak r_K(d,M,\tau).\label{eq:weighted_lower}
\end{align}
The rank constants use the same supplied caps. The structural radius is
\[
h_0=M\left(\sum_{s=1}^{s_{\rm cap}}(w_s^*)^2\right)^{1/2},
\]
with $s_{\rm cap}$ truncated at $m$.

\begin{theorem}[Loading-dependent full-data bounds]\label{thm:weighted_bounds}
For fixed $K\ge3$, fixed positive caps and fixed $0<\alpha<1/4$, an honest interval obtained by residual projection in the rotated block satisfies
\[
 |C_w(Z)|\le\min(2h_0,U_w)
\]
for every observation. Conversely every interval honest on the full ambient cap class satisfies
\[
 \E_0|C(Z)|\ge c_{K,\alpha}L_w.
\]
The lower bound allows use of every ambient tensor entry. These bounds are not asserted to match for arbitrary nonuniform weights.
\end{theorem}

\begin{proposition}[Sharp rates for a flat orthogonal loading]\label{prop:flat_loading}
If $w_i=\epsilon_i/\sqrt m$ with prescribed $\epsilon_i\in\{-1,1\}$, then both minimax expected-length criteria have order $\mathfrak r_K(m,M,\tau)$, and minimax squared error has its square, independently of the larger ambient mode dimensions. The statement includes dense orthogonal rotations of the loading. At $m=1$ it gives the bounded-normal-mean scale $\min(M,\tau)$.
\end{proposition}

\paragraph{Scope of the general-loading bounds.}
The displayed envelopes need not match for arbitrary unequal weights. Results L1 and W2 below establish additional small-budget and fixed-decay regimes. Results A1, G2 and P1 separately describe fixed-cap rate adaptation, a growing-outer-class obstruction and strong saturated pointwise inference. General loading profiles and complete pointwise adaptation remain open.

\section{Robust noise, computation, and rank adaptation}\label{sec:robust_adaptive}

This section separates extensions of validity from sharp minimax claims. It also separates adaptation on a spectrally separated class from honesty on the original class containing arbitrarily weak components. Lower-bound comparisons use a fixed error level below $1/4$. All statements have proofs in Appendix~\ref{app:robust_adaptive}.

\subsection{Beyond Gaussian observations}

For a centered error tensor $E$, the following joint directional condition is sufficient:
\begin{equation}\label{eq:joint_sg}
 \E\exp(t\ip UE)\le\exp(t^2v^2\norm U_F^2/2)
 \quad(U\text{ deterministic},\ t\in\R).
\end{equation}
It allows dependence and nonidentical coordinates. Independent centered sub-Gaussian coordinates with common proxy $v$ satisfy it, as does Gaussian covariance bounded above by $v^2I$. Marginal coordinate bounds alone do not imply it.

\begin{theorem}[Full-data robustness under directional tails]\label{thm:sg_full}
In $Z=A+E$, replace $\tau$ by a known $v>0$ in the scan thresholds and residual-fit tolerance. Under \eqref{eq:joint_sg}, the norm-sensitive diameter, uniform squared-risk upper bound, and orthogonal-loading upper bounds remain valid. For fixed order/caps, their equal-diagonal orders are $\mathfrak r_K(p,M,v)$ and its square. If the noise law is an unknown member of the family satisfying \eqref{eq:joint_sg}, these rates are minimax when coverage and the supremum risk/length range over that family. This lower-bound statement need not hold for each fixed non-Gaussian law separately.
\end{theorem}

For a finite-variance extension retain independent identically distributed full-tensor replicates $W_1,\ldots,W_n$, with mean $A$ and
$\operatorname{Cov}(\operatorname{vec}W_j)\preceq\sigma^2I$ for known $\sigma>0$.
Put $\tau=\sigma/\sqrt n$, and let
\[
 H_s=s\log(ep/s)+Ks\log(1+4K).
\]
For a desired error $0<\beta<1/4$, choose $b_s$ as the smallest odd integer at least
\begin{equation}\label{eq:mom_blocks}
 \max\{9,[H_s+\log(p/\beta)]/\log2\},\qquad n\ge2\max_s b_s.
\end{equation}
For each unit product-net direction $U$ supported on $S^K$, $|S|=s$, let $d_U$ be the median of its $b_s$ disjoint block averages, using block size $k_s=\lfloor n/b_s\rfloor$. The block partitions may differ across $s$ but are deterministic. Set $\lambda_s^{\rm mom}=4\sigma/\sqrt{k_s}$.

\begin{theorem}[Full scans with only finite-variance replicates]\label{thm:mom_full}
Use the feasible constraints
$|d_U-\ip UB|\le\lambda_s^{\rm mom}$ in place of the Gaussian residual constraints. Their projection onto the equal-diagonal target is honest at level $1-\beta$, with the norm-sensitive diameter using $\lambda_s^{\rm mom}$. For fixed $K$, caps and confidence, both minimax expected-length criteria, with a supremum over all the stated replicate laws, have order $\mathfrak r_K(p,M,\sigma/\sqrt n)$ whenever \eqref{eq:mom_blocks} holds.

For squared error use $\beta=1/8$ and approximately minimize
\[
 F(B)=\max\bigl(0,\max_U\{|d_U-\ip UB|-\lambda_{s(U)}^{\rm mom}\}\bigr)
\]
over a fixed countable dense cap-class sequence, within tolerance $\tau$. The resulting scalar estimator has minimax squared risk of order $\mathfrak r_K^2$ over the same family. This is an existence construction. The sufficient replicate count is $n\ge C_{K,\beta}p$, and the full directional search remains exponential.
\end{theorem}

The moments $d_U$ need not be contractions of a single tensor estimate. Their separate linear constraints suffice. No Gaussian likelihood is used for the robust upper bounds; the lower bounds use the Gaussian subfamily of replicate distributions. This is a finite-direction use of the classical median-of-means principle \citep{devroye2016}.

\begin{proposition}[Two limits of a noise relaxation]\label{prop:noise_limits}
(i) In one dimension, known noise taking values $\pm v$ with equal probabilities permits exact recovery of every mean $|A|\le M<v$, although it has variance $v^2$ and satisfies \eqref{eq:joint_sg}. Thus a Gaussian lower bound cannot be transferred to every noise law from a variance bound alone.
(ii) For $E=v\xi\mathbf1$ with one Rademacher sign $\xi$, every coordinate has proxy $v$, but a flat unit product contraction on a $p^K$ block has variance $v^2p^K$. Coordinatewise calibration cannot certify a full scan or a dense rotated loading.
\end{proposition}

\subsection{Approximate ranks and growing caps}

\begin{theorem}[Sharp length under a Frobenius approximation allowance]\label{thm:approx_sharp}
In the original Gaussian experiment let
$\Theta^{\varepsilon_A}=\Theta(M,\bar r)+\varepsilon_A\mathbb B_F$, $\varepsilon_A\ge0$. For the equal diagonal, fixed $K\ge3$, fixed positive caps and fixed $0<\alpha<1/4$,
\[
 \mathcal L_0(\Theta^{\varepsilon_A})\asymp\mathcal L_{\max}(\Theta^{\varepsilon_A})
 \asymp \min\{\tau,\mathfrak r_K(p,M,\tau)+\varepsilon_A\}.
\]
The upper interval inflates every scan threshold by $\varepsilon_A$ and both projected endpoints by $\varepsilon_A$, with a deterministically selected direct Gaussian fallback. Thus approximate low rank is allowed, but the original rate requires $\varepsilon_A=O(\mathfrak r_K)$.
\end{theorem}

\begin{proposition}[A growing-rank lower bound and the unrestricted endpoint]\label{prop:growing_caps}
Let $r_*$ be the minimum supplied cap, truncated at $p$. Under equal loading every honest interval obeys
\[
 \E_0|C|\ge\frac{c_\alpha}{\sqrt p}
 \max_{1\le s\le r_*}
 \min\{M\sqrt s,\tau s\sqrt{\log(1+p/(4s^2))}\}.
\]
The finite upper envelopes of the preceding sections remain valid for dimension-dependent supplied caps with their displayed cap constants; these constants are not uniformly bounded. If all rank restrictions are removed, then for any prescribed unit loading both length criteria have order $\min(M,\tau)$, attained in polynomial time. The vanishing tensor rate therefore does not persist under unrestricted ranks.
\end{proposition}

\subsection{Arbitrary loading tensors and efficient regimes}

For a nonzero, otherwise arbitrary deterministic tensor $L$, a useful sufficient description is a certified decomposition $L=L_0+R$, where $L_0$ has a tractable orthogonal loading representation and the remainder $R$ is known exactly. Normalization is not required in this subsection.

\begin{proposition}[A confidence interval for any certified loading decomposition]\label{prop:loading_transfer}
Suppose $C_0=[a_0,b_0]$ is honest for $\ip{L_0}A$ at error $\alpha_0$, and
$\PP(|\ip RE|>q_R)\le\alpha_1$ uniformly. Then
\[
 C_L=[a_0+\ip RZ-q_R,\ b_0+\ip RZ+q_R]
\]
is honest at error $\alpha_0+\alpha_1$ and has length $|C_0|+2q_R$.
Under \eqref{eq:joint_sg}, one may take
\[
 q_R=v\norm R_F\sqrt{2\log(2/\alpha_1)}.
\]
For Gaussian noise use the normal quantile. No independence between the two terms is required. A decomposition or partition selected from loading and supplied constants alone needs no selection penalty. Selection using $Z$ does require simultaneous calibration.

This also supplies additional unequal-weight upper bounds: partition the orthogonal terms into prescribed blocks, apply their separate weighted bounds with error budgets summing to $\alpha$, and add the intervals. Include the unsplit construction among the choices. The minimum among their deterministic diameter bounds is achievable by choosing the construction before observing data.
\end{proposition}

\begin{theorem}[Two efficiently attained minimax regimes]\label{thm:efficient_regimes}
(i) For the equal diagonal, fixed order/caps and $M\le\tau\sqrt{\log(ep)}$, the structural interval is minimax up to fixed constants, with length of order $M/\sqrt p$.
(ii) Let a parameter class contained in $\{A:\norm A_F\le M\}$ contain every rank-one tensor of Frobenius norm at most $M$. If a supplied unit product tensor $U$ certifies
$|\ip LU|\ge c\norm L_F$ for a fixed $c>0$, then under Gaussian observations both length criteria have order $\norm L_F\min(M,\tau)$. Choose between the norm-only structural interval and the direct Gaussian interval. The algorithm uses a contraction and scalar comparisons, without rank recovery or tensor optimization. This covers arbitrary loadings with a certified constant product overlap, regardless of supplied rank caps.
\end{theorem}

The general loading transfer is a valid upper bound, not a characterization of the optimum for every tensor $L$. The near-flat difficult contrasts do not satisfy a dimension-independent product-overlap condition.

\begin{proposition}[Why full convexification loses the weak-signal gain]\label{prop:convex_loss}
Any convex parameter set containing the rank-one Frobenius ball of radius $M$ also contains $\{tL:|t|\le M/\sqrt p\}$ for the equal diagonal loading. Consequently, every interval honest on that convex set satisfies
\[
 \E_0|C|\ge c_\alpha\min(M/\sqrt p,\tau).
\]
At $M=\tau\sqrt p$ this is order $\tau$, exceeding the original full-data rate by an unbounded factor as $p\to\infty$. The conclusion concerns procedures guaranteeing coverage on the whole convex outer class; it is not a computational lower bound for every algorithm using a convex subroutine.
\end{proposition}

\subsection{Actual rank: an obstruction and a computable separated case}

\begin{theorem}[An adaptation obstruction at nonzero tensors]\label{thm:shifted_adaptation}
Suppose $A_0$ is supported on $\{1,\ldots,q\}^K$, $q<p$, has norm $a<M$, and every supplied unfolding cap is at least its actual rank plus one. Put $m=p-q$ and $M_0=(M^2-a^2)^{1/2}$. For the equal diagonal every interval honest on the outer cap class satisfies
\[
 \E_{A_0}|C|\ge c_{K,\alpha}\sqrt{m/p}\,
           \mathfrak r_K(m,M_0,\tau).
\]
Thus arbitrarily weak extra components obstruct shortening at some nonzero lower-rank tensors, not only at zero. This is a lower bound, not a complete characterization of adaptation on all rank subclasses.
\end{theorem}

An explicit positive result uses two independent observations
$Z^{(j)}=A+E^{(j)}$, with directional proxies $v_j$ as in \eqref{eq:joint_sg}. Let $d_k=\prod_{\ell\ne k}p_\ell$, choose $0<\delta<\alpha$, and set
\begin{equation}\label{eq:pilot_eta}
 \eta_k=2v_1\sqrt{2\{(p_k+d_k)\log9+\log(2K/\delta)\}}.
\end{equation}
The allowed class has $\norm A_F\le M$ and, in each mode, either rank zero or smallest nonzero singular value strictly greater than $2\eta_k$. No ranks or rank caps are supplied to the algorithm. Separation is a model assumption; observed singular values alone cannot certify the absence of undetected weak components.
Compute $\widehat r_k=\#\{j:\sigma_j(Z^{(1)}_{(k)})>\eta_k\}$ and let $Q_k$ project onto those left singular vectors. If any $\widehat r_k=0$, return $\{0\}$. Otherwise put
\[
 d_k^*=\min\left\{1,\frac{2\eta_k}{\sigma_{\widehat r_k}(Z^{(1)}_{(k)})-\eta_k}\right\}.
\]
Here $d_k^*$ is an error bound, distinct from the unfolding column dimension $d_k$.
For the commuting mode projectors define the orthogonal projector
\[
 T=\bigotimes_k Q_k+
       \sum_k (I-Q_k)\otimes\bigotimes_{\ell\ne k}Q_\ell.
\]
The tensor factors in this expression occupy their respective modes. This first-order product-space projector may be larger than the exact Tucker manifold tangent projector at ranks above one. Set
\begin{align}
 B_L&=M\norm{(I-T)L}_F
      \min\left\{1,\left(\sum_{k<\ell}(d_k^*d_\ell^*)^2\right)^{1/2}\right\},\label{eq:adaptive_bias}\\
 h_T&=B_L+v_2\norm{TL}_F\sqrt{2\log(2/(\alpha-\delta))}.
\end{align}

\begin{theorem}[Finite-sample rank calibration with a computable bias correction]\label{thm:adaptive_tangent}
The interval centered at $\ip{TL}{Z^{(2)}}$ with radius $h_T$ is honest at level $1-\alpha$ on the separated class just specified, for any prescribed loading $L$. One may choose, using only the pilot and known quantities, the shortest-radius option among it, the direct second-observation interval with error $\alpha-\delta$, and the norm-only structural interval. It may also be intersected with the structural interval. For Gaussian second-observation noise the square-root logarithm can be replaced by $z_{1-(\alpha-\delta)/2}$.

With probability at least $1-\delta$, all selected mode ranks equal the actual ranks and \eqref{eq:adaptive_bias} bounds the functional bias. The method consists of unfolding SVDs and tensor projections and is polynomial in the explicit input size. Its noise radius depends on $\norm{TL}_F$, and its additional cost is the displayed bias term. General minimax optimality and adaptation without the spectral separation are not asserted.
\end{theorem}

For one Gaussian observation, independent artificial $G'$ gives
$Z^{(1)}=Z+\tau G'$ and $Z^{(2)}=Z-\tau G'$, with independent errors of scale $\sqrt2\tau$. For non-Gaussian errors, this trick need not give independence; genuinely independent replicate groups suffice. The rank statement concerns Tucker mode ranks, and applies to TT tensors meeting the mode-gap conditions, but does not identify their TT bond ranks. Subspace projection is established methodology \citep{xu2025}; the present contribution is an explicit conservative finite-sample certificate in this setting.

\paragraph{Assumption audit.}
Gaussian isotropy gives the exact diagonal sequence and the likelihood-ratio identities. The threshold interval alone extends whenever a valid simultaneous coordinate noise envelope is supplied, but the direct alternative and lower bounds require their own calibration. Orthonormal loading families produce Bessel bounds and equal noise variances; general loading results require the additional conditions stated later in this supplement. Canonical tensor expansions avoid a spectral gap but yield potentially large rank constants. Supplied caps and $M$ identify the outer honesty class; replacing them with unvalidated estimates is not covered. None of these assumptions should be described as a routine approximation for dependent LLM representations.

\appendix
\section{Complete proofs}\label{app:proofs}

The proofs apply in finite dimensions. Constants used in the main text are explicit. All tensors and Gaussian variables are real.

\subsection{Canonical product expansions}

\begin{lemma}[Controlled squared component weights]\label{lem:cp}
For a tensor in either cap class there is a finite representation
\[
 A=\sum_{j=1}^{J_A}x_{1j}\otimes\cdots\otimes x_{Kj},\qquad
 a_j^{\mathrm{cp}}=\prod_{k=1}^K\norm{x_{kj}}_2,
\]
such that $J_A\le J$ and $\sum_j(a_j^{\mathrm{cp}})^2=\norm A_F^2$.
\end{lemma}


\noindent\emph{Roadmap.} Choose orthonormal Tucker mode bases or left-orthogonal TT cores; expand product terms and sum their squared weights. The TT argument uses a prefix identity, not orthogonality of individual paths.
\begin{proof}
The zero tensor has a zero-term representation. For a nonzero Tucker tensor, choose orthonormal bases $U_k$ for the column spaces of its mode unfoldings. Orthogonal projection onto each such space leaves $A$ unchanged. Contracting these bases gives a core $\mathcal G$ with $A=\mathcal G\times_1U_1\cdots\times_KU_K$ and $\norm{\mathcal G}_F=\norm A_F$. The actual mode ranks satisfy $r_k\le\bar r_k$.

Fix a mode $k$. Expand all core indices other than $k$. Each tuple $j_{-k}$ gives unit factors $U_\ell[:,j_\ell]$ in modes $\ell\ne k$ and factor $U_k\mathcal G[:,j_{-k}]$ in mode $k$. Its squared product weight is $\norm{\mathcal G[:,j_{-k}]}_2^2$. These slice squares sum to $\norm{\mathcal G}_F^2$. The number of terms is $\prod_{\ell\ne k}r_\ell$. Choosing $k$ to minimize the corresponding cap product proves the Tucker assertion.

For a TT tensor, sequential matrix factorizations yield cores $G_k\in\R^{r_{k-1}\times p_k\times r_k}$, with $r_0=r_K=1$ and $r_k\le\bar r_k$. At each cut, factor the relevant matrix through a basis of its column space. Equivalently, left-orthogonalize the cores by QR and absorb the triangular factors into the next core. Thus the $(r_{k-1}p_k)\times r_k$ matrix associated with $G_k$ has orthonormal columns for $k<K$, and $\norm{G_K}_F=\norm A_F$. These steps use no inverse lower bound on a singular value.

Expand the contraction over all internal paths $(a_1,\ldots,a_{K-1})$. A path contributes the product tensor with physical vectors $G_k[a_{k-1},:,a_k]$. There are at most $\prod_{k<K}r_k\le J$ paths. To sum their squared weights, let
\[
 t_k(b)=\sum_{a_1,\ldots,a_{k-1}}
 \prod_{\ell=1}^k\norm{G_\ell[a_{\ell-1},:,a_\ell]}_2^2,
 \qquad a_k=b.
\]
Column normalization gives $t_1(b)=1$. If $t_{k-1}(a)=1$ for all $a$, then
\[
 t_k(b)=\sum_a t_{k-1}(a)\norm{G_k[a,:,b]}_2^2=1
\]
for $k<K$. Hence the sum of all squared path weights is
\[
 \sum_a t_{K-1}(a)\norm{G_K[a,:,1]}_2^2=\norm{G_K}_F^2=\norm A_F^2.
\]
The individual path terms need not be orthogonal; the prefix identity is what proves the required equality.
\end{proof}

\subsection{Proof of Theorem~\ref{thm:diagonal}}


\noindent\emph{Roadmap.} Apply subadditivity, modewise H\"older and Bessel, then control the component-weight sum.
\begin{proof}
Let $q=2/K\le1$ and use Lemma~\ref{lem:cp}. Subadditivity of the $q$th power of absolute value gives
\[
 \sum_i|\ip A{V_i}|^q
 \le\sum_j\sum_i\prod_{k=1}^K|\ip{x_{kj}}{v_{ki}}|^{2/K}.
\]
H\"older's inequality with $K$ equal exponents, followed by Bessel's inequality for each orthonormal family, yields
\begin{align*}
 \sum_i\prod_k|\ip{x_{kj}}{v_{ki}}|^{2/K}
 &\le\prod_k\left(\sum_i|\ip{x_{kj}}{v_{ki}}|^2\right)^{1/K}\\
 &\le\prod_k\norm{x_{kj}}_2^{2/K}=(a_j^{\mathrm{cp}})^q.
\end{align*}
Finally, the power-mean inequality (or H\"older with exponents $2/q$ and $2/(2-q)$) gives
\[
 \sum_j(a_j^{\mathrm{cp}})^q
 \le J_A^{1-q/2}\left\{\sum_j(a_j^{\mathrm{cp}})^2\right\}^{q/2}
 \le J^{1-q/2}M^q.
\]
This is \eqref{eq:compress}. For $J_A=0$ both sides of the first inequality are zero.
\end{proof}

\subsection{Exact rank-one diagonal image}

\begin{lemma}[Exact diagonal reduction]\label{lem:image}
Under \eqref{eq:loading}, the image of the rank-one tensor Frobenius ball of radius $M$ under $A\mapsto(\ip A{V_i})_{i=1}^m$ is exactly
\[
 \left\{a\in\R^m:\sum_i|a_i|^{2/K}\le M^{2/K}\right\}.
\]
\end{lemma}


\noindent\emph{Roadmap.} Use the diagonal bound for one inclusion and construct signed product factors for the converse.
\begin{proof}
One inclusion follows from Theorem~\ref{thm:diagonal} with $J=1$. Conversely, for any vector in the displayed ball, put
\[
 x_1=\sum_i\operatorname{sign}(a_i)|a_i|^{1/K}v_{1i},\qquad
 x_k=\sum_i|a_i|^{1/K}v_{ki},\quad k\ge2,
\]
and set $A=x_1\otimes\cdots\otimes x_K$. Orthonormality gives $\ip A{V_i}=a_i$ and
\[
 \norm A_F=\prod_k\norm{x_k}_2
 =\left(\sum_i|a_i|^{2/K}\right)^{K/2}\le M.
\]
The tensor has every Tucker and TT rank at most one. The zero vector is handled by the zero tensor.
\end{proof}

\subsection{Threshold error and structural radius}

\begin{lemma}[Scalar threshold envelope]\label{lem:soft}
If $|e|\le\lambda$ and $s=\soft_\lambda(x+e)$, then $s$ lies between zero and $x$ and
\[
 |s-x|\le\min(|x|,2\lambda).
\]
For $q\in(0,1]$, $t,u\ge0$, one also has $\min(t,u)\le t^q u^{1-q}$, with the continuous interpretation at zero.
\end{lemma}


\noindent\emph{Roadmap.} Check positive and negative scalar inputs, then interpolate the two deterministic error bounds.
\begin{proof}
For $x\ge0$, $x+e\ge-\lambda$, so thresholding cannot produce a negative value. Also $x+e-\lambda\le x$, so $0\le s\le x$. If $x>2\lambda$, then $x+e>\lambda$ and $x-s=\lambda-e\le2\lambda$. Otherwise $x-s\le x\le2\lambda$. The negative case follows by oddness. If $t\le u$, then $t=t^q t^{1-q}\le t^q u^{1-q}$; if $u\le t$, then $u=u^q u^{1-q}\le t^q u^{1-q}$.
\end{proof}

\begin{lemma}[Matrix functional bound]\label{lem:structural}
If $\rank(B)\le r$ and $H$ has singular values $s_1(H)\ge s_2(H)\ge\cdots$, then
\[
 |\ip HB|\le\norm B_F\left(\sum_{i\le r}s_i(H)^2\right)^{1/2}.
\]
\end{lemma}


\noindent\emph{Roadmap.} Project onto the matrix column space and maximize the spectral trace over projections of bounded rank.
\begin{proof}
Let $P$ project onto the column space of $B$. Then $PB=B$ and Cauchy--Schwarz gives $|\ip HB|=|\ip{PH}B|\le\norm{PH}_F\norm B_F$. Write $HH^\top=\sum_i\lambda_i u_i u_i^\top$, with decreasing nonnegative eigenvalues. The numbers $b_i=u_i^\top P u_i$ satisfy $0\le b_i\le1$ and $\sum_i b_i=\operatorname{tr}(P)\le r$. Therefore $\norm{PH}_F^2=\sum_i\lambda_i b_i\le\sum_{i\le r}\lambda_i$: moving any weight from a smaller eigenvalue to an unfilled larger one cannot reduce this sum. Since $\lambda_i=s_i(H)^2$, the result follows.
\end{proof}

\subsection{Proof of Theorem~\ref{thm:coverage}}


\noindent\emph{Roadmap.} Construct a uniform coordinate-noise event, bound the threshold error pathwise, verify the other intervals and use the deterministic selection rule.
\begin{proof}
The product atoms are orthonormal because
\[
 \ip{V_i}{V_j}=\prod_k\ip{v_{ki}}{v_{kj}}=\mathbf1\{i=j\}.
\]
Gaussian isotropy then gives the independent normals in \eqref{eq:diagonal}. For a standard normal $g$, Chernoff's inequality gives $\PP(g>t)\le\exp(-t^2/2)$ by minimizing $\exp(-ut+u^2/2)$ at $u=t$. Apply the same inequality to $-g$ and take a union bound to obtain
\[
 \PP\{\max_i|\tau g_i|>\lambda\}
 \le2m\exp\{-\lambda^2/(2\tau^2)\}=\alpha.
\]
On the complementary event, Lemma~\ref{lem:soft} and Theorem~\ref{thm:diagonal} give
\begin{align*}
 |\widehat\psi_{\NL}-\psi(A)|
 &\le\norm w_\infty\sum_i\min(|a_i|,2\lambda)\\
 &\le\norm w_\infty(2\lambda)^{1-q}\sum_i|a_i|^q\le h_{\NL}.
\end{align*}
This proves the uniform coverage of $I_{\NL}$.

Because $\norm L_F=1$, $\ip LZ-\psi(A)\sim N(0,\tau^2)$, so $I_G$ has exact coverage. For any constrained unfolding, the loading matrix has singular values $|w_i|$: the product vectors on each side form orthonormal families. Lemma~\ref{lem:structural} thus gives $|\psi(A)|\le M(\sum_{i\le r}w_{(i)}^2)^{1/2}$. Choosing the smallest cap proves $|\psi(A)|\le h_0$, so $I_0$ has coverage one. The selected index in \eqref{eq:select} is deterministic, hence the selected interval retains its individual coverage without a multiple-selection correction. Its length is exactly \eqref{eq:length}.
\end{proof}

\subsection{Proof of Theorem~\ref{thm:affine}}


\noindent\emph{Roadmap.} Force a variance cost at zero and a bias cost at signed coordinate tensors, then combine the two necessary inequalities.
\begin{proof}
Set $t=\ip WL$. At $A=0$ the center error is normal with mean $b$ and standard deviation $v=\tau\norm W_F$. For fixed $h$ and $v>0$, the probability that this variable belongs to $[-h,h]$ is maximized at $b=0$. Indeed, for $b\ge0$ its derivative is
\[
 v^{-1}\{\phi((h+b)/v)-\phi((h-b)/v)\}\le0,
\]
since $|h+b|\ge|h-b|$; the function is even in $b$. Honesty at zero therefore implies $h\ge z\tau\norm W_F\ge z\tau|t|$. If $v=0$, the same last inequality is automatic.

At $A=\pm M E_i$, the mean center error is
\[
 b\pm M(\ip W{E_i}-p^{-1/2}).
\]
At least one sign has absolute mean at least $M|\ip W{E_i}-p^{-1/2}|$. A symmetric interval whose radius is strictly smaller than this mean magnitude has probability below $1/2$ under a nondegenerate normal distribution: it lies strictly to one side of the median. In the degenerate case the coverage is zero. Since $1-\alpha>1/2$, honesty forces
\[
 h\ge M\max_i|\ip W{E_i}-p^{-1/2}|
 \ge\frac M{\sqrt p}|1-t|.
\]
The last inequality follows from $t-1=p^{-1/2}\sum_i(\ip W{E_i}-p^{-1/2})$. Put $a=M/\sqrt p$ and $b'=z\tau$. The two bounds imply
\[
 1\le|1-t|+|t|\le h/a+h/b',
\]
which rearranges to \eqref{eq:affinelower}. On the rank-one tensor ball, $h_0=M/\sqrt p$, and the direct/structural choice gives the matching order. No result about arbitrary data-dependent widths is invoked.
\end{proof}

\subsection{A mixture inequality for confidence intervals}

\begin{lemma}[Common-target alternatives]\label{lem:mixture}
Let $P_0$ have target zero and let $Q$ be a mixture of alternatives whose targets all equal $\delta>0$. If $C$ covers zero under $P_0$ and covers the target at every alternative with probability at least $1-\alpha$, then
\begin{equation}\label{eq:mixtureCI}
 \E_0|C|\ge\delta\{1-2\alpha-\TV(P_0,Q)\}.
\end{equation}
If $Q\ll P_0$, then $\TV(P_0,Q)\le\frac12\sqrt{\chi^2(Q,P_0)}$.
\end{lemma}


\noindent\emph{Roadmap.} Transfer coverage from the common-target mixture to the null using total variation, then intersect the two coverage events.
\begin{proof}
Mixture coverage gives $Q(\delta\in C)\ge1-\alpha$. By the definition of total variation, $P_0(\delta\in C)\ge1-\alpha-\TV(P_0,Q)$. Intersecting with $\{0\in C\}$ gives probability at least $1-2\alpha-\TV$. An interval containing both points has length at least $\delta$, proving \eqref{eq:mixtureCI}. The divergence bound follows from Cauchy--Schwarz applied to $\frac12\E_0|dQ/dP_0-1|$. Appending an independent random seed leaves total variation unchanged, so randomized intervals are covered as well.
\end{proof}

For Gaussian shifts $P_B=N(\operatorname{vec}(B),\tau^2I_D)$, direct integration of their likelihood ratios gives
\begin{equation}\label{eq:LR}
 \E_0\left[\frac{dP_B}{dP_0}\frac{dP_C}{dP_0}\right]
 =\exp\{\ip BC/\tau^2\}.
\end{equation}
To verify this, each ratio is $\exp\{\ip ZB/\tau^2-\norm B_F^2/(2\tau^2)\}$. Multiply the ratios and use the moment generating function of $\ip Z{B+C}$, whose variance under $P_0$ is $\tau^2\norm{B+C}_F^2$. The exponent simplifies to \eqref{eq:LR}.

\subsection{Proof of Theorem~\ref{thm:unfolding}}


\noindent\emph{Roadmap.} Construct rank-one matrix alternatives within the chosen unfolding and bound their full Gaussian likelihood second moment using an auxiliary Gaussian.
\begin{proof}
For the given bipartition, let $f_i=\bigotimes_{k\in S}e_i$ and $g_i=\bigotimes_{k\notin S}e_i$. Both families are orthonormal. Let $u=p^{-1/2}(s_1,\ldots,s_p)$ with independent uniform signs, and define
\[
 \theta=\min(M,\tau\sqrt p/4),\qquad
 A_u=\theta\left(\sum_i u_i f_i\right)\otimes
                    \left(\sum_j u_j g_j\right).
\]
Every alternative has rank one in this unfolding, norm $\theta$, and functional value $\delta=\theta/\sqrt p$. Let $Q$ average their Gaussian laws. If $v$ is an independent copy of $u$, then \eqref{eq:LR} gives
\[
 1+\chi^2(Q,P_0)=\E_{u,v}\exp\{\theta^2(u^\top v)^2/\tau^2\}
 =\E\exp\{a S_p^2/p\},
\]
where $S_p$ is a sum of $p$ independent signs and $a=\theta^2/(\tau^2p)\le1/16$.

Introduce an independent standard normal $H$. Its moment generating function yields
\begin{align*}
 \E\exp(aS_p^2/p)
 &=\E_H\left\{\cosh(\sqrt{2a/p}\,H)\right\}^p\\
 &\le\E_H\exp(aH^2)=(1-2a)^{-1/2}\le\sqrt{8/7}.
\end{align*}
The inequality uses $\log\cosh(x)\le x^2/2$: for $x\ge0$, its derivative $\tanh(x)$ is at most $x$, and both functions are even. The Gaussian-square expectation follows by integrating $\exp\{-(1-2a)h^2/2\}/\sqrt{2\pi}$; it is finite since $a<1/2$. Nonnegative integrands justify exchanging the expectations.

Lemma~\ref{lem:mixture} now gives \eqref{eq:unfoldlower}. The constant is positive for $\alpha<1/4$. On $\Theta_S$, Lemma~\ref{lem:structural} with loading singular values all equal to $p^{-1/2}$ gives $|\psi(A)|\le M/\sqrt p$. Choosing between this structural interval and the direct interval attains the lower order. All calculations use the entire Gaussian experiment; the result does not arise from hiding off-diagonal observations.
\end{proof}

\subsection{Proof of Theorem~\ref{thm:separation}}


\noindent\emph{Roadmap.} Substitute the specified dimensions and norm scaling, then apply the two lower-bound theorems with their confidence-level restrictions.
\begin{proof}
Substitute $M=\tau\sqrt p$, $q=2/K$ and $\norm w_\infty=p^{-1/2}$ into \eqref{eq:nl}. This gives
\[
 h_{\NL}=\tau J^{1-1/K}
 p^{1/K-1/2}\{8\log(2p/\alpha)\}^{(K-2)/(2K)},
\]
which is the radius in \eqref{eq:separation}. For fixed $K\ge3$, its exponent $\beta=(K-2)/(2K)$ is positive and $(\log p/p)^\beta\to0$. The cap classes contain all $\pm M E_i$ and zero, so Theorem~\ref{thm:affine} gives radius at least $z\tau/(1+z)$. For any retained constrained unfolding, the corresponding matrix-rank cap class contains $\Theta_S$, and Theorem~\ref{thm:unfolding} gives expected length at zero at least $c_\alpha\tau/4$. The same comparisons under the stated growing-cap condition follow by retaining the factor $J^{1-1/K}$.
\end{proof}

\subsection{Proof of Theorem~\ref{thm:jointlower}}


\noindent\emph{Roadmap.} Mix orthogonal coordinate shifts, compute the full likelihood second moment exactly and apply the interval mixture lemma.
\begin{proof}
Let $a=\min\{M,\tau\sqrt{\log(1+p/4)}\}$ and mix uniformly over $A_j=aE_j$, $j=1,\ldots,p$. Each is in the rank-one tensor class and has target $\delta=a/\sqrt p$. Orthogonality of the $E_j$ and \eqref{eq:LR} imply
\[
 1+\chi^2(Q,P_0)=p^{-2}\{p\exp(a^2/\tau^2)+p(p-1)\}.
\]
Thus $\chi^2(Q,P_0)=\{\exp(a^2/\tau^2)-1\}/p\le1/4$ and $\TV(P_0,Q)\le1/4$. Lemma~\ref{lem:mixture} gives \eqref{eq:jointlower}. This construction does not establish that these alternatives are least favorable.
\end{proof}


\subsection{Proof of Proposition~\ref{prop:diagonal_minimax}}


\noindent\emph{Roadmap.} Construct sparse diagonals inside the rank-one tensor class, bound the diagonal likelihood mixture using factorial moments and balance the support size. The full-data likelihood is kept separate.
\begin{proof}
For an integer $1\le s\le p$, let $S$ be a uniformly chosen size-$s$ subset of $\{1,\ldots,p\}$. Define
\[
 A_S=b_s\left(\sum_{i\in S}e_i\right)^{\otimes K}.
\]
This is a rank-one tensor, $\norm{A_S}_F=b_s s^{K/2}\le M$, and every alternative has target $\delta_s=s b_s/\sqrt p$. Its observed diagonal has mean $b_s\mathbf1_S$. Let $Q_d$ be the uniform mixture of these diagonal Gaussian laws, and let $P_{0,d}$ be the zero-mean diagonal law. Applying the Gaussian shift calculation \eqref{eq:LR} in dimension $p$ gives
\begin{equation}\label{eq:diag_chisquare}
 1+\chi^2(Q_d,P_{0,d})
 =\E_{S,T}\exp\{b_s^2|S\cap T|/\tau^2\},
\end{equation}
where $S,T$ are independent uniform size-$s$ subsets. This is the sparse-mixture mechanism used in classical sequence lower bounds \citep{collier2017}; the following counting argument records the needed bound explicitly.

Conditional on $S$, $H=|S\cap T|$ is hypergeometric with population size $p$, success count $s$ and sample size $s$. Write $(x)_j=x(x-1)\cdots(x-j+1)$, with $(x)_0=1$. Counting the $j$-element subsets of $S$ contained in $T$ yields
\[
 \E\binom Hj=\binom sj\frac{(s)_j}{(p)_j}
 \le\binom sj(s/p)^j,\qquad 0\le j\le s.
\]
For $j\ge1$, the inequality follows factor by factor from
$(s-\ell)/(p-\ell)\le s/p$ for $0\le\ell<j$, since $s\le p$.
Set $t=\exp(b_s^2/\tau^2)-1\ge0$. Expanding $(1+t)^H$ and using the preceding factorial moments gives
\begin{align*}
 \E\exp(b_s^2H/\tau^2)
 &=\sum_{j=0}^s t^j\E\binom Hj\\
 &\le (1+st/p)^s
 \le \exp(s^2t/p)
 \le e^{1/4}.
\end{align*}
The last inequality uses the definition of $b_s$. Equation~\eqref{eq:diag_chisquare} and Cauchy--Schwarz therefore give
\[
 \TV(Q_d,P_{0,d})\le\frac12\sqrt{e^{1/4}-1}.
\]
Lemma~\ref{lem:mixture}, now applied to the diagonal experiment, proves \eqref{eq:diagonal_lower}. The constant is positive because $\alpha<1/4$ and $\frac12\sqrt{e^{1/4}-1}<1/2$. The argument also covers independent randomization.

To prove \eqref{eq:diagonal_rate}, let
$x_p=(p/\log p)^{1/K}$ and $s_p=\lfloor x_p\rfloor$. For sufficiently large $p$,
$x_p/2\le s_p\le x_p$ and $s_p\ge1$. When $M=\tau\sqrt p$, the norm term in $b_{s_p}$ is at least $\tau\sqrt{\log p}$. The noise term is at most this quantity, since $1+p/(4s_p^2)\le1+p/4\le p$ for $p\ge2$. Thus the minimum equals the noise term. Since $K>2$ and $s_p\le p^{1/K}$ eventually,
\[
 \log\{1+p/(4s_p^2)\}
 \ge\log\{1+p^{1-2/K}/4\}
 \ge\tfrac12(1-2/K)\log p
\]
for sufficiently large $p$. Consequently,
\[
 \frac{s_p b_{s_p}}{\sqrt p}
 \ge \frac12\sqrt{\frac{1-2/K}{2}}\,
 \tau\left(\frac{\log p}{p}\right)^{(K-2)/(2K)}.
\]
Theorem~\ref{thm:coverage} with $J=1$ gives an upper bound of the same order and deterministic length. Since the worst-case criterion is at least the criterion at zero, both rates follow. The fixed-cap extension follows by class inclusion for the lower bound and a constant $J^{1-1/K}$ for the upper bound.

For completeness, the corresponding \emph{full-tensor} mixture would have
$\ip{A_S}{A_T}=b_s^2|S\cap T|^K$ in \eqref{eq:LR}. Its likelihood second moment is not \eqref{eq:diag_chisquare}, and the preceding bound does not apply to it. The proposition is restricted to diagonal observations throughout.
\end{proof}


\section{Proofs of the full-data refinements}\label{app:full_data_proofs}

\begin{lemma}[Finite product nets and injective norm]\label{lem:net_scan}
The nets used in \eqref{eq:scan_threshold} exist and satisfy
$\rho(B)\le2N_S(B)$ for any tensor on $S^K$.
Furthermore, a tensor $B$ in either rank-cap class satisfies
$\norm B_F\le\sqrt J\,\rho(B)$.
\end{lemma}
\noindent\emph{Roadmap.} Use a volume bound and product telescoping, then apply the canonical squared-weight identity.
\begin{proof}
A maximal $\varepsilon$-separated subset of the unit sphere is an $\varepsilon$-net. Balls of radius $\varepsilon/2$ around its points are disjoint and contained in the ball of radius $1+\varepsilon/2$. The ratio of their Euclidean volumes bounds the cardinality by $(1+2/\varepsilon)^s$. Compactness gives a finite maximal set. Take $\varepsilon=1/(2K)$.

Approximate each factor of any unit product direction by a net point. Telescoping the product one factor at a time, each contraction difference has magnitude at most $\varepsilon\rho(B)$, by multilinearity and the definition of $\rho$. The sum is at most $K\varepsilon\rho(B)=\rho(B)/2$. Taking the supremum gives $\rho(B)\le N_S(B)+\rho(B)/2$.

For the second assertion use Lemma~\ref{lem:cp}, writing its product terms as $X_j$. Then
\[
 \norm B_F^2
 =\sum_j\ip B{X_j}
 \le\rho(B)\sum_j\norm{X_j}_F
 \le\rho(B)\sqrt J
       \left(\sum_j\norm{X_j}_F^2\right)^{1/2}
 =\rho(B)\sqrt J\,\norm B_F.
\]
The middle inequality is Cauchy--Schwarz, and the last identity is the squared-weight conclusion of Lemma~\ref{lem:cp}. Divide when $B\ne0$; the zero case is immediate.
\end{proof}

\subsection{Proof of Theorem~\ref{thm:scan}}
\noindent\emph{Roadmap.} A finite Gaussian union bound controls every net contraction. Test acceptance on that event bounds each ordered diagonal coordinate. Summation gives the interval and intersection gives the two expected-length guarantees.
\begin{proof}
There are at most $(ep/s)^s$ subsets of size $s$, since
$\binom ps\le p^s/s!$ and
$\log(s!)\ge\int_1^s\log x\,dx\ge s\log s-s$.
Each unit product contraction of $\tau G$ has distribution $N(0,\tau^2)$. Therefore the Gaussian Chernoff bound gives
\begin{align*}
 \PP\{\exists S:N_S(\tau G)>\lambda_{|S|}^{\rm scan}\}
 &\le\sum_{s=1}^p
 2(ep/s)^s(1+4K)^{Ks}
     \exp\{-(\lambda_s^{\rm scan})^2/(2\tau^2)\}\\
 &\le\gamma.
\end{align*}
Let $\mathcal E_\gamma$ be the complementary event. No independence between the scanned contractions is required.

On $\mathcal E_\gamma$, if the test accepts, the triangle inequality for the finite maximum implies $N_S(A)\le2\lambda_s^{\rm scan}$. Lemma~\ref{lem:net_scan} gives $\rho(A_{S,\ldots,S})\le4\lambda_s^{\rm scan}$. Coordinate restriction cannot increase any constrained unfolding rank, because it multiplies an unfolding by selection matrices. Theorem~\ref{thm:diagonal} and Lemma~\ref{lem:net_scan} thus imply
\[
 \sum_{i\in S}|A_{i,\ldots,i}|^q
 \le J^{1-q/2}\norm{A_{S,\ldots,S}}_F^q
 \le J\,\rho(A_{S,\ldots,S})^q.
\]
Let $|a|_{(s)}$ be the $s$th largest absolute diagonal entry. Choose $S$ to contain the $s$ largest entries. It follows that
\[
 s|a|_{(s)}^q\le J(4\lambda_s^{\rm scan})^q,
 \qquad
 |a|_{(s)}\le4J^{K/2}\lambda_s^{\rm scan}s^{-K/2}.
\]
Summing these inequalities and using
$|\psi(A)|\le p^{-1/2}\sum_i|a_i|$
proves the bound $|\psi(A)|\le R_\gamma$ in \eqref{eq:scan_radius}. On rejection the structural interval covers surely. Therefore $C_{\rm scan}$ covers on $\mathcal E_\gamma$ and is honest at level $1-\gamma$.

Both intersected intervals cover with probability at least $1-\alpha_0-\gamma$. On that event the intersection is nonempty and covers; replacing empty intersections by $\{0\}$ preserves this argument. Its length is always at most $\ell_*$. At zero, the rejection probability is at most $\gamma$, and on acceptance the intersection has length at most $2R_\gamma$. This gives \eqref{eq:scan_expectation}, without independence between the intervals.

For the specified $\gamma$, elementary logarithm bounds give
$\lambda_s^{\rm scan}\le C_K\tau\sqrt{s\log(2p/\alpha)}$.
For example, use $s\log(ep/s)\le s(1+\log p)$,
$\log(2p/\gamma)\le2\log(2p/\alpha)$ and $s\ge1$.
Equation~\eqref{eq:scan_radius} now proves \eqref{eq:scan_rate}; the extra term $\ell_*\gamma$ is smaller in order by Theorem~\ref{thm:separation}.
If $K\ge4$, the sum is at most $1+2/(K-3)$ by integration of $x^{-(K-1)/2}$ over $[1,\infty)$. If $K=3$, it is at most $1+\log p$. Combining the upper bound for $K\ge4$ with Theorem~\ref{thm:jointlower} proves \eqref{eq:zero_sharp}.
\end{proof}

\begin{lemma}[Uniform cubic moment for a weighted sign sum]\label{lem:cubic_moment}
For independent uniform signs $\epsilon_i$, let
$S_p=\sum_{i=1}^p i^{-2/3}\epsilon_i$. Then, uniformly in $p\ge1$,
\[
 \PP(|S_p|>x)\le2e^{-x^3/192}\quad(x\ge12),
 \qquad
 \E\exp\{|S_p|^3/40000\}\le20/19+1/100<5/4.
\]
\end{lemma}
\noindent\emph{Roadmap.} Split the weighted sign sum into a bounded head and a sub-Gaussian tail, then integrate its cubic tail probability.
\begin{proof}
For $x\ge12$, put $m=\lfloor(x/6)^3\rfloor$. If $m\ge p$, then
$|S_p|\le3p^{1/3}\le x/2$.
Otherwise, the first $m$ terms have absolute sum at most
$3m^{1/3}\le x/2$. The remaining independent sign sum has variance proxy
\[
 V_m=\sum_{i>m}i^{-4/3}\le3m^{-1/3}\le24/x.
\]
Here $m\ge(x/6)^3/2$ and $18\cdot2^{1/3}<24$.
The product inequality $\prod_i\cosh(ti^{-2/3})\le\exp(t^2V_m/2)$
and Chernoff optimization give a two-sided tail bound
$2\exp\{-(x/2)^2/(2V_m)\}\le2e^{-x^3/192}$.

Put $u=1/40000$ and $Y=|S_p|^3$. Tonelli's theorem applied to
$e^{uY}=1+\int_0^Y ue^{ut}\,dt$ yields
\[
 \E e^{uY}
 \le e^{1728u}
 +\frac{2u}{1/192-u}
       e^{-(1/192-u)1728}
 \le \frac{20}{19}+\frac1{100}.
\]
The last step uses $1728u<1/20$, $e^x\le(1-x)^{-1}$ for $0\le x<1$,
and $1/192-u\ge1/200$. The exponential factor in the second term is at most one. These estimates are uniform in $p$.
\end{proof}

\subsection{Proof of Theorem~\ref{thm:harmonic}}
\noindent\emph{Roadmap.} Couple three rank-one factor signs so that the diagonal is deterministic and harmonic. Control the full likelihood mixture by the cubic moment lemma and H\"older, without assuming independence of the three inner-product sums.
\begin{proof}
For independent uniform sign vectors $u,v\in\{-1,1\}^p$, define
\[
 (x_u)_i=u_i i^{-1/3},\qquad
 A_{u,v}=\frac{\tau}{200}x_u\otimes x_v\otimes x_{uv},
\]
where $(uv)_i=u_iv_i$. Every alternative is rank one and
\[
 \norm{A_{u,v}}_F
 =\frac{\tau}{200}\left(\sum_i i^{-2/3}\right)^{3/2}
 \le\frac{3^{3/2}}{200}\tau\sqrt p\le M.
\]
The sum bound follows from
$1+\int_1^p x^{-2/3}\,dx\le3p^{1/3}$.
Its diagonal is $\tau/(200i)$, so the common target is
$\delta=\tau H_p/(200\sqrt p)$.

Let $Q$ mix the full Gaussian laws over $u,v$. For independent copies
$u',v'$, set $\epsilon_i=u_iu_i'$ and $\eta_i=v_iv_i'$.
The full tensor inner product divided by $\tau^2$ is
$S_1S_2S_3/40000$, where
\[
 S_1=\sum_i i^{-2/3}\epsilon_i,\quad
 S_2=\sum_i i^{-2/3}\eta_i,\quad
 S_3=\sum_i i^{-2/3}\epsilon_i\eta_i.
\]
Each has the marginal distribution in Lemma~\ref{lem:cubic_moment}.
They are dependent. The likelihood identity \eqref{eq:LR}, the arithmetic--geometric mean inequality and H\"older give, with $u_0=1/40000$,
\begin{align*}
 1+\chi^2(Q,P_0)
 &=\E e^{u_0S_1S_2S_3}\\
 &\le\E e^{u_0(|S_1|^3+|S_2|^3+|S_3|^3)/3}\\
 &\le\prod_{j=1}^3\{\E e^{u_0|S_j|^3}\}^{1/3}<5/4.
\end{align*}
Thus $\TV(Q,P_0)<1/4$. Lemma~\ref{lem:mixture} gives
$\E_0|C|\ge(3/4-2\alpha)\delta$, proving \eqref{eq:harmonic_lower}.
Since $H_p\asymp\log p$, this also proves the claimed obstruction to the earlier zero-point conjecture.
\end{proof}


\section{Supporting refinements}\label{app:refinements}

These separately calibrated extensions do not change the observation model of a lower bound by implication. The simulations retain the original Gaussian procedure.

\subsection{Diagonal geometry}
\begin{proposition}[Smaller rank coefficient]\label{prop:budget}
In $\sum_i|a_i|^{2/K}\le C_{\mathrm{diag}} M^{2/K}$ one may take
\[
 C_{\mathrm{diag},\Tu}=\min\left\{J_{\Tu}^{1-1/K},m^{1-1/K},
                  (\prod_k\bar r_k)^{1/K}\right\}.
\]
For TT caps put $\bar r_0=\bar r_K=1$ and
$\rho_k=\min(p_k,\bar r_{k-1}\bar r_k)$; one may take
\[
 C_{\mathrm{diag},\TT}=\min\left\{J_{\TT}^{1-1/K},m^{1-1/K},
                  (\prod_k\rho_k)^{1/K}\right\}.
\]
The nonlinear radius remains valid with coefficient $C_{\mathrm{diag}}$.
\end{proposition}
\noindent\emph{Roadmap.} Add ambient and mode-subspace bounds to the original expansion bound.
\begin{proof}
Bessel gives $\sum_i|a_i|^2\le M^2$; H\"older in dimension $m$ gives the second coefficient. For Tucker mode projectors $P_k$ put $\ell_{ki}=\norm{P_kv_{ki}}_2^2$. Projection leaves $A$ unchanged, so Cauchy--Schwarz and H\"older give
\[
 \sum_i|a_i|^{2/K}
 \le M^{2/K}\sum_i\prod_k\ell_{ki}^{1/K}
 \le M^{2/K}\prod_k(\sum_i\ell_{ki})^{1/K}
 \le M^{2/K}\prod_k\bar r_k^{1/K}.
\]
The last step uses $\sum_i v_{ki}^\top P_kv_{ki}\le\operatorname{tr}P_k$.
For TT cores, fixing the two bond indices of core $k$ expresses its physical fibers through at most $\bar r_{k-1}\bar r_k$ vectors. The mode rank is therefore at most $\rho_k$; the Tucker argument applies. Theorem~\ref{thm:diagonal} gives the first coefficient.
\end{proof}
For equal Tucker caps $r$ this supplies coefficient $r$, instead of $r^{(K-1)^2/K}$; optimal rank dependence is not established.
For unit loading vectors with frame bounds $\sum_i v_{ki}v_{ki}^\top\preceq F_kI$, the original proof instead gives coefficient $J^{1-1/K}\prod_kF_k^{1/K}$. This follows by replacing each Bessel inequality by its frame bound. The direct variance must use the actual $\norm L_F^2$; the orthogonal structural formula is not asserted for these frames.

\subsection{Exact affine benchmark}
Let $Q_\alpha(\mu,v)$ be the smallest $h\ge0$ such that
$\PP\{|N(\mu,v^2)|\le h\}\ge1-\alpha$, for $\mu,v\ge0$, with
$Q_\alpha(\mu,0)=\mu$.

\begin{proposition}[Structural range and fixed affine optimum]\label{prop:affine_exact}
The range of $\psi$ on either cap class under \eqref{eq:loading} is exactly $[-h_0,h_0]$. Under \eqref{eq:equalL} and $0<\alpha<1/2$, set
$a_0=M\sqrt{s_{\rm cap}/p}$. Among all honest full-data fixed-radius affine intervals, the optimal radius is
\[
 h_{\rm aff}^*=\min_{0\le t\le1}Q_\alpha(a_0(1-t),\tau t).
\]
Center $t\ip LZ$ attains the optimum. If $a_0\le z_{1-\alpha}\tau$, then $h_{\rm aff}^*=a_0$.
\end{proposition}
\noindent\emph{Roadmap.} Attain the structural bound, force bias and variance for any affine center, and reduce to a scalar optimization.
\begin{proof}
Let $T$ index the $s_{\rm cap}$ largest weights. The tensor
$M\sum_{i\in T}w_iV_i/(\sum_{i\in T}w_i^2)^{1/2}$ has norm $M$, all constrained ranks at most $s_{\rm cap}$, and target $h_0$. Its signed scalar multiples attain the full interval; Lemma~\ref{lem:structural} gives the reverse inclusion.

For center $\ip WZ+b$ put $t=\ip WL$ and $g_i=\ip W{E_i}-p^{-1/2}$.
All $s_{\rm cap}$-sparse diagonal tensors of norm at most $M$ belong to the class. Opposite such tensors force an absolute mean error at least
\[
 |b|+M(\sum_{i=1}^{s_{\rm cap}}g_{(i)}^2)^{1/2}\ge a_0|1-t|.
\]
Indeed, the top-square sum is at least $s_{\rm cap}/p$ times the full sum, and
$\norm g_2\ge|\sum_i g_i|/\sqrt p=|t-1|$.
The noise standard deviation is $\tau\norm W_F\ge\tau|t|$.
Confidence greater than one half forces $h$ to be at least the absolute mean error, by the normal median argument.
For $h\ge\mu\ge0$ the probability
$F(h;\mu,v)=\Phi((h-\mu)/v)+\Phi((h+\mu)/v)-1$
decreases with $\mu$ and with $v>0$. For $\mu$, compare the normal densities. For $v$, its derivative is
\[
 -\{(h-\mu)\phi((h-\mu)/v)+(h+\mu)\phi((h+\mu)/v)\}/v^2\le0.
\]
Thus $h\ge Q_\alpha(a_0|1-t|,\tau|t|)$; zero variance follows directly.
For $t<0$ this is at least $a_0$, and for $t>1$ it is at least
$z_{1-\alpha/2}\tau$. Both endpoints are achievable. On $[0,1]$, center $t\ip LZ$ has exactly the asserted maximum bias and variance and attains the quantile bound. Continuity, including the zero-variance limit, gives a minimum.
Finally $F(h;\mu,v)\le\Phi((h-\mu)/v)$ implies
$Q_\alpha(\mu,v)\ge\mu+z_{1-\alpha}v$.
If $a_0\le z_{1-\alpha}\tau$, every $t\in[0,1]$ requires radius at least $a_0$, attained at $t=0$.
\end{proof}
At cap one, $M=\tau\sqrt p$ and $\alpha=0.05$, the exact affine length is $2\tau$.
The scalar bounded-normal-mean mechanism is classical \citep{stark1992}; the tensor reduction is proved above.

\subsection{Noise envelopes and approximate ranks}
Define $\Theta^{\varepsilon_A}=\{A:\exists B\in\Theta(M,\bar r),\ \norm{A-B}_F\le\varepsilon_A\}$ for supplied $\varepsilon_A\ge0$.
\begin{proposition}[Envelope and approximation calibration]\label{prop:envelope}
Suppose $d_i=a_i+e_i$ and, uniformly over the asserted class,
$\PP(\max_i|e_i|\le\lambda_0)\ge1-\alpha$.
For any measurable $\Lambda\ge\lambda_0$, the threshold center has an honest radius
\[
 \norm w_\infty C_{\mathrm{diag}} M^{2/K}(2\Lambda)^{1-2/K}+\varepsilon_A
\]
on $\Theta^{\varepsilon_A}$. The structural radius is $h_0+\varepsilon_A$.
Data-dependent selection between these two intervals is valid.
A direct Gaussian alternative needs its own calibration and deterministic selection or joint error allocation.
\end{proposition}
\noindent\emph{Roadmap.} Use one event for every allowed threshold and a residual Lipschitz bound.
\begin{proof}
On the envelope event, Lemma~\ref{lem:soft} holds simultaneously for all $\Lambda\ge\lambda_0$, proving the assertion at $\varepsilon_A=0$.
For fixed $e$, the residual map
$f_e(x)=\soft_\Lambda(x+e)-x=e-\operatorname{clip}(x+e,[-\Lambda,\Lambda])$
is one-Lipschitz. For $a_i=\ip A{V_i}$ and $b_i=\ip B{V_i}$,
\[
 |\sum_iw_i\{f_{e_i}(a_i)-f_{e_i}(b_i)\}|
 \le\norm w_2\norm{a-b}_2\le\norm{A-B}_F.
\]
The last step is Bessel. Add the threshold error at $B$.
Also $|\psi(A)-\psi(B)|\le\norm L_F\norm{A-B}_F\le\varepsilon_A$.
Both intervals cover on the same event, proving the selection assertion.
\end{proof}
For iid Gaussian errors the exact threshold is
$\tau\Phi^{-1}(\{1+(1-\alpha)^{1/m}\}/2)$,
by multiplying the $m$ probabilities $2\Phi(\lambda_0/\tau)-1$.
Dependent Gaussian margins with variances at most $\tau^2$ admit
$\tau\Phi^{-1}(1-\alpha/(2m))$ by Bonferroni.
Marginal bounds $\E e^{te_i}\le e^{t^2v^2/2}$ give
$v\sqrt{2\log(2m/\alpha)}$ by Chernoff and a union bound.
Bounds $\E|e_i|^a\le v^a$ give $v(m/\alpha)^{1/a}$ by Markov and a union bound.
Independence is unnecessary for these three union-bound calibrations.

Writing $\mathcal D_L A=(\ip A{V_i})_i$, the same proof gives
\[
 |\widehat\psi-\psi(A)|\le\inf_B\left\{
 \norm w_\infty C_{\mathrm{diag},B}\norm B_F^{2/K}(2\Lambda)^{1-2/K}
 +\norm{\mathcal D_L(A-B)}_2\right\}.
\]
The infimum ranges over tensors with their indicated valid budget $C_{\mathrm{diag},B}$.
Its unknown value is not an adaptive confidence radius.

For Tucker approximation let $P_k$ be any leading mode singular-space projector of dimension $\min(\bar r_k,p_k)$ and put $B=(\bigotimes_kP_k)A$.
The lifted projections commute. The identity
$I-\prod_kP_k=\sum_k(\prod_{\ell<k}P_\ell)(I-P_k)$
has summands with mutually orthogonal ranges.
Their squared residual norms add, each at most $\norm{(I-P_k)A}_F^2$.
Thus
\[
 \norm{A-B}_F^2\le\sum_k\sum_{j>\bar r_k}\sigma_j(A_{(k)})^2,
 \qquad \norm B_F\le\norm A_F.
\]
This holds without a spectral gap, including ties.
The distance-class proposition also covers TT approximation.

\subsection{Finite variance and bounded row corruption}
\begin{proposition}[Replicate median calibration]\label{prop:mom}
Let clean tensor replicates be independent with common mean $A$, and every prescribed contraction have variance at most known $v^2$. Allow at most $c$ adversarial replacements of entire replicates. Coordinates within a replicate may be dependent. Choose an odd block count $b$ with
\[
 b\ge8c,\qquad b\ge32\log(m/\alpha),\qquad n\ge2b,
 \qquad k=\lfloor n/b\rfloor.
\]
Partition the first $bk$ replicates into fixed disjoint blocks.
For coordinatewise medians $d_i$ of the block means,
\[
 \PP(\max_i|d_i-a_i|>2v/\sqrt k)\le\alpha.
\]
Thus Proposition~\ref{prop:envelope} applies. For the smallest admissible odd $b$,
the threshold has order
\[
 O\{v\sqrt{c/n+\log(m/\alpha)/n}\}.
\]
\end{proposition}
\noindent\emph{Roadmap.} Bound clean-block errors, count corrupted blocks deterministically and concentrate the original bad-block count.
\begin{proof}
Chebyshev bounds a clean block's error probability at $2v/\sqrt k$ by $1/4$.
For a fixed coordinate let $I_j$ be its original clean-block error indicators.
They are independent and $\sum_j\E I_j\le b/4$.
At most $c$ summaries change even for an adaptive adversary.
If the final median errs, $\sum_j I_j\ge b/2-c\ge3b/8$.
For a Bernoulli $I$, $h(t)=\log\E e^{t(I-\E I)}$ satisfies
$h(0)=h'(0)=0$ and $h''(t)\le1/4$, the variance of a tilted Bernoulli.
Hence $h(t)\le t^2/8$.
Independence and Markov give
\[
 \PP\{\sum_j(I_j-\E I_j)\ge b/8\}
 \le e^{-tb/8+bt^2/8}=e^{-b/32}\quad(t=1/2).
\]
A union bound over $m$ coordinates proves the claim. Since $k\ge n/(2b)$,
the rate follows, including the fixed rounding constants in $b$.
\end{proof}
Median-of-means calibration is classical \citep{devroye2016}.
For independent Bernoulli replacements of probability $\epsilon$, let
$c=\lceil2n\epsilon+2\log(1/\delta)\rceil$.
Their count satisfies $\E2^N=(1+\epsilon)^n\le e^{n\epsilon}$, giving
$\PP(N>c)\le e^{n\epsilon-c\log2}\le\delta$.
With envelope error budget $\alpha_0$, coverage is at least
$1-\alpha_0-\delta$ when the block conditions hold.
These sufficient corruption conditions are not a sharp breakdown or contamination-minimax result.

\subsection{Unknown Gaussian variance}
\begin{proposition}[Independent residual calibration]\label{prop:variance}
Retain $n\ge2$ iid Gaussian full-tensor replicates with unknown noise standard deviation $\sigma>0$.
Let $d_{ji}=\ip{W_j}{V_i}$, $\bar d_i=n^{-1}\sum_jd_{ji}$, $\nu=m(n-1)$, and
\[
 S^2=\nu^{-1}\sum_i\sum_j(d_{ji}-\bar d_i)^2,\qquad
 \sigma_U=S\sqrt{\nu/\chi^2_{\nu,\delta}},\quad 0<\delta<\alpha.
\]
The denominator is the lower $\delta$ chi-square quantile.
Compute the original three intervals with $\tau_U=\sigma_U/\sqrt n$ and
$\alpha_0=\alpha-\delta$, choosing the smallest using only $S$ and supplied quantities.
The selected interval is honest at level $1-\alpha$.
For fixed $\delta$ and $\nu\to\infty$, $\sigma_U/\sigma\to1$ in probability,
and the nonlinear expected-length order is preserved.
\end{proposition}
\noindent\emph{Roadmap.} Orthogonally separate means and residuals, then condition on a conservative scale estimate.
\begin{proof}
An orthogonal transformation in the replicate index separates each sample mean from its $n-1$ residual coordinates. Joint Gaussianity and loading orthonormality give independence of all residual and mean coordinates.
Therefore $\nu S^2/\sigma^2\sim\chi^2_\nu$, independent of the vector of means, and
$\PP(\sigma_U\ge\sigma)=1-\delta$.
Conditional on $S$ with $\sigma_U\ge\sigma$, both stochastic intervals have failure at most $\alpha_0$ and the selected index is fixed. The structural interval covers surely. Integrating conditional failure and adding $\delta$ proves honesty.

The normalized chi-square has mean one and variance $2/\nu$.
Chebyshev gives convergence to one and convergence of its fixed-$\delta$ quantile to one. Thus $\sigma_U/\sigma\to1$ and
$\E(\sigma_U/\sigma)^2=\nu/\chi^2_{\nu,\delta}\to1$.
The second-moment bound gives uniform integrability for powers at most one.
Since the nonlinear radius is proportional to $\tau_U^{1-2/K}$, its expected order is preserved.
\end{proof}
Retained residuals are essential to this construction; it is not an estimate of variance from $Z$ alone.

\subsection{Confidence levels and covariance}
\begin{proposition}[Calibrated lower bounds]\label{prop:levels}
The fixed-confidence rates of Theorems~\ref{thm:jointlower}, \ref{thm:unfolding}
and Proposition~\ref{prop:diagonal_minimax} extend to every $0<\alpha<1/2$
after changing constants.
If known Gaussian covariance satisfies $\Sigma\succeq s_0^2I$ with $s_0>0$,
all-procedure isotropic lower bounds transfer with $\tau=s_0$.
\end{proposition}
\noindent\emph{Roadmap.} Tune the mixture distance below $1-2\alpha$; compare covariance models by adding independent noise.
\begin{proof}
Fix $0<\varepsilon_{\mathrm{TV}}<1-2\alpha$. For the coordinate mixture use
$a=\min\{M,\tau\sqrt{\log(1+4\varepsilon_{\mathrm{TV}}^2p)}\}$, giving $\chi^2\le4\varepsilon_{\mathrm{TV}}^2$.
For the unfolding mixture use $\theta=\min(M,c_{\varepsilon_{\mathrm{TV}}}\tau\sqrt p)$ with
$c_{\varepsilon_{\mathrm{TV}}}^2=\{1-(1+4\varepsilon_{\mathrm{TV}}^2)^{-2}\}/2$.
Its second-moment bound is at most $1+4\varepsilon_{\mathrm{TV}}^2$.
For the diagonal support mixture use
\[
 b_s=\min\{Ms^{-K/2},\
 \tau\sqrt{\log[1+\log(1+4\varepsilon_{\mathrm{TV}}^2)p/s^2]}\}.
\]
The factorial-moment bound again gives second moment at most $1+4\varepsilon_{\mathrm{TV}}^2$.
Total variation is at most $\varepsilon_{\mathrm{TV}}$ in each case; the interval lemma has positive factor $1-2\alpha-\varepsilon_{\mathrm{TV}}$. The fixed-$\alpha$ rates are unchanged.

Add independent $N(0,\Sigma-s_0^2I)$ noise to $N(A,s_0^2I)$.
Applying any $\Sigma$-procedure to the result gives an isotropic randomized procedure with the same coverage and length distributions.
The isotropic all-procedure lower bound therefore applies.
\end{proof}
Upper envelope calibration uses the marginal variances
$\operatorname{vec}(V_i)^\top\Sigma\operatorname{vec}(V_i)$;
the direct radius uses
$z_{1-\alpha/2}\sqrt{\operatorname{vec}(L)^\top\Sigma\operatorname{vec}(L)}$.
A lower covariance bound is needed for a universal positive noise lower scale:
if $\Sigma=\tau^2(I-\operatorname{vec}(L)\operatorname{vec}(L)^\top)$,
the target is observed without noise.
Cross terms $B^\top\Sigma^{-1}C$ cannot in general be ordered by a covariance bound; the randomization proof avoids that invalid step.


\section{Proofs of the sharp rates and computational statements}\label{app:gap_proofs}

\subsection{Proof of Theorem~\ref{thm:sparse_harmonic}}
\noindent\emph{Roadmap.} Embed a harmonic rank-one profile on a random support; bound its full tensor likelihood using support overlap.
\begin{proof}
For a size-$m$ support $S$, assign values $j^{-1/3}$, $j=1,\ldots,m$, in increasing coordinate order and zeros elsewhere; call this vector $x_S$. Set $A_S=b_mx_S^{\otimes3}$. Then
\[
 \norm{A_S}_F=b_m\left(\sum_{j=1}^m j^{-2/3}\right)^{3/2}
 \le b_m\sqrt{27m}\le M,\qquad \psi(A_S)=b_mH_m/\sqrt p.
\]
The sum bound follows by integrating $x^{-2/3}$.
All alternatives have rank one and the same target.
For $H=|S\cap T|$, Cauchy--Schwarz and the decreasing profile give
\[
 0\le\ip{x_S}{x_T}
 \le\left(\sum_{i\in S\cap T}x_{S,i}^2
           \sum_{i\in S\cap T}x_{T,i}^2\right)^{1/2}
 \le\sum_{j=1}^Hj^{-2/3}\le3H^{1/3}.
\]
At $H=0$ every term is zero. Therefore
$\ip{A_S}{A_T}\le27b_m^2H$.
The full Gaussian likelihood identity and the hypergeometric bound already proved in Proposition~\ref{prop:diagonal_minimax} imply
\[
 1+\chi^2(Q,P_0)\le\E e^{27b_m^2H/\tau^2}
 \le\exp\{(m^2/p)(e^{27b_m^2/\tau^2}-1)\}\le e^{1/4}.
\]
Only the hypergeometric inequality is reused, not the diagonal likelihood.
Lemma~\ref{lem:mixture} proves \eqref{eq:sparse_harmonic}.

For $p\ge256$, $m=\lfloor p^{1/4}\rfloor\ge p^{1/4}/2$,
so $H_m\ge\log m\ge(\log p)/8$.
Also $p/(4m^2)\ge\sqrt p/4\ge p^{1/4}$, and
\[
 \tfrac14\log p\le\log(1+p/(4m^2))\le\log p\le\sqrt p\le p/m.
\]
For $p\ge4$, $\log p\le\sqrt p$ follows from positivity at four and the nonnegative derivative of $\sqrt p-\log p$.
Thus the noise term is the minimum in $b_m$ and is at least
$\tau\sqrt{\log p}/(2\sqrt{27})$. Multiplication proves the stated constant.
\end{proof}

\subsection{Proof of Theorem~\ref{thm:full_minimax}}
\noindent\emph{Roadmap.} Verify measurable extrema and coverage, then bound differences of feasible tensors using doubled caps.
\begin{proof}
The parameter class is compact: its norm ball is bounded and closed, and rank constraints are closed polynomial conditions on unfolding minors. Each finite-net maximum is continuous. Hence every nonempty feasible section is compact and its functional extrema are attained.

The feasible-data domain is closed. Indeed, if $Z_j\to Z$ and $B_j$ are feasible, compactness supplies a subsequence $B_{j_k}\to B\in\Theta$; continuity preserves every constraint.
For any real $u$, the set of $Z$ admitting a feasible $B$ with $\psi(B)\ge u$ is also closed by this argument, as is the analogous set with $\psi(B)\le u$.
The extrema on the feasible-data domain, extended by zero elsewhere, are therefore Borel functions.
The noise event at error level $\alpha$ makes $A$ feasible, uniformly with probability at least $1-\alpha$. On that event $C_F$ covers the target.

For $B,B'\in{\cal F}(Z)$, put $\Delta=B-B'$.
Matrix-rank subadditivity gives the doubled unfolding caps; coordinate restriction cannot increase them. Also
$N_S(\Delta)\le N_S(Z-B)+N_S(Z-B')\le2\lambda_s$,
so $\rho(\Delta_S)\le4\lambda_s$.
The diagonal bound and injective-norm comparison with $J_2$ give
\[
 \sum_{i\in S}|\Delta_{i,\ldots,i}|^{2/K}
 \le J_2^{1-1/K}\norm{\Delta_S}_F^{2/K}
 \le J_2\rho(\Delta_S)^{2/K}.
\]
Taking the $s$ largest diagonal entries yields
$|\Delta|_{(s)}\le4J_2^{K/2}\lambda_s s^{-K/2}$.
Summation proves $|\psi(B)-\psi(B')|\le W_\alpha$.
Apply it to the attained extrema; the structural range gives the other bound $2h_0$. Empty sections have length zero.

For fixed confidence, $\lambda_s\le C_{K,\alpha}\tau\sqrt{s\log p}$ for large $p$.
The sum of $s^{-(K-1)/2}$ is bounded for $K\ge4$ and at most $1+\log p$ for $K=3$.
This proves the worst-case upper bounds.
Theorem~\ref{thm:jointlower} matches the higher-order bound and
Theorem~\ref{thm:sparse_harmonic} matches the cubic bound, already at zero and within the rank-one subclass. Since $\calL_0\le\calL_{\max}$, both rates follow.
\end{proof}

\subsection{Proof of Proposition~\ref{prop:mse_sharp}}
\noindent\emph{Roadmap.} Use a smaller failure budget and bounded target range for the upper risk; reduce estimation to common-target testing for the lower risk.
\begin{proof}
For $p\ge2$, $\alpha_p=p^{-2}$ still gives $W_{\alpha_p}=O(r_K(p))$.
The midpoint is in $[-h_0,h_0]$, including the empty-set convention.
On coverage its error is at most $W_{\alpha_p}/2$; otherwise the error is at most $2h_0$. Thus
\[
 \sup_A\E_A(\widehat\psi-\psi(A))^2
 \le W_{\alpha_p}^2/4+4h_0^2p^{-2}=O(r_K(p)^2),
\]
since $h_0=\tau\sqrt{s_{\rm cap}}$.
For a common-target prior with target $\delta$, test whether an estimator is at least $\delta/2$.
If its worst-case squared risk is $R$, Markov bounds the total null-plus-mixture testing error by $8R/\delta^2$.
Total variation bounds this error below by $1-\TV(P_0,Q)$.
Hence $R\ge\delta^2(1-\TV)/8$.
The coordinate and sparse harmonic mixtures have bounded total variation and targets of the required respective orders.
\end{proof}

\subsection{Proof of Proposition~\ref{prop:truncated}}
\noindent\emph{Roadmap.} Combine leading-entry scan bounds with a quasi-norm tail bound and an explicit grid-net construction.
\begin{proof}
On acceptance and the common noise event,
$|a|_{(s)}\le4J^{K/2}\lambda_s s^{-K/2}$ for $s\le t$.
The remaining magnitudes are at most $a_t$, so
\[
 \sum_{s>t}|a|_{(s)}
 \le a_t^{1-q}\sum_{s>t}|a|_{(s)}^q\le a_t^{1-q}B_q.
\]
The tail vanishes for $t=p$. This proves the radius.
On rejection the structural interval covers surely.
Intersection gives honesty by error allocation and length at most $\ell_*$.
At zero rejection has probability at most $\gamma$, giving the stated expectation bound.

For the rate, substitute $\lambda_s\lesssim\tau\sqrt{s\log p}$ and
$B_q=J^{1-1/K}\tau^{2/K}p^{1/K}$.
The tail becomes a constant times
$\tau p^{1/K-1/2}(\log p)^{(1-2/K)/2}t^{-(K-1)(1-2/K)/2}$,
as asserted.

For a computable net, split each coordinate of $[-1,1]^s$ into
$N=\lceil4K\sqrt s\rceil$ equal intervals and normalize every nonzero product-grid vector.
Rounding a unit vector incurs error at most $\sqrt s/N\le1/(4K)$.
The rounded vector is nonzero, and normalization adds at most the original error, giving a $1/(2K)$-net.
Its size is at most $(N+1)^s\le(3+4K\sqrt s)^s$.
Use log cardinality $Ks\log(3+4K\sqrt s)$ in the Gaussian union bound.
Since $s\le p$, this still gives $\lambda_s\lesssim_K\tau\sqrt{s\log p}$.
Each contraction costs $O(s^K)$ operations, proving the complexity bound.
All completed levels cover on one event; adaptive continuation and intersection add no scan error. The separately calibrated $C_*$ is intersected only once.
\end{proof}

\subsection{Proof of Theorem~\ref{thm:low_degree}}
\noindent\emph{Roadmap.} Expand in Hermite polynomials, bound overlap moments and control the possible mean shift of a low-degree estimator by its null variance.
\begin{proof}
Scale by $\tau$ and let $h_\nu$ be normalized multivariate probabilists' Hermite polynomials over the $p^K$ coordinates.
The generating identity
$e^{tz-t^2/2}=\sum_{j\ge0}H_j(z)t^j/j!$
gives $\E H_j(g)H_k(g)=j!\mathbf1_{\{j=k\}}$:
multiply two generating functions and compare coefficients of their expectation $e^{tu}$.
Coordinate independence gives multivariate orthonormality.
The shift calculation similarly gives
$\E h_\nu(g+a)=a^\nu/\sqrt{\nu!}$.
The coefficient identities follow by differentiation at zero of Gaussian entire moment-generating functions.

For the finite support prior, the degree-$d$ coefficient sum is
\[
 \sum_{|\nu|=d}\{\E_Qh_\nu(Z/\tau)\}^2
 =\frac1{d!}\E_{S,T}
      \left(\frac{\ip{A_S}{A_T}}{\tau^2}\right)^d
 =\frac{\lambda^{2d}}{d!s^{Kd}}\E H^{Kd},
\]
by the multinomial theorem, where $H=|S\cap T|$.
This uses the full likelihood.

Let $\mathrm S(m,j)$ count partitions of an $m$-element set into $j$ nonempty unlabeled sets.
Counting maps according to their number of image elements gives
$H^m=\sum_{j=1}^m\mathrm S(m,j)(H)_j$.
There are at most $j^m$ labeled maps, so
$\mathrm S(m,j)\le j^m/j!$.
For $j\le s$, the hypergeometric factorial moments obey
\[
 \E(H)_j=(s)_j^2/(p)_j\le(s^2/p)^j=\mu^j,
\]
using $(s-\ell)/(p-\ell)\le s/p$ and $s-\ell\le s$ factor by factor.
For $j>s$ the moment is zero. Therefore
\[
 \E H^m\le\sum_{j=1}^m j^m\mu^j/j!
 \le m^m(e^\mu-1)\le\mu e^\mu m^m.
\]
The centered projection omits degree zero. Summing the remaining degrees,
\[
 \chi^2_{\le \mathsf d}
 \le\mu e^\mu\sum_{d=1}^{\mathsf d}
      \frac{\lambda^{2d}(Kd)^{Kd}}{d!s^{Kd}}
 \le\mu e^\mu\sum_{d=1}^{\mathsf d} B_{\mathsf d}^d
 \le\mu e^\mu B_{\mathsf d}/(1-B_{\mathsf d})=\eta_{\mathsf d}^2.
\]

For a degree-$\mathsf d$ estimator with worst-case squared risk $R$ on zero and the alternatives,
$|\E_0T|\le\sqrt R$ and $|\E_QT-\delta|\le\sqrt R$.
Orthonormal projection and Cauchy--Schwarz give
\[
 |\E_QT-\E_0T|\le
 \sqrt{\operatorname{Var}_0T}\sqrt{\chi^2_{\le \mathsf d}}\le\sqrt R\,\eta_{\mathsf d}.
\]
Thus $\delta\le(2+\eta_{\mathsf d})\sqrt R$, proving the risk lower bound.
For the stated exponents,
$\mu=O(p^{2a-1})\to0$ and $B_{\mathsf d}=O(p^{1-aK+bK})\to0$;
floor factors remain bounded for large $p$.
The target is $\delta=\tau s^{1-K/2}$.
Since $a(K-2)<1$, its square is polynomially larger than
$\tau^2/p$ times either logarithmic factor of the unrestricted minimax risk.
The alternatives have norm $M$ and rank one in every unfolding, so belong to either positive cap class.
\end{proof}

\subsection{Proof of Proposition~\ref{prop:ci_detection}}
\noindent\emph{Roadmap.} Combine mixture coverage with null coverage and Markov's inequality for length.
\begin{proof}
Mixture coverage gives $Q(\delta\notin C)\le\alpha$.
If $0\in C$ and $|C|<\delta$, the interval cannot contain $\delta$.
Consequently
\[
 P_0(\delta\in C)\le P_0(0\notin C)+P_0(|C|\ge\delta)
 \le\alpha+\E_0|C|/\delta.
\]
Add the errors. Polynomial-time interval evaluation and membership evaluation yield a polynomial-time test, proving the conditional implication.
No unrestricted testing-hardness conclusion is deduced from the low-degree squared-risk bound.
\end{proof}

\section{Proofs for norm budgets and loading geometry}\label{app:norm_loading}

\subsection{Proof of Theorem~\ref{thm:norm_diameter}}
\noindent\emph{Roadmap.} Bound each ordered diagonal difference twice, then take the minimum before summing.
\begin{proof}
The coverage, compactness, and measurable-extrema argument of Theorem~\ref{thm:full_minimax} applies for every finite $M>0$. For two feasible tensors put $\Delta=B-B'$. Then $\norm\Delta_F\le2M$ and every constrained unfolding rank is at most twice its original cap. Writing $a_i=\Delta_{i,\ldots,i}$ and $q=2/K$, Theorem~\ref{thm:diagonal} gives
\[
 s|a|_{(s)}^q\le\sum_i|a_i|^q
       \le J_2^{1-1/K}(2M)^q.
\]
The norm argument and the residual-scan proof give, respectively,
\[
|a|_{(s)}\le P_2M s^{-K/2},\qquad
|a|_{(s)}\le Q_2\lambda_s(\alpha)s^{-K/2}.
\]
Both inequalities hold for the same ordered entry. Summing their minimum bounds $|\psi(B)-\psi(B')|$ by \eqref{eq:norm_diameter}. The structural range gives $2h_0$. Empty feasible sets produce a singleton and therefore length zero.
\end{proof}

\begin{lemma}[A norm-sensitive harmonic sum]\label{lem:norm_sum}
For $p\ge1$, $x>0$ and $K\ge3$, put
\[
 S_{K,p}(x)=\sum_{s=1}^p\min\{x s^{-K/2},s^{-(K-1)/2}\}.
\]
Then $S_{3,p}(x)\le4g_p(x)$, and $S_{K,p}(x)\le3\min(x,1)$ for $K\ge4$.
For fixed $K$ and $0<\alpha<1$,
$U_{K,p}(M;\alpha)\le C_{K,\bar r,\alpha}\mathfrak r_K(p,M,\tau)$ uniformly in $p,M,\tau$.
\end{lemma}
\noindent\emph{Roadmap.} Split the cubic sum at $s=x^2$ and use summability in higher order.
\begin{proof}
For $K\ge3$, $S_{K,p}(x)\le x\sum_{s\ge1}s^{-3/2}\le3x$. For $K\ge4$ the second term gives $S_{K,p}(x)\le\sum_{s\ge1}s^{-3/2}\le3$.

For $K=3$, $S_{3,p}(x)\le H_p\le\ell_p$. If $x\le1$, the bound $3x$ suffices for $g_p(x)=x$. If $1<x<\sqrt p$, put $n=\lfloor x^2\rfloor$. Since $n\ge x^2/2$,
\[
 S_{3,p}(x)\le H_n+x\sum_{s>n}s^{-3/2}
 \le1+2\log x+2x/\sqrt n
 \le4(1+\log x).
\]
If $x\ge\sqrt p$, the bound $H_p\le1+\log p\le2(1+\log x)$ gives the same conclusion. Together with $3x$ and $\ell_p$, this proves the first statement.

The threshold satisfies
\[
 \lambda_s(\alpha)\le c_{K,\alpha}\tau\sqrt{s\ell_p},\qquad
 c_{K,\alpha}=\{2[2+K\log(1+4K)+\log(2/\alpha)]\}^{1/2}.
\]
Indeed, bound $\log(ep/s)\le\ell_p$, $\log p\le\ell_p$, $s\ge1$, and $\ell_p\ge1$ in \eqref{eq:scan_threshold}. The inequality
$\min(au,bv)\le\max(a,b)\min(u,v)$ for nonnegative arguments shows
\[
 U_{K,p}(M;\alpha)\le
 \max(P_2,Q_2c_{K,\alpha})\frac{\tau\sqrt{\ell_p}}{\sqrt p}
 S_{K,p}(x).
\]
The established sum bounds prove the claim.
\end{proof}

\subsection{Proof of Theorem~\ref{thm:all_norm_rates}}
\noindent\emph{Roadmap.} Combine the new diameter with a coordinate mixture and a norm-tuned harmonic support mixture.
\begin{proof}
The upper bound follows from Theorem~\ref{thm:norm_diameter} and Lemma~\ref{lem:norm_sum}. Since the upper diameter is deterministic it controls both expected-length criteria.

For the lower bound put $a=\min\{M,\tau\sqrt{\log(1+p/4)}\}$ and mix uniformly over $aE_i$, $i=1,\ldots,p$. These rank-one alternatives have target $a/\sqrt p$. The full Gaussian likelihood identity gives
\[
 \chi^2(Q,P_0)=\frac{e^{a^2/\tau^2}-1}{p}\le\tfrac14,
 \qquad\TV(Q,P_0)\le\tfrac14.
\]
The common-target interval lemma implies
\begin{equation}\label{eq:norm_coordinate_lower}
 \E_0|C|\ge(3/4-2\alpha)\frac{\min\{M,\tau\sqrt{\log(1+p/4)}\}}{\sqrt p}.
\end{equation}
For $p\ge1$, $\log(1+p/4)\ge\ell_p/5$. To verify this, the difference is positive at one and its derivative is
$1/(p+4)-1/(5p)\ge0$. Thus \eqref{eq:norm_coordinate_lower} is at least a constant times
$\tau\sqrt{\ell_p}\min(x,1)/\sqrt p$. This proves the lower bound for every $K\ge4$, and for $K=3$ when $x\le1$.

For the remaining cubic case assume first $p\ge256$ and $x>1$, and take
\[
 u=\min(p^{1/4},x^2),\qquad m=\lfloor u\rfloor\ge1.
\]
Use the sparse harmonic mixture of Theorem~\ref{thm:sparse_harmonic}, with the original general-budget coefficient $b_m$. Since $m\le x^2$,
$M/\sqrt m\ge\tau\sqrt{\ell_p}$. Also $m\le p^{1/4}$ implies
\[
 \log(1+p/(4m^2))\ge\tfrac14\log p\ge\ell_p/5,
\]
where the first inequality holds for $p\ge256$ and the second for $\log p\ge4$. Hence
$b_m\ge\tau\sqrt{\ell_p}/\sqrt{135}$.
Furthermore $H_m\ge1$ and $H_m\ge\log(m+1)\ge\log u$.
Writing $t=\log x$, if $2t\le(\log p)/4$ then
$H_m\ge\max(1,2t)\ge(1+t)/2$. Otherwise
$H_m\ge(\log p)/4\ge\ell_p/5$.
In both cases
\[
 H_m\ge\tfrac15\min(1+\log x,\ell_p)
       =\tfrac15 g_p(x),
\]
where $1+\log x\le x$ for $x\ge1$ was used. Consequently
\begin{equation}\label{eq:norm_cubic_lower}
 \E_0|C|\ge\frac{c_{\alpha,d}}{5\sqrt{135}}
            \mathfrak r_3(p,M,\tau).
\end{equation}
All likelihoods here are for the full tensor; the support-mixture proof holds for the newly chosen $m$ and $b_m$ without modification.

For $p<256$, $\ell_p<7$ and $g_p(x)\le7\min(x,1)$, so the coordinate mixture suffices with a changed absolute constant. Finally $\mathcal L_0\le\mathcal L_{\max}$ completes the proof. Every lower alternative has rank one; higher supplied caps enter only the upper constants.
\end{proof}

\begin{lemma}[Residual fit and its noise tail]\label{lem:residual_fit}
The estimator in \eqref{eq:residual_fit} is Borel measurable. With $T_0=T(\tau G)$ and $\beta=1/8$,
\[
 \PP(T_0>t)\le\beta e^{-t^2/(2\tau^2)}\quad(t\ge0),
 \qquad \E T_0^2\le2\beta\tau^2.
\]
Put $z_K=\sum_{s=1}^{\infty}s^{-K/2}\le3$. Then
\begin{equation}\label{eq:residual_error}
 |\widehat\psi-\psi(A)|\le U_{K,p}(M;\beta)
              +\frac{Q_2z_K}{\sqrt p}(T_0+\tau/2),
 \qquad |\widehat\psi-\psi(A)|\le2h_0.
\end{equation}
\end{lemma}
\noindent\emph{Roadmap.} Use a Gaussian union bound with a common additive slack, then compare an approximate residual minimizer to the truth.
\begin{proof}
Each unit net contraction is one-Lipschitz in the Frobenius norm; finite maxima and the positive part preserve this property. Thus $T$ is continuous and one-Lipschitz. Density of $(B_j)$ implies
$d(Z)=\min_{B\in\Theta}T(Z-B)$, although the estimator only uses a positive tolerance. The function $d$ is one-Lipschitz as an infimum of functions with the same Lipschitz constant. The sets $\{Z:T(Z-B_j)\le d(Z)+\tau\}$ are Borel, and their union is the entire observation space. Selecting the first such index is measurable. Hence so is $\widehat\psi$.

For $t\ge0$, apply the same finite Gaussian union bound as for the original scan at thresholds $\lambda_s(\beta)+t$. Since
\[
 e^{-(\lambda_s+t)^2/(2\tau^2)}
 \le e^{-\lambda_s^2/(2\tau^2)}e^{-t^2/(2\tau^2)},
\]
the total probability is at most $\beta e^{-t^2/(2\tau^2)}$. Integrating the tail of $T_0^2$ proves its moment bound.

For $\widehat A=B_{j(Z)}$, comparison with $A$ gives
$T(Z-\widehat A)\le d(Z)+\tau\le T_0+\tau$. Therefore for every support,
\[
 N_S(\widehat A-A)\le2\lambda_s(\beta)+2T_0+\tau.
\]
The net comparison and doubled-rank argument bound its ordered diagonal entry by
$Q_2(\lambda_s(\beta)+T_0+\tau/2)s^{-K/2}$.
It is also bounded by $P_2M s^{-K/2}$. Use
$\min(a,b+c)\le\min(a,b)+c$, sum, and divide by $\sqrt p$ to obtain \eqref{eq:residual_error}. Both tensors belong to the original cap class, so their targets lie in $[-h_0,h_0]$, giving the second bound.
\end{proof}

\subsection{Proof of Theorem~\ref{thm:all_norm_mse}}
\noindent\emph{Roadmap.} Control residual-fit tails without an $M$-dependent failure event and reuse the same full-data mixtures for the lower bound.
\begin{proof}
Lemma~\ref{lem:residual_fit} and $(a+b)^2\le2a^2+2b^2$ give
\begin{equation}\label{eq:norm_mse_finite}
 \sup_A\E_A(\widehat\psi-\psi(A))^2
 \le\min\left\{4h_0^2,
 2U_{K,p}(M;\beta)^2
 +(8\beta+1)Q_2^2z_K^2\frac{\tau^2}{p}\right\}.
\end{equation}
If $x\le1$, use $h_0^2\le s_{\rm cap}M^2/p$ and
$\mathfrak r_K=M/\sqrt p$. If $x>1$, then
$\mathfrak r_K\ge\tau\sqrt{\ell_p}/\sqrt p\ge\tau/\sqrt p$; Lemma~\ref{lem:norm_sum} absorbs both terms in the second bound into a constant times $\mathfrak r_K^2$.

The lower constructions in Theorem~\ref{thm:all_norm_rates} have common target $\delta\ge c\mathfrak r_K$ and total variation bounded by a constant below one, independently of $p,M,\tau$. For any estimator with worst-case risk $R$, the test $\mathbf1\{\widetilde\psi\ge\delta/2\}$ has total null-plus-mixture error at most $8R/\delta^2$ by Markov's inequality. Testing error is at least $1-\TV(P_0,Q)$. Thus
$R\ge\delta^2(1-\TV)/8\gtrsim\mathfrak r_K^2$.
\end{proof}

\subsection{Proof of Theorem~\ref{thm:weighted_bounds}}
\noindent\emph{Roadmap.} Rearrange weights against ordered diagonal bounds; equalize the target of a hard mixture by a contraction in one mode.
\begin{proof}
Extend each prescribed orthonormal loading family to an orthonormal mode basis. Mode transformations preserve Gaussian isotropy and every constrained unfolding rank; restriction to the loading block cannot increase rank or Frobenius norm. Apply residual projection to that $m^K$ block. Every pair of feasible block tensors satisfies the ordered-entry bound from Theorem~\ref{thm:norm_diameter}, with $p=m$. The rearrangement inequality gives
\[
 |\psi(B)-\psi(B')|
 \le\sum_{s=1}^m w_s^*|\operatorname{diag}(B-B')|_{(s)}\le U_w.
\]
The structural bound and the earlier measurable-extrema argument give the other assertions about this interval.

For the lower bound fix $d\le m$ and work in the coordinates corresponding to the $d$ largest absolute weights. These weights are nonzero. Take whichever coordinate or nonnegative cubic harmonic mixture supplies the lower rate in dimension $d$ in Theorem~\ref{thm:all_norm_rates}. Its common normalized diagonal target is at least $c\mathfrak r_K(d,M,\tau)$. In the first tensor mode apply the diagonal linear map with multipliers $w_d^*/w_i$ on these coordinates. Its operator norm is at most one, so every alternative remains rank one with Frobenius norm at most $M$. Its weighted target is now the same for every alternative and equals $w_d^*\sqrt d$ times the original normalized target.

For each pair of alternatives, the first-mode factor inner product becomes
\[
 \sum_i (w_d^*/w_i)^2 u_i v_i\le\sum_i u_i v_i,
\]
because the original factors are nonnegative. The other factor inner products are unchanged and nonnegative. Thus the full tensor inner product does not increase. The Gaussian likelihood second moment is no larger than that of the original hard mixture. Embedding these alternatives in the full ambient tensor adds only independent mean-zero noise coordinates and leaves the likelihood second moment unchanged. The common-target interval lemma therefore gives
$\E_0|C|\ge c_{K,\alpha}w_d^*\sqrt d\,\mathfrak r_K(d,M,\tau)$.
Maximize over the finite set of $d$ to obtain \eqref{eq:weighted_lower}.
\end{proof}

\subsection{Proof of Proposition~\ref{prop:flat_loading}}
\noindent\emph{Roadmap.} Rotate to the loading block, absorb known signs, and apply the matching dimension-$m$ theorem.
\begin{proof}
Replace each first-mode loading vector by $\epsilon_i v_{1i}$, preserving orthonormality. The loading now has equal positive weights $m^{-1/2}$. The reduced $m^K$ Gaussian tensor has the same norm and cap constraints, so Theorems~\ref{thm:all_norm_rates} and \ref{thm:all_norm_mse} give the upper bounds. For the lower bounds embed their rank-one alternatives in the loading block and set all other mean entries to zero. The extra independent noise carries no parameter information; the full likelihood second moments equal the reduced ones. Equivalently, the interval lower bound follows by taking $d=m$ in Theorem~\ref{thm:weighted_bounds}, and the same mixtures give the squared-risk lower bound. This includes $m=1$ and arbitrary ambient dimensions at least $m$.
\end{proof}

\section{Proofs for robustness, loading transfer, and adaptation}\label{app:robust_adaptive}

\subsection{Proof of Theorem~\ref{thm:sg_full}}
\noindent\emph{Roadmap.} Identify the tail inequalities actually used by the upper proofs, then use a Gaussian subfamily for family-wise minimax lower bounds.
\begin{proof}
For every unit net direction, Chernoff applied to \eqref{eq:joint_sg} gives
$\PP(|\ip UE|>t)\le2e^{-t^2/(2v^2)}$.
The finite union bound proving scan coverage and the additive-slack bound proving Lemma~\ref{lem:residual_fit} use no independence between directions. Both therefore hold with $v$ replacing $\tau$. All subsequent diameter and squared-error inequalities are deterministic functions of these events and moments. Orthogonal changes of coordinates preserve \eqref{eq:joint_sg}. The weighted proofs apply unchanged. The unknown-law family includes $E\sim N(0,v^2I)$. Honesty over the family implies honesty over that subexperiment, and the worst-case risk or expected length over the family is at least its Gaussian value. The Gaussian lower bounds give the matching order.
\end{proof}

\subsection{Proof of Theorem~\ref{thm:mom_full}}
\noindent\emph{Roadmap.} Control every median by a binomial tail, obtain a common excess-tail moment, and reuse the ordered-difference argument.
\begin{proof}
For a direction at level $s$, each block average has variance at most $\sigma^2/k_s$. Chebyshev gives a block-error probability at most $\sigma^2/(k_st^2)$. If the median errs by more than $t$, at least half its blocks err. Independence of blocks, followed by a union over their subsets, gives
\[
 \PP(|d_U-\ip UA|>\lambda_s^{\rm mom}+t)
 \le 2^{-b_s}(1+t/\lambda_s^{\rm mom})^{-b_s},\quad t\ge0.
\]
This follows by substituting $\lambda_s^{\rm mom}=4\sigma/\sqrt{k_s}$ into
$2^{b_s}\{\sigma^2/[k_s(\lambda_s^{\rm mom}+t)^2]\}^{b_s/2}$.
There are at most $e^{H_s}$ directions at level $s$. By the definition of $b_s$,
$\sum_s e^{H_s}2^{-b_s}\le\beta$.
Thus all constraints cover the true tensor with probability at least $1-\beta$.

Since $k_s\ge n/(2b_s)$,
$\lambda_s^{\rm mom}\le4\sqrt2\tau\sqrt{b_s}$.
Define
$T_0=\max_U(|d_U-\ip UA|-\lambda_{s(U)}^{\rm mom})_+$.
For $a\ge0$, the function $b\log(1+a/\sqrt b)$ is increasing for $b>0$: its derivative is
$\log(1+u)-u/[2(1+u)]\ge0$, $u=a/\sqrt b$.
Because every $b_s\ge9$, the union bound yields
\begin{equation}\label{eq:mom_excess}
 \PP(T_0>t)\le\beta(1+t/(12\sqrt2\tau))^{-9},\qquad
 \E T_0^2\le\frac{72}{7}\beta\tau^2.
\end{equation}
For the moment, integrate $2t$ times the tail and use
$\int_0^\infty u(1+u)^{-9}du=1/(8\cdot7)$.

For two feasible tensors every directional difference is at most $2\lambda_s^{\rm mom}$. The net comparison, doubled caps and norm bound prove exactly the previous norm-sensitive diameter. Medians are measurable finite order statistics. The extrema of their compact feasible sections are Borel by the same compact-section argument as before, now treating $(d_U)_U$ as the finite observation vector. The smallest-odd-integer rule gives
$b_s\le11+[H_s+\log(p/\beta)]/\log2\le C_{K,\beta}s\ell_p$.
The sum lemma therefore bounds the diameter by $C\mathfrak r_K$.

For estimation, $F$ is continuous and one-Lipschitz in $B$. Selecting the first dense-sequence element with value at most $\inf_jF(B_j)+\tau$ is measurable. Its comparison with $A$ gives directional differences at most $2\lambda_s^{\rm mom}+2T_0+\tau$. Lemma~\ref{lem:residual_fit}'s deterministic argument gives
\[
 |\widehat\psi-\psi(A)|\le U^{\rm mom}
       +Q_2z_K(T_0+\tau/2)/\sqrt p,
 \qquad |\widehat\psi-\psi(A)|\le2h_0.
\]
Equation~\eqref{eq:mom_excess} and the small-$M$ structural bound prove risk $O(\mathfrak r_K^2)$ exactly as in Theorem~\ref{thm:all_norm_mse}. For the lower bound restrict the replicate family to iid Gaussian errors with covariance $\sigma^2I$. Orthogonal transformation in the replicate index separates the sample mean, whose noise level is $\sigma/\sqrt n$, from mean-independent residuals. Thus the full replicate likelihood second moments equal the sample-mean ones. The earlier Gaussian lower bounds apply to any procedure using all replicates, including randomized ones. Finally, $H_s$ is increasing in $s$, so the replicate condition has the stated sufficient order $p$ for fixed $K,\beta$.
\end{proof}

\subsection{Proof of Proposition~\ref{prop:noise_limits}}
\noindent\emph{Roadmap.} Exhibit exact scalar recovery and a correlated direction with inflated variance.
\begin{proof}
If $|A|<v$, then $\operatorname{sign}(A+v\xi)=\xi$. The estimator
$Z-v\operatorname{sign}(Z)$ is identically $A$, proving (i). The Rademacher moment-generating function obeys $\cosh(tv)\le e^{t^2v^2/2}$.
For (ii), take the unit product tensor with every entry $p^{-K/2}$. Its contraction with $E$ is $v\xi p^{K/2}$, with variance $v^2p^K$. Each individual entry still has variance and sub-Gaussian proxy $v^2$.
\end{proof}

\subsection{Proof of Theorem~\ref{thm:approx_sharp}}
\noindent\emph{Roadmap.} Inflate scan constraints for a bounded approximation residual, then combine two nested lower-bound subclasses.
\begin{proof}
Write $A=B_0+R$, $B_0\in\Theta$, $\norm R_F\le\varepsilon_A$; compactness ensures such a decomposition exists. On the Gaussian scan event every residual contraction of $Z-B_0$ is at most $\lambda_s+\varepsilon_A$, so $B_0$ is feasible. Inflate its projected target interval by $\varepsilon_A$ on both sides. If the feasible set is empty use a singleton; this cannot occur on the coverage event. Since $\norm L_F=1$, the enlarged interval covers $\psi(A)$.
For feasible differences the diameter argument gives
\[
 |C|\le U_{K,p}(M;\alpha)
       +\{Q_2z_K/\sqrt p+2\}\varepsilon_A,
\]
using $\min(a,b+c)\le\min(a,b)+c$. It also gives $|C|\le2h_0+2\varepsilon_A$.
Choose between this procedure and the direct Gaussian interval using their deterministic length bounds. This proves the upper order
$\min\{\tau,\mathfrak r_K+\varepsilon_A\}$.

The class contains $\Theta$, giving lower bound $c\mathfrak r_K$ at zero. It also contains $\{tL:|t|\le\varepsilon_A\}$ because $0\in\Theta$. Comparing zero with $tL$ for $t=\min(\varepsilon_A,\tau/2)$ gives total variation at most
$\tfrac12\sqrt{e^{1/4}-1}<1/2$ and target difference $t$. Hence the lower bound is also $c_\alpha\min(\varepsilon_A,\tau)$. Finally
$\min(\tau,r+\varepsilon_A)\le r+\min(\tau,\varepsilon_A)\le2\max\{r,\min(\tau,\varepsilon_A)\}$
for every $r\ge0$. This establishes both criteria up to fixed constants.
\end{proof}

\subsection{Proof of Proposition~\ref{prop:growing_caps}}
\noindent\emph{Roadmap.} Embed sparse diagonal tensors of rank $s$, then compare with the full norm ball.
\begin{proof}
For a uniform size-$s$ support $S$, set
$A_S=b\sum_{i\in S}E_i$ with
\[
 b=\min\{M/\sqrt s,\tau\sqrt{\log(1+p/(4s^2))}\}.
\]
All constrained ranks are $s$, its norm is $b\sqrt s\le M$, and its target is $sb/\sqrt p$. The full tensor inner product is $b^2|S\cap T|$; the hypergeometric factorial-moment inequality gives
$1+\chi^2(Q,P_0)\le\exp\{(s^2/p)(e^{b^2/\tau^2}-1)\}\le e^{1/4}$.
The common-target lemma proves the lower bound for every $s\le r_*$. Maximize over $s$.
If there are no rank restrictions, the line $\{tL:|t|\le M\}$ belongs to the Frobenius ball. The two-point argument with $t=\min(M,\tau/2)$ gives lower order $\min(M,\tau)$. The direct interval and $[-M,M]$, chosen deterministically by radius, give the matching upper bound for every mean in that ball. No asymptotic fixed-cap argument applies to this enlarged class.
\end{proof}

\subsection{Proof of Proposition~\ref{prop:loading_transfer}}
\noindent\emph{Roadmap.} Add an independently calibrated error event, without assuming stochastic independence.
\begin{proof}
On the intersection of $\ip{L_0}A\in C_0$ and $|\ip RE|\le q_R$, adding the second target contribution gives the asserted interval. A union bound controls failure by $\alpha_0+\alpha_1$. The length is the displayed sum. Chernoff or the Gaussian distribution of $\ip RE$ supplies the calibrations. For a prescribed partition, repeat the argument over blocks using the sum of their allocated failures; their centers may be dependent. Selecting among complete constructions from deterministic diameter bounds uses no data and retains its own coverage guarantee. Applying the result with $L_0=0$ also covers every $L$, although it then reduces to the direct interval.
\end{proof}

\subsection{Proof of Theorem~\ref{thm:efficient_regimes}}
\noindent\emph{Roadmap.} Match elementary constructive upper bounds with the existing small-budget lower bound and a product-direction two-point test.
\begin{proof}
In (i), the rate formula has $x\le1$ and is $M/\sqrt p$. The structural radius is $M\sqrt{s_{\rm cap}/p}$, with fixed cap factor. Theorem~\ref{thm:all_norm_rates} gives the matching lower bound.
For (ii), choose the sign of $U$ so that $\ip LU\ge c\norm L_F$ and compare zero with $tU$, $t=\min(M,\tau/2)$. The full likelihood second moment is $e^{t^2/\tau^2}\le e^{1/4}$ and the target difference is at least $ct\norm L_F$. The common-target lemma gives the lower order at zero. Cauchy--Schwarz supplies structural radius $M\norm L_F$, while the direct Gaussian radius is $z_{1-\alpha/2}\tau\norm L_F$. Their deterministic selection controls worst-case length and is polynomial in the number of observed coordinates. The cap class is used only for the lower inclusion; fixed rank caps are unnecessary here.
\end{proof}

\subsection{Proof of Proposition~\ref{prop:convex_loss}}
\noindent\emph{Roadmap.} Find a one-dimensional Gaussian submodel inside every such convex relaxation.
\begin{proof}
The average of $ME_i$, $i=1,\ldots,p$, is $(M/\sqrt p)L$. Both signs and zero belong to the convex hull, hence so does the asserted line segment. Compare zero with $tL$ for $t=\min(M/\sqrt p,\tau/2)$ and apply the same two-point calculation. At $M=\tau\sqrt p$ it forces order $\tau$. The previous minimax rates divided by $\tau$ tend to zero for each fixed $K\ge3$. This argument changes the honesty class and makes no complexity assumption.
\end{proof}

\subsection{Proof of Theorem~\ref{thm:shifted_adaptation}}
\noindent\emph{Roadmap.} Shift the full-data hard mixtures into an unused orthogonal block.
\begin{proof}
Embed the dimension-$m$ rank-one alternatives from Theorem~\ref{thm:all_norm_rates}, with norm budget $M_0$, in coordinates $\{q+1,\ldots,p\}^K$. Denote them by $B$. Their supports are disjoint from $A_0$, so
$\norm{A_0+B}_F^2=\norm{A_0}_F^2+\norm B_F^2\le M^2$.
Unfolding rank subadditivity shows $A_0+B$ is in the outer cap class. Translation of Gaussian likelihoods by $A_0$ leaves their second moments unchanged. All alternatives have a common increment in the normalized $p$-diagonal target equal to $\sqrt{m/p}$ times their dimension-$m$ target. Their total variation from the law at $A_0$ is bounded by a fixed constant less than $1-2\alpha$. Apply the common-target lemma with baseline target $\psi(A_0)$, or translate the interval by that fixed number. This proves the lower bound.
\end{proof}

\subsection{Proof of Theorem~\ref{thm:adaptive_tangent}}
\noindent\emph{Roadmap.} Recover ranks on one spectral event, bound a second-order projection residual, and condition on the independent pilot.
\begin{proof}
For each unfolding use $1/4$-nets of its two unit spheres, of total cardinality at most $9^{p_k+d_k}$. Product telescoping bounds its operator norm by twice the maximum net contraction. Chernoff and a union bound over the $K$ unfoldings give an event $\mathcal E$ of probability at least $1-\delta$ on which every pilot error operator norm is at most $\eta_k$. No independence between unfoldings is used.
The singular-value perturbation inequality
$|\sigma_j(Z^{(1)}_{(k)})-\sigma_j(A_{(k)})|\le\eta_k$
follows from the variational characterization of singular values: for every unit vector the two image norms differ by at most $\eta_k$, and taking the relevant min/max over subspaces preserves this bound. Thus zero singular values produce observed values at most $\eta_k$, while nonzero values exceeding $2\eta_k$ produce observed values greater than $\eta_k$. All selected ranks are exact on $\mathcal E$. If any is zero then $A=0$ and the singleton covers.

For a positive selected rank $r_k$, write $A_{(k)}=U_kS_kV_k^\top$ using its true nonzero singular values. Because $Q_k$ is the leading observed rank-$r_k$ projector,
\[
 \norm{(I-Q_k)A_{(k)}}_{\rm op}
 \le\sigma_{r_k+1}(Z^{(1)}_{(k)})+\eta_k\le2\eta_k.
\]
For full row rank the residual is zero, and the same inequality holds. Right multiplication by $V_kS_k^{-1}$ gives
\[
 \norm{(I-Q_k)U_k}_{\rm op}
 \le\frac{2\eta_k}{\sigma_{r_k}(A_{(k)})}
 \le\frac{2\eta_k}{\sigma_{r_k}(Z^{(1)}_{(k)})-\eta_k}.
\]
Clipping at one yields $d_k^*$. This calculation uses the gap only on the stated separated class, with a strictly positive observed denominator.

Expand the identity into the $2^K$ mutually orthogonal products of $Q_k$ and $I-Q_k$. The projector $T$ retains exactly the terms with at most one complement. Each term discarded by $T$ appears in at least one product $(I-Q_k)(I-Q_\ell)$. Orthogonality therefore implies
\[
 \norm{(I-T)A}_F^2
 \le\sum_{k<\ell}\norm{(I-Q_k)(I-Q_\ell)A}_F^2
 \le M^2\sum_{k<\ell}(d_k^*d_\ell^*)^2.
\]
The last inequality uses $A=(\bigotimes_kU_k)G$ with $\norm G_F=\norm A_F$ and the operator norm of a product of mode maps. Also $\norm{(I-T)A}_F\le M$. Since $T$ is orthogonal, Cauchy--Schwarz bounds
$|\ip L{(I-T)A}|$ by $B_L$.

Condition on the pilot. The vector $TL$ is now deterministic, and independence makes the second error contraction sub-Gaussian with proxy $v_2\norm{TL}_F$. Conditional on any pilot in $\mathcal E$, the projected interval has error at most $\alpha-\delta$. The direct alternative has the same conditional failure bound and the structural alternative never fails. A choice determined by the pilot thus also has conditional failure at most $\alpha-\delta$. Intersecting with a surely valid structural interval preserves coverage; empty intersections may be replaced by a singleton. Integrating and adding $\PP(\mathcal E^c)\le\delta$ proves honesty. SVDs, finite sums, and tensor matrix products give the stated arithmetic complexity.

For Gaussian artificial splitting the two centered errors are jointly Gaussian with zero cross covariance, each of covariance $2\tau^2I$, hence independent. Zero covariance would not suffice for non-Gaussian errors. No minimax length conclusion on the original unseparated class follows from this proof.
\end{proof}


\section{Optimization contracts, loadings and cap-valid inference}
Observe the entire tensor \(Z=A+\tau G\), with iid standard Gaussian
entries, known \(\tau>0\), and \(A\in\Theta(M,\bar r)\), a Tucker or TT
Frobenius ball. The functional is \(\psi_L(A)=\langle L,A\rangle\);
normalization of \(L\) is unnecessary. Write \(J=J(\bar r)\) for the
existing canonical expansion bound: \(A=\sum_{j\le J}a_j U_j\), each
\(U_j\) a unit product tensor, and \(\sum_j a_j^2=\|A\|_F^2\). Zero
terms can be appended. Assume positive caps.

For reference, the previously established equal/flat-loading rate is
\(\mathfrak r_K(m,M,\tau)=\tau\sqrt{\ell_m/m}\,g_K(x)\), where
\(\ell_m=\log(em)\), \(x=M/(\tau\sqrt{\ell_m})\), and

\[g_K(x)=\begin{cases}
\min\{x,1+\max(0,\log x),\ell_m\},&K=3,\\
\min\{x,1\},&K\ge4.
\end{cases}\]

The original equal diagonal has \(m=p\) and weights \(1/\sqrt p\).
Length criteria below mean the infimum expected length at zero
(\(\mathcal L_0\)), or the infimum worst-case expected length
(\(\mathcal L_{\max}\)), over intervals honest on the stated parameter
class.

\subsection{C1. A quantitative contract for approximate minimax
optimization}\label{c1.-a-quantitative-contract-for-approximate-minimax-optimization}\label{sup:C1}

Use every support and every direction in the calibrated finite product
nets. For their unit directions \(U\) and thresholds \(\lambda_{s(U)}\),
put

\[F_Z(B)=\max\{0,\max_U(|\langle U,Z-B\rangle|-\lambda_{s(U)})\}.\]

This function is 1-Lipschitz in Frobenius norm. Suppose an algorithm
returns \(\widehat B\in\Theta\) and numbers
\(l\le\inf_{B\in\Theta}F_Z(B)\le
F_Z(\widehat B)=u\), with \(u-l\le\varepsilon\). On the simultaneous
noise event, the true \(A\) is feasible, so \(F_Z(A)=0\) and
\(u\le\varepsilon\). Consequently, for every support size \(s\),

\[N_S(\widehat B-A)\le 2\lambda_s+\varepsilon.\]

With \(J_2=J(2\bar r)\), the preceding ordered-diagonal proof gives

\[|\widehat B_{i,\ldots,i}-A_{i,\ldots,i}|_{(s)}
\le \min\{2J_2^{(K-1)/2}M,
4J_2^{K/2}(\lambda_s+\varepsilon/2)\}s^{-K/2}.\]

For known orthogonal loading weights \(w_i\), sorted by absolute value,
the sum of the right side against \(w_s^*\) is an honest radius about
\(\psi_L(\widehat B)\). It may be intersected with the structural range.
Choosing \(\varepsilon\le c\tau\) retains the existing all-budget
minimax orders for equal/flat loadings. Indeed every \(\lambda_s\) is
bounded below by a positive constant times \(\tau\), so the added term
changes the existing minimum by only a fixed factor, including in the
norm-limited regime.

For squared error set the scan error budget to \(1/8\). If the optimizer
guarantees \(\varepsilon\le c\tau\) for \emph{every observation}, the
same excess tail proof as before applies:
\(F_Z(\widehat B)\le F_Z(A)+\varepsilon\) and
\(N_S(\widehat B-A)\le2\lambda_s+2F_Z(A)+\varepsilon\). Thus the
existing all-budget MSE bound follows. A budget-limited run with a
larger reported gap does not establish that MSE guarantee.

\textbf{Endpoint version.} Let \(\mathcal F=\{B\in\Theta:F_Z(B)=0\}\).
If certified outer endpoints enclose \(\psi_L(\mathcal F)\), they define
an honest interval, even at an arbitrary computation stop. If feasible
witnesses establish each outer endpoint within \(\varepsilon\) of the
exact endpoint, its length exceeds the ideal projection length by at
most \(2\varepsilon\). An inner optimization solution alone does not
give an honest interval. On a certified empty feasible set use
\(\{0\}\), as before.

\textbf{Proof.} Lipschitzness follows from \(\|U\|_F=1\), the reverse
triangle inequality, and stability of finite maxima. The objective-gap
implication uses \(F_Z\ge0\). Rank caps double under subtraction; the
norm budget doubles as well. The ordered bound is exactly the earlier
net comparison and diagonal quasi-norm inequality with \(2\lambda_s\)
replaced by \(2\lambda_s+\varepsilon\). The Gaussian union event has
probability at least \(1-\alpha\). Endpoint containment proves coverage
on the same event. For the length assertion, use an inner feasible
witness at each end. The squared-risk argument invokes the already
supplied excess-tail integration, with only its deterministic tolerance
constant changed. No algorithmic complexity conclusion follows from this
contract.

\subsection{C2. An implemented global solver on the binary rank-one
class}\label{c2.-an-implemented-global-solver-on-the-binary-rank-one-class}\label{sup:C2}

For \(p=2\) and any fixed order \(K\ge3\), every rank-one tensor of norm
at most \(M\) has a representation

\[B(a,\theta)=a\bigotimes_{k=1}^K(\cos\theta_k,\sin\theta_k),
\quad a\in[-M,M],\quad\theta_k\in[0,\pi].\]

Each vector is represented up to sign on a semicircle; the product of
signs is absorbed into \(a\). The zero tensor is included. There is no
grid restriction on the parameter space.

The code includes both singleton supports and the full support. For the
latter it uses \(n_K=\lceil\pi K\rceil\) equally spaced angles in
\([0,\pi)\) per mode. Every unit vector is within Euclidean distance
\(1/(2K)\) of a net point up to sign: its angular discrepancy modulo
\(\pi\) is at most \(\pi/(2n_K)\). Absolute contractions remove the
signs. The same telescoping proof therefore gives \(\rho\le2N\). There
are \(\mathsf d_1=2\) singleton directions and \(\mathsf d_2=n_K^K\)
full-support directions. Set

\[\lambda_s=\tau\sqrt{2\log(4\mathsf d_s/\alpha)},\qquad s=1,2.\]

The sum of the two Gaussian union bounds is at most \(\alpha\).

For each parameter box, exact real-arithmetic ranges of cosine and
products enclose every scanned contraction and tensor entry. If their
contraction range is \([b_U^-,b_U^+]\), then

\[l_{\rm box}=\max\{0,\max_U[
\operatorname{dist}(\langle U,Z\rangle,[b_U^-,b_U^+])-\lambda_{s(U)}]\}\]

is a lower bound for \(F_Z\) throughout the box. Centers provide actual
rank-one upper witnesses. A largest-scaled-side split covers the parent
with two children. Branch and bound retains global lower and upper
bounds; the objective run ends at its requested gap or reports a budget
stop. It implements the ideal residual criterion itself on this class.

For confidence projection reject a box only when its objective lower
bound is positive. Loading-entry interval arithmetic bounds every
surviving box's target; their hull contains the ideal confidence
projection. Feasible amplitude intervals at fixed central factors supply
inner witnesses. An endpoint gap is reported only when such witnesses
exist. Endpoints need not reach a requested tolerance within a finite
node budget, especially near singular feasibility boundaries. This is
recorded, not treated as a successful exact optimization.

The code supports arbitrary loadings in the binary tensor, but caps
above one and mode dimensions above two are not implemented. Exhaustive
scans and global subdivision are exponential. Tests check range
containment and global certificates against independent finite
calculations. Floating point evaluations include conservative guards;
these are numerical certificates under the derived inequalities, not
formal directed-rounding interval proofs. The real-arithmetic method,
rather than any particular floating point tolerance, is the mathematical
result.

\subsection{L1. Sharp small-budget rates for every loading
tensor}\label{l1.-sharp-small-budget-rates-for-every-loading-tensor}\label{sup:L1}

Define the product overlap

\[\rho(L)=\sup_{\|u_k\|_2=1}|\langle L,\bigotimes_k u_k\rangle|.\]

For fixed positive caps, fixed \(0<\alpha<1/4\), arbitrary mode
dimensions, and \emph{every} known loading tensor \(L\), if
\(0<M\le\tau\), then

\[\mathcal L_0(\Theta;L)\asymp\mathcal L_{\max}(\Theta;L)
\asymp M\rho(L),\qquad
\inf_{\widehat\psi}\sup_\Theta\mathbb E_A
(\widehat\psi-\psi_L(A))^2\asymp M^2\rho(L)^2.\]

Constants depend only on the fixed caps/order and confidence, not on
\(L\) or the ambient dimensions. The zero loading has identically zero
risk and length. More generally, for all \(M,\tau>0\),

\[c_\alpha\rho(L)\min(M,\tau)
\le \mathcal L_0\le\mathcal L_{\max}
\le 2\min\{\sqrt J M\rho(L),\ z_{1-\alpha/2}\tau\|L\|_F\}.\]

\textbf{Proof.} The canonical expansion and Cauchy--Schwarz give
\(|\psi_L(A)|\le\rho(L)\sum_j|a_j|\le\sqrt J M\rho(L)\). This gives a
deterministic interval and the risk of estimator zero. The direct
Gaussian interval gives the other upper bound; choose between the two
using known quantities. Compactness gives a maximizing product \(U\);
choose its sign so \(\langle L,U\rangle=\rho(L)\). Both \(0\) and \(tU\)
belong to every positive cap class, for \(t=\min(M,\tau)/2\). Their
Gaussian total variation is \(2\Phi(t/(2\tau))-1\le2\Phi(1/4)-1<1/4\).
Honesty at both parameters implies
\(\mathbb E_0|C|\ge t\rho(L)(1-2\alpha-\mathrm{TV})\). The two-point
squared-loss lower bound is a fixed positive multiple of
\(t^2\rho(L)^2\): classify an estimate by its nearest target. A
classification error entails squared error at least \(t^2\rho(L)^2/4\),
while the sum of the two testing errors is at least \(1-\mathrm{TV}\).
Therefore the larger squared risk is at least
\(t^2\rho(L)^2(1-\mathrm{TV})/8\). For \(M\le\tau\) the bounds match.
This proof is a direct specialization of standard Gaussian two-point
reasoning, not a novel general inference principle.

\textbf{Limits.} Computing \(\rho(L)\) for a general high-order loading
is itself a global optimization problem. A certified upper bound can
replace it in the interval; a certified product witness supplies a lower
bound. Without a constant ratio between them, an efficient sharp
algorithm has not been obtained. This theorem covers all loading tensors
in the stated budget regime; it does not characterize their
larger-budget rates. The flat loading small-budget theorem extends to
\(M\lesssim\tau\sqrt{\log m}\), a larger regime than the
arbitrary-loading result here.

\subsection{L2. An optimized common-target lower bound for unequal
weights}\label{l2.-an-optimized-common-target-lower-bound-for-unequal-weights}\label{sup:L2}

For an orthogonal product loading with decreasing nonzero absolute
weights \(w_1\ge\cdots\ge w_m>0\), choose \(d\le m\) and
\(0<\Delta\le Mw_d\). For signed original weights \(\widetilde w_i\),
set \(A_i=(\Delta/\widetilde w_i)U_i\), \(1\le i\le d\). All
alternatives have target \(\Delta\), satisfy the norm budget, and have
rank one. They are orthogonal as full tensors. Let

\[v_i=\exp\{\Delta^2/(\tau^2w_i^2)\}-1,\quad
H_d(\Delta)=\sum_{i=1}^d v_i^{-1},\quad
\pi_i=v_i^{-1}/H_d(\Delta).\]

The optimal chi-square divergence among mixtures on these alternatives
is

\[\chi^2(Q,P_0)=\min_{\pi\in\mathcal S_d}\sum_i\pi_i^2v_i
=H_d(\Delta)^{-1}.\]

Thus every interval honest on the full tensor class satisfies

\[\mathbb E_0|C|\ge\Delta\left[1-2\alpha-
\frac{1}{2\sqrt{H_d(\Delta)}}\right]_+.\]

\textbf{Proof.} The Gaussian likelihood identity gives overlap 1 for
distinct orthogonal alternatives and \(1+v_i\) for equal alternatives.
Subtracting 1 leaves \(\sum\pi_i^2v_i\). Cauchy--Schwarz bounds this
below by \((\sum v_i^{-1})^{-1}\) and the displayed weights attain
equality. Apply \(\mathrm{TV}\le\sqrt{\chi^2}/2\) and the common-target
interval lemma. This uses the full observation; there is no
diagonal-data restriction.

\texttt{weighted\_axis\_lower} enumerates \(d\) and uses scalar root
finding for the largest \(\Delta\le Mw_d\) with \(\chi^2\le1/4\). Every
reported lower bound uses the actual recomputed divergence and a
feasible prior. This optimizes the prior at each chosen target within
the axis family. It is not an optimization over all tensor alternatives.
For equal weights it recovers the one-spike logarithmic obstruction. It
cannot recover the cubic harmonic obstruction on its own and is not
claimed sharp for all weights.

\subsection{A1. Unseparated rate adaptation across positive fixed
caps}\label{a1.-unseparated-rate-adaptation-across-positive-fixed-caps}\label{sup:A1}

Fix an outer positive Tucker or TT cap vector \(\bar r\). Let
\(\mathcal R\) contain any collection of positive cap vectors
componentwise below it. For equal or flat orthogonal loadings, the
\emph{same} full-data outer-class interval satisfies

\[\sup_{A\in\Theta(M,r)}\mathbb E_A|C_{\bar r}|
\le C_{K,\bar r,\alpha}\,\mathfrak r_K(m,M,\tau),\qquad r\in\mathcal R.\]

Every smaller class contains the rank-one ball, whose lower bound is
\(c_{K,\alpha}\mathfrak r_K\). Therefore the interval is simultaneously
minimax in rate over every such class, uniformly in \(m,M,\tau\),
without knowing the actual ranks or imposing singular-value separation.
The analogous MSE conclusion holds for the outer-class ideal residual
fit.

\textbf{Proof.} Restrict the existing outer supremum upper bound to each
subset and apply the same rank-one lower bound inside each smaller
class. No selection step or new probability argument is required.

This observation resolves the fixed-positive-cap, rate-only statistical
adaptation question. It does not give optimal rank-dependent constants,
growing-cap rates, computation, exact-rank selection, or pointwise
oracle lengths. Rank zero must be excluded: its known-model oracle
interval has length zero, whereas outer-class honesty forces positive
length at zero. The preceding ledgers used the word adaptation broadly;
these distinct objectives must be separated.

\subsection{A2. A computable projection certificate without
separation}\label{a2.-a-computable-projection-certificate-without-separation}\label{sup:A2}

Let \(Z^{(1)}=A+E^{(1)}\) and \(Z^{(2)}=A+E^{(2)}\) have independent iid
Gaussian errors of known scales \(\tau_1,\tau_2\). Only \(\|A\|_F\le M\)
is assumed here; supplied rank caps are unnecessary. Let
\(0<\delta<\alpha\). Use the simultaneous unfolding operator envelopes
from the previous stage:

\[\eta_k=2\tau_1\sqrt{2\{(p_k+\prod_{l\ne k}p_l)\log9+
\log(2K/\delta)\}}.\]

On one event of probability at least \(1-\delta\) all
\(\|E^{(1)}_{(k)}\|_{op}\le\eta_k\). For any pilot-chosen ranks \(r_k\),
let \(Q_k\) project onto the first \(r_k\) left singular vectors of
\(Z^{(1)}_{(k)}\). Set \(d_k=\sigma_{r_k+1}(Z^{(1)}_{(k)})+\eta_k\) if
\(r_k<p_k\), and \(d_k=0\) if \(Q_k=I\). Singular values beyond the
matrix rank are zero. For \(P=\bigotimes_k Q_k\), decompose

\[R=(I-P)L=\sum_{k=1}^K R_k,\quad
R_k=\left[\bigotimes_{l<k}Q_l\right]\otimes(I-Q_k)
\otimes\left[\bigotimes_{l>k}I\right]L.\]

Then the observable bound

\[B=\min\left\{M\|R\|_F,\ \sum_k d_k\|(R_k)_{(k)}\|_*\right\}\]

satisfies \(|\langle R,A\rangle|\le B\) on that event. The interval
centered at \(\langle PL,Z^{(2)}\rangle\), with radius

\[B+z_{1-(\alpha-\delta)/2}\tau_2\|PL\|_F,\]

is honest over the whole norm ball. Any finite set of pilot-selected
projection choices can be compared by their radii on this same event;
choose the smallest before using validation data. Include \(P=I\) for
the direct interval and \(P=0\) for the structural interval. Intersect
with \([-M\|L\|_F,M\|L\|_F]\) and use \(\{0\}\) for an empty
intersection.

\textbf{Proof.} Residual SVD optimality and the triangle inequality give
\(\|(I-Q_k)A_{(k)}\|_{op}\le d_k\) on the common event, even for
selected \(r_k\). The definition of \(R_k\) gives
\((I-Q_k)(R_k)_{(k)}=(R_k)_{(k)}\). Operator/nuclear norm duality bounds
its contraction with \(A\) by \(d_k\|(R_k)_{(k)}\|_*\). Sum and combine
with the norm-only bound. Conditional on the pilot, validation noise in
the selected loading is Gaussian with variance \(\tau_2^2\|PL\|_F^2\).
No multiplicity correction over the candidate ranks is required: all
bias certificates hold on one uniform pilot event, and the choice
precedes validation. The union bound gives total error at most
\(\alpha\). Unfolding SVDs, tensor projections and nuclear norms are
polynomial in the explicit input size.

This algorithm never interprets an undetected component as identically
zero. It does not certify the actual rank. Its bias term can be large,
and no weak-signal or pointwise minimax optimality is claimed. Honesty
on the full norm ball cannot yield the low-rank weak-signal minimax
gain: the earlier unrestricted-class lower bound still applies.
Independent artificial Gaussian splitting of one observation is valid
with \(\tau_1=\tau_2=\sqrt2\tau\); using both halves at scale \(\tau\)
is invalid.

\subsubsection{Why the first-order repair does not solve the original
problem}\label{why-the-first-order-repair-does-not-solve-the-original-problem}

For the A2 construction, if \(\min_k\eta_k\ge z_{\mathrm{val}}\tau_2\),
where \(z_{\mathrm{val}}=z_{1-(\alpha-\delta)/2}\), no projected
candidate improves the smaller of the untruncated direct and norm-only
radii. In fact, nonzero \(R_k\) have \(d_k\ge\eta_k\), so triangle
inequalities and nuclear/Frobenius norm comparison give
\(B\ge\min(M,\min_k\eta_k)\|R\|_F\). Hence its radius is at least
\(\min(M,z_{\mathrm{val}}\tau_2)(\|PL\|_F+\|R\|_F)\), which is at least
\(\min(M,z_{\mathrm{val}}\tau_2)\|L\|_F\). The explicit envelopes meet
this condition in the equal-noise experiments below. This is a proved
limitation of this particular repair, not a universal computational
lower bound.

\subsection{A3. The adaptation obstruction at an arbitrary
orientation}\label{a3.-the-adaptation-obstruction-at-an-arbitrary-orientation}\label{sup:A3}

Let a base tensor \(A_0\) have norm \(a<M\), and leave at least one
unused unit of every supplied Tucker or TT rank cap. Choose mode
subspaces \(W_k\) such that \(A_0\) is orthogonal to every tensor in
\(\mathcal W=\bigotimes_k W_k\). One sufficient choice is to take at
least one \(W_k\) orthogonal to the column space of \(A_{0,(k)}\). Put
\(M_0=\sqrt{M^2-a^2}\) and \(L_{\mathcal W}=P_{\mathcal W}L\). Then

\[\mathbb E_{A_0}|C|\ge c_\alpha\min(M_0,\tau)\rho(L_{\mathcal W})\]

for every interval honest on the outer class. If the restriction of the
loading to an \(m\)-dimensional orthogonal diagonal block in
\(\mathcal W\) equals \(b\) times a flat unit loading, the stronger
bound is

\[\mathbb E_{A_0}|C|\ge c_{K,\alpha}|b|\,
\mathfrak r_K(m,M_0,\tau).\]

Here \(\rho\) is taken within the indicated product subspace. The zero
overlap case gives a zero lower bound and says nothing further.

\textbf{Proof.} Add rank-one perturbations \(B\in\mathcal W\) with norm
at most \(M_0\). Rank subadditivity keeps \(A_0+B\) in every outer
unfolding cap; orthogonality gives
\(\|A_0+B\|_F^2=a^2+\|B\|_F^2\le M^2\). Translate the Gaussian
experiment by \(A_0\). The target difference is
\(\langle L_{\mathcal W},B\rangle\). Apply the two-point construction of
L1 or the existing flat-loading mixtures, all of which have rank one.
Unused observed coordinates contain identical independent Gaussian noise
under both hypotheses and do not change the likelihood calculation. This
is a rotation-invariant extension of the preceding block-supported
shifted bound, not a complete characterization at every mean.

\subsection{A4. A cap-valid tangent interval on the original unseparated
class}\label{a4.-a-cap-valid-tangent-interval-on-the-original-unseparated-class}\label{sup:A4}

There is a more useful repair when mode caps \(\bar r_k\le p_k\) are
supplied. Use the same independent Gaussian pilot and validation data
and envelopes as A2. Let \(Q_k\) project onto the leading \(\bar r_k\)
pilot left singular vectors. Define

\[e_k=\begin{cases}
0,&\bar r_k=p_k,\\
\min\{1,2\eta_k/(\sigma_{\bar r_k}(Z^{(1)}_{(k)})-\eta_k)\},
&\bar r_k<p_k,\ \sigma_{\bar r_k}(Z^{(1)}_{(k)})>\eta_k,\\
1,&\text{otherwise}.
\end{cases}\]

With the first-order product-space projector from the previous stage,

\[T=\bigotimes_k Q_k+\sum_k(I-Q_k)\otimes\bigotimes_{l\ne k}Q_l,\]

put \(R=(I-T)L\) and

\[B_T=M\|R\|_F\min\{1,(\sum_{k<l}e_k^2e_l^2)^{1/2}\}.\]

The interval centered at \(\langle TL,Z^{(2)}\rangle\), with radius
\(B_T+z_{\mathrm{val}}\tau_2\|TL\|_F\), is honest on \textbf{every}
tensor with these supplied mode caps and norm bound, including zero,
arbitrarily weak signals, and actual ranks strictly below the caps. No
singular-value separation is assumed. Choose the smallest
pilot-measurable radius among this candidate, the direct interval, and
the structural norm interval, as before.

\textbf{Proof.} Let \(V_k\) be the actual left singular subspace of
\(A_{(k)}\). On the simultaneous noise event we claim
\(\|(I-Q_k)P_{V_k}\|_{op}\le e_k\). This is immediate for \(e_k=1\) or
\(Q_k=I\). Otherwise \(\sigma_{\bar r_k}(Z^{(1)}_{(k)})>\eta_k\) forces
the true mode rank to equal the supplied cap: a smaller true rank would
give that observed singular value at most \(\eta_k\) by Weyl's
inequality. The pilot tail after the cap is at most \(\eta_k\), so
\(\|(I-Q_k)A_{(k)}\|_{op}\le2\eta_k\). Multiplying by the inverse
nonzero singular values and using \(\sigma_{\bar r_k}(A_{(k)})\ge
\sigma_{\bar r_k}(Z^{(1)}_{(k)})-\eta_k\) proves the claim.

Decompose \(A\) into the mutually orthogonal mode-projector blocks
indexed by which modes fall in \(I-Q_k\). The complement of \(T\)
consists of blocks with at least two such modes. Counting each block
once is bounded by counting it for every contained pair. Consequently

\[\|(I-T)A\|_F^2\le\sum_{k<l}
\|(I-Q_k)(I-Q_l)A\|_F^2
\le M^2\sum_{k<l}e_k^2e_l^2.\]

The last inequality restricts each operator to its true mode subspace
and uses that operators on distinct modes commute. Combine with the
trivial bound \(\|(I-T)A\|_F\le M\) and Cauchy--Schwarz against \(R\).
Conditional Gaussian validation and the common pilot event now prove
coverage and valid radius selection exactly as in A2.

This algorithm is polynomial and honest for the original weak-signal
class. Strong signals can make the second-order bias correction small,
even though separation is not a coverage assumption. It is \textbf{not}
a polynomial weak-signal minimax solution or a method of estimating
actual rank. At smaller actual ranks, the unresolved cap singular values
trigger conservative error bounds. For a supplied TT cap vector, valid
mode caps are \(\min(p_k,\bar r_{k-1}\bar r_k)\), with boundary TT ranks
one; the interval does not estimate TT bond ranks. Artificial splitting
must use its increased noise scales.

For the original single observation,
\texttt{single\_observation\_interval} creates \(Z\pm\tau G'\) with
independent artificial \(G'\), uses scales \(\sqrt2\tau\) and error
\(\alpha/2\) for the cap-valid construction, and intersects with the
direct interval based on \(Z\) at error \(\alpha/2\). The union bound
proves honesty at error \(\alpha\) despite the dependence between these
two intervals. The direct component caps the reported length at
\(2z_{1-\alpha/4}\tau\|L\|_F\). The split is randomization within the
original experiment, not an extra observed replicate.

\subsection{Remaining problems}\label{remaining-problems}

\begin{enumerate}
\def\labelenumi{\arabic{enumi}.}
\tightlist
\item
  Polynomial minimax intervals for diffuse weak-signal classes at
  general dimensions: open. The new optimizer establishes a working
  global benchmark only for binary rank-one tensors; it does not remove
  the computational obstruction.
\item
  Sharp larger-budget rates for every loading, including every unequal
  weight sequence: open. L1 resolves all loadings only for \(M\le\tau\);
  L2 strengthens the usable lower-bound machinery.
\item
  General rank-dependent and pointwise adaptation: open, with explicit
  impossibility results for some oracle demands. Rate adaptation over
  fixed positive nested caps is already attained statistically (A1).
\item
  General Tucker/TT global optimizers, validated directed-rounding
  arithmetic, and high-dimensional runtime scaling: not implemented.
\end{enumerate}


\section{Growing caps, decaying loadings and pointwise inference}
Observe \(Z=A+\tau G\) with independent standard Gaussian tensor
entries, known \(\tau>0\), norm bound \(\|A\|_F\le M\), and supplied
Tucker or TT caps. Unless specified otherwise, dimensions are equal to
\(p\) and \(K\ge3\). For \(L=p^{-1/2}\sum_{i=1}^p E_i\), write
\(\psi(A)=\langle L,A\rangle\), where \(E_i=e_i^{\otimes K}\). The
criteria \(\mathcal L_0\) and \(\mathcal L_{\max}\) are optimal expected
length at zero and optimal worst-case expected length among intervals
honest over the whole stated class. Fix \(0<\alpha<1/4\).

\subsection{G1. Sharp growing-cap rates with a polynomial
algorithm}\label{g1.-sharp-growing-cap-rates-with-a-polynomial-algorithm}\label{sup:G1}

Let \(r=\min(p,\min_k\bar r_k)\), where the minimum is over the supplied
constrained unfoldings (Tucker modes or TT cuts). Suppose
\(r\ge c\sqrt p\) for a fixed \(c>0\). Uniformly over all remaining
caps, all \(p\), and every \(M,\tau>0\),

\[\mathcal L_0\asymp_{\alpha,c}\mathcal L_{\max}
\asymp_{\alpha,c}\min\{M\sqrt{r/p},\tau\}.\]

Minimax squared risk has the square of this rate. The constants do not
depend on growing ranks or tensor order. Choose deterministically
between the structural interval \([-h_0,h_0]\), \(h_0=M\sqrt{r/p}\), and
the direct Gaussian interval centered at \(\psi(Z)\) with radius
\(z_{1-\alpha/2}\tau\). This takes \(O(p)\) accessed entries and
arithmetic for the coordinate loading. For a rotated flat loading,
obtaining the contractions has its additional polynomial input-access
cost. No factor estimation or exponential scan is required.

\textbf{Upper proof.} An unfolding attaining the smallest cap has
loading matrix singular values \(1/\sqrt p\), repeated \(p\) times.
Matrix singular-value duality and Cauchy--Schwarz bound its contraction
with any rank-at-most-\(r\) matrix of norm \(M\) by \(M\sqrt{r/p}\).
Thus the structural interval covers surely. The direct error is
\(N(0,\tau^2)\). Deterministic radius selection proves the length upper
bound. Choosing estimator zero when \(h_0\le\tau\) and the direct
estimator otherwise gives squared risk at most \(\min(h_0^2,\tau^2)\).

\textbf{Lower proof.} For a uniform size-\(r\) subset \(S\) of
\(\{1,\ldots,p\}\) set

\[A_S=b\sum_{i\in S}E_i,\qquad
b=\min\{M/\sqrt r,\tau\sqrt{\log(1+p/(4r^2))}\}.\]

All constrained unfolding ranks are \(r\), the norm is at most \(M\),
and the common target is \(\Delta=rb/\sqrt p\). For independent supports
\(S,T\), the full Gaussian likelihood second moment equals
\(\mathbb E\exp\{b^2|S\cap T|/\tau^2\}\). To bound it, for
\(H=|S\cap T|\) factorial moments give
\(\mathbb E\binom Hj=\binom rj(r)_j/(p)_j\le\binom rj(r/p)^j\). Expand
\((1+t)^H\) at \(t=e^{b^2/\tau^2}-1\) to obtain

\[1+\chi^2(Q,P_0)\le(1+rt/p)^r\le\exp(r^2t/p)\le e^{1/4}.\]

Put \(x=p/(4r^2)\le1/(4c^2)\). Since \(x/(1+x)\le\log(1+x)\le x\),

\[\frac{\tau}{2\sqrt{1+1/(4c^2)}}
\le\frac{\tau r}{\sqrt p}\sqrt{\log(1+p/(4r^2))}\le\frac\tau2.\]

Consequently \(\Delta\asymp_c\min(M\sqrt{r/p},\tau)\). Total variation
is at most \(\eta=\sqrt{e^{1/4}-1}/2<1/2\). Coverage of both zero and
the common target implies \(\mathbb E_0|C|\ge\Delta(1-2\alpha-\eta)\),
which is positive for \(\alpha<1/4\). The nearest-target testing
argument also gives squared risk at least \(\Delta^2(1-\eta)/8\). This
proves all claims.

This closes the large-cap regime using the earlier sparse-diagonal lower
bound. It is a consequence of that bound with a newly matched upper
bound, not a new sparse-mixture technique. Under \(M=\tau\sqrt p\) the
rate is \(\tau\), so the vanishing fixed-cap improvement has already
disappeared by rank order \(\sqrt p\), well below unrestricted rank
\(p\).

\subsection{G2. General adaptation across positive growing caps is
impossible}\label{g2.-general-adaptation-across-positive-growing-caps-is-impossible}\label{sup:G2}

Set \(M=\tau\sqrt p\) and let the outer smallest cap satisfy
\(r\ge c\sqrt p\). Every interval honest on that outer class has
\(\mathbb E_0|C|\ge c_{\alpha,c}\tau\). For a fixed positive inner cap
vector, the previously established oracle minimax rate is

\[\mathfrak r_K(p)=\frac\tau{\sqrt p}
\begin{cases}(\log(ep))^{3/2},&K=3,\\
\sqrt{\log(ep)},&K\ge4.
\end{cases}\]

Zero belongs to every such inner class. Its worst-case expected length
under outer honesty is therefore at least order \(\tau\), a diverging
factor \(\sqrt p/(\log(ep))^{3/2}\) or \(\sqrt{p/\log(ep)}\) larger than
its known-class oracle rate. No algorithm, including an exponential one,
can be simultaneously minimax in rate over these positive fixed and
growing cap classes. This is stronger than excluding only the
known-rank-zero oracle. It does not conflict with earlier adaptation
across \emph{fixed} outer caps, whose constants were allowed to depend
on that outer cap.

\textbf{Nonzero version.} Suppose \(A_0\) is supported on the first
\(q\) coordinates in every mode, with norm \(a<M\) and actual
constrained ranks \(r_{0,k}\). Let \(m=p-q\),
\(s=\min(m,\min_k(\bar r_k-r_{0,k}))\) and \(M_0=\sqrt{M^2-a^2}\). If
\(s\ge c\sqrt m\), then

\[\mathbb E_{A_0}|C|\gtrsim_{\alpha,c}
\min\{M_0\sqrt{s/p},\tau\sqrt{m/p}\}.\]

\textbf{Proof.} Add the G1 alternatives on the disjoint remaining block.
Norm squares add; every unfolding rank increases by at most \(s\).
Translation by \(A_0\) preserves Gaussian likelihood overlaps, and
normalization of the target multiplies the dimension-\(m\) lower bound
by \(\sqrt{m/p}\). This gives a positive nonzero-base obstruction when
weak extra rank fits inside the supplied caps. Arbitrarily oriented
versions follow by using orthogonal product blocks when the transformed
loading is flat on that block.

\subsection{W1. A dyadic loading bound using the full weight
profile}\label{w1.-a-dyadic-loading-bound-using-the-full-weight-profile}\label{sup:W1}

Return to a prescribed orthogonal product loading
\(L_w=\sum_{i=1}^m w_i U_i\), with decreasing absolute weights and fixed
positive caps. The ambient mode dimensions may exceed \(m\). Partition
the ordered indices into \(B_0=\{1\}\) and
\(B_j=\{2^{j-1}+1,\ldots,\min(2^j,m)\}\) for \(j\ge1\), stopping when
empty. Let \(a_j=\max_{i\in B_j}|w_i|\), \(d_j=|B_j|\), and
\(\alpha_j=\alpha/2^{j+1}\). Define

\[H_K(w)=\sum_j a_j(j+1)^{\nu_K},\qquad
\nu_3=3/2,\quad \nu_K=1/2\ (K\ge4).\]

There is a full-data honest interval satisfying the deterministic bound

\[|C_w|\le C_{K,\bar r,\alpha}
\min\{M\|w\|_\infty,\tau H_K(w)\}.\]

The minimax squared risk is at most a constant times the square of the
same expression. This is a statistical upper bound using the ideal block
residual fits, not a polynomial implementation of those fits.

\textbf{Proof.} Rotate the known orthogonal factors into coordinates.
Restrict the data to each \(d_j^K\) block; this cannot increase its norm
or rank caps. Apply the earlier weighted residual confidence projection
at error \(\alpha_j\). With \(J_2=J(2\bar r)\), its block diameter is at
most

\[4J_2^{K/2}a_j\sum_{s=1}^{d_j}\lambda_s^{(j)}s^{-K/2},\]

where

\[\lambda_s^{(j)}=\tau\sqrt{2\{s\log(ed_j/s)+Ks\log(1+4K)
+\log(2d_j/\alpha_j)\}}.\]

For \(s\ge1\), \(\log(ed_j)\) and \(\log(1/\alpha_j)\) are bounded by a
fixed constant times \(j+1\). Thus
\(\lambda_s^{(j)}\le C_{K,\alpha}\tau\sqrt{s(j+1)}\). The remaining sum
is harmonic for \(K=3\), bounded by \(C(j+1)\), and bounded by a
\(K\)-dependent constant for \(K\ge4\). Add the block intervals and use
\(\sum_j\alpha_j\le\alpha\). Their lengths sum to at most
\(C\tau H_K(w)\). The old structural radius is at most
\(M\sqrt{\min\bar r_k}\|w\|_\infty\), which is \(O(M\|w\|_\infty)\) for
fixed caps. Select the smaller deterministic bound before seeing the
data.

For squared error use block residual fits with error budgets
\(\beta_j=2^{-j-4}\) and deterministic optimization tolerance \(\tau\).
The earlier excess-tail proof gives each block error an \(L_2\) norm at
most \(C\tau a_j(j+1)^{\nu_K}\): the deterministic diameter term has the
bound just derived, and the excess term has \(L_2\) norm \(O(\tau a_j)\)
because \(\sum_s s^{-K/2}<\infty\) and
\(\mathbb E T_{0,j}^2\le2\beta_j\tau^2\). Minkowski's inequality bounds
the \(L_2\) norm of their summed error by \(C\tau H_K(w)\), without
requiring independence. Compare with estimator zero and its structural
risk bound. The finite number of blocks makes measurability and
summation immediate.

\subsection{W2. Sharp rates at every budget for polynomially decreasing
loadings}\label{w2.-sharp-rates-at-every-budget-for-polynomially-decreasing-loadings}\label{sup:W2}

Suppose for fixed \(\beta>0\) and \(C_0<\infty\) that the nonzero,
decreasing weights satisfy

\[|w_i|\le C_0|w_1|i^{-\beta},\qquad i=1,\ldots,m.\]

For fixed positive caps and order, uniformly over all \(m,M,\tau>0\),

\[\mathcal L_0(\Theta;L_w)\asymp\mathcal L_{\max}(\Theta;L_w)
\asymp |w_1|\min(M,\tau),\]

and minimax squared risk is \(\asymp |w_1|^2\min(M,\tau)^2\). Constants
may depend on \(\beta,C_0,K,\bar r,\alpha\), but not on dimension or
budget. This includes all signed power-law weights, whether normalized
or not. In particular, \(0<\beta<1/2\) permits diffuse unit-Frobenius
loadings with \(|w_1|\to0\); the result is not limited to a constant
product overlap with a unit-Frobenius loading.

\textbf{Proof.} For \(j\ge1\), \(a_j\le C_0|w_1|2^{-\beta(j-1)}\). The
series \(\sum_j2^{-\beta(j-1)}(j+1)^{\nu_K}\) is finite for every fixed
\(\beta>0\). W1 therefore gives the matching length and risk upper
bounds. For the lower bounds compare zero with \(tU_1\), where
\(t=\min(M,\tau)/2\) and its sign agrees with \(w_1\). These rank-one
means belong to every positive cap class, their target difference is
\(t|w_1|\), and their total variation is at most \(2\Phi(1/4)-1<1/4\).
The common-target interval argument and two-point squared-loss argument
give the asserted lower orders. No lower power-envelope assumption on
the other weights is needed.

As \(\beta\downarrow0\) the constants deteriorate: the dyadic sum is of
order at most \(1+\beta^{-(\nu_K+1)}\) up to fixed constants. Thus this
theorem does not cover a dimension-dependent exponent approaching zero
uniformly, and does not contradict the extra logarithmic factors for
flat loadings. Completely arbitrary unequal weight sequences, including
many slowly varying profiles, remain uncharacterized.

\textbf{Nonorthogonal perturbations.} Suppose \(L=L_w+R\), with the same
weight condition and a known otherwise arbitrary tensor \(R\) satisfying
\(\|R\|_F\le\epsilon|w_1|\) for a fixed \(0\le\epsilon<1\). The same
sharp rates hold with constants depending also on \(1-\epsilon\). For
the lower bound, \(|\langle L,U_1\rangle|\ge(1-\epsilon)|w_1|\). For
large budgets add a direct interval for \(\langle R,A\rangle\) at a
separate error allocation, costing \(O(\tau\|R\|_F)\). For small
budgets, \(\rho(L_w)=|w_1|\) by generalized Hölder and
\(\rho(R)\le\|R\|_F\); the canonical expansion bounds the structural
radius by \(\sqrt J M(1+\epsilon)|w_1|\). The same reasoning applies to
squared error, using the direct residual estimator and Minkowski. This
gives an all-budget result for a class of loadings without an exact
orthogonal representation, but is not a theorem for every loading.

\subsection{W3. A restricted computational separation for these newly
matched
loadings}\label{w3.-a-restricted-computational-separation-for-these-newly-matched-loadings}\label{sup:W3}

The W2 rate does not automatically have an efficient estimator. To make
this distinction precise, let \(w_i=w_1i^{-\beta}>0\) with fixed
\(0<\beta<1/K\), and set \(M=\tau\sqrt p\). Choose constants

\[\frac1K<a<\min\left\{\frac12,\frac{1-2\beta}{K-2}\right\},
\qquad 0<b<a-\frac1K.\]

For \(s=\lfloor p^a\rfloor\) and \(\mathsf d=\lfloor p^b\rfloor\), every
polynomial estimator of total degree at most \(\mathsf d\) in the full
observed tensor has worst-case squared risk polynomially larger than the
unrestricted \(\tau^2w_1^2\) rate of W2.

\textbf{Finite statement and proof.} For a uniform size-\(s\) support
\(S\), put \(x_S=s^{-1/2}\mathbf1_S\) and let \(\mathsf d_w\) be the
diagonal map with entries \(w_p/w_i\le1\). Use the rank-one alternatives

\[A_S=\tau\lambda(\mathsf d_wx_S)\otimes x_S^{\otimes(K-1)},\qquad
0<\tau\lambda\le M.\]

They have norm at most \(M\) and common weighted target
\(\Delta=\tau\lambda w_p s^{1-K/2}\). All factors are nonnegative, so
for two supports \(0\le\langle A_S,A_T\rangle/\tau^2\le
\lambda^2|S\cap T|^K/s^K\). The full Gaussian Hermite expansion
therefore bounds the squared norm of the centered likelihood ratio
through degree \(\mathsf d\) by

\[\chi^2_{\le \mathsf d}\le\sum_{d=1}^{\mathsf d}
\frac{\lambda^{2d}\mathbb E H^{Kd}}{d!s^{Kd}}
\le\frac{\mu e^\mu B_{\mathsf d}}{1-B_{\mathsf d}}=\eta_{\mathsf d}^2,
\quad \mu=s^2/p,\quad B_{\mathsf d}=\frac{\lambda^2(K\mathsf d)^K}{s^K}<1.\]

For completeness, counting maps by their nonempty fibers gives
\(H^q=\sum_{j=1}^q S(q,j)(H)_j\), where \(S(q,j)\) counts partitions of
\(q\) elements into \(j\) nonempty sets. We have \(S(q,j)\le j^q/j!\)
and \(\mathbb E(H)_j\le\mu^j\). Hence
\(\mathbb E H^q\le q^q(e^\mu-1)\le\mu e^\mu q^q\). Substitute \(q=Kd\)
and sum the geometric upper bound to get the display. For a
degree-\(\mathsf d\) estimator with worst-case squared risk \(R\),
orthogonality of Hermite polynomials and Cauchy--Schwarz give
\(|\mathbb E_QT-\mathbb E_0T|\le\sqrt R\,\eta_{\mathsf d}\). The two
target means differ by \(\Delta\) and their estimation biases are each
at most \(\sqrt R\), so

\[R\ge\Delta^2/(2+\eta_{\mathsf d})^2.\]

Set \(\lambda=\sqrt p\). Then \(\mu\to0\) and
\(B_{\mathsf d}=O(p^{1-Ka+Kb})\to0\), whereas

\[\frac{\Delta}{\tau w_1}
\asymp p^{1/2-\beta-a(K-2)/2}\longrightarrow\infty\]

at a polynomial rate by the choice of \(a\). Such an \(a\) exists
precisely under the stated sufficient condition \(\beta<1/K\). This
proves the restricted separation, using the earlier general low-degree
argument with a new loading-compatible prior. It includes diffuse
power-law loadings with sharp all-budget rates from W2.

This is not a lower bound for all polynomial-time algorithms or all
confidence intervals. An interval honest on zero and these alternatives
would give the test \(\mathbf1\{\Delta\in C(Z)\}\) with total error at
most \(2\alpha+\mathbb E_0|C|/\Delta\). Excluding such efficient tests
requires an additional computational hypothesis; the
polynomial-estimator result alone does not supply one.

\subsection{P1. Efficient pointwise optimality at strong saturated
Tucker
means}\label{p1.-efficient-pointwise-optimality-at-strong-saturated-tucker-means}\label{sup:P1}

This result concerns \textbf{Tucker mode caps all equal to \(r\)},
\(1\le r\le p\). It does not substitute TT bond caps for mode caps. Let

\[A_0=\sum_{i\in S}a_i E_i,\qquad |S|=r,\quad
a=\min_{i\in S}|a_i|>0,\quad
\|A_0\|_F+\tau\le M\le c_0a\sqrt r,\]

where \(c_0>1\) is fixed. The support and signed amplitudes are not
supplied to the algorithm. These means saturate every supplied mode cap.
There is a polynomial interval honest on the entire unseparated Tucker
class such that, in the strong-signal regime specified below,

\[\mathbb E_{A_0}|C|\asymp_{K,\alpha,c_0}\tau\sqrt{r/p},\]

and this matches the infimum expected length at \(A_0\) under that
\emph{same outer-class honesty requirement}. Both \(p\) and \(r\) may
grow. This is pointwise optimality in rate, not optimal leading
constants.

\textbf{Algorithm and noise budgets.} Make the valid artificial split
\(Z^{(1)}=Z+\tau G'\), \(Z^{(2)}=Z-\tau G'\), with independent errors of
scale \(v=\sqrt2\tau\). Set \(\delta=\alpha r/(8p)\) and

\[\eta=2v\sqrt{2\{(p+p^{K-1})\log9+\log(2K/\delta)\}}.\]

Use the preceding cap-valid tangent interval with pilot failure budget
\(\delta\) and total failure budget \(\alpha/2\). Its validation
quantile is \(z_{\mathrm{val}}=z_{1-(\alpha/2-\delta)/2}\). Intersect it
with the direct interval on \(Z\) at error \(\alpha/2\) and the
structural range. All choices are valid on the unseparated class. The
total reported length is bounded by \(2z_{1-\alpha/4}\tau\), including
on pilot failure.

A sufficient condition for the pointwise rate is

\[a\ge 4\eta,\qquad a\ge \eta^2\sqrt p/\tau.\]

For fixed order/confidence this is satisfied at sufficiently large
constant multiple of \(\tau p^{K-1/2}\). This requirement is
conservative and much stronger than a claim of optimal computational
signal strength.

\textbf{Upper proof.} On the simultaneous pilot event, the top-\(r\)
projector \(Q_k\) satisfies \(\|Q_k-P_k\|_{op}\le d=4\eta/a\), where
\(P_k\) projects onto the coordinates in \(S\). The cap-valid
certificate has each angle bound at most \(d\) by Weyl's inequality and
\(a\ge4\eta\). Write \(T_Q\) and \(T_P\) for their first-order
product-space projectors. The representation

\[T_Q=\sum_k\bigotimes_{l\ne k}Q_l-(K-1)\bigotimes_l Q_l\]

with an identity in the omitted mode, and product telescoping, gives
\(\|T_Q-T_P\|_{op}\le2K(K-1)d\). For the diagonal loading,
\(T_PL=p^{-1/2}\sum_{i\in S}E_i\), so

\[\|T_QL\|_F\le\sqrt{r/p}+2K(K-1)d.\]

The cap-valid bias is bounded by \(M\sqrt{\binom K2}d^2\le
16c_0\sqrt{\binom K2}\sqrt r\,\eta^2/a\). The assumed amplitude makes
this \(O(\tau\sqrt{r/p})\). It also gives
\(d\le4\tau/(\eta\sqrt p)=O(\sqrt{r/p})\), because \(\eta\gtrsim\tau\)
and \(r\ge1\). Thus the selected/intersected interval has length
\(O(\tau\sqrt{r/p})\) on the pilot event. Off that event its length is
\(O(\tau)\), with probability at most
\(\delta=O(r/p)\le O(\sqrt{r/p})\). The expected-length upper bound
follows. The dimension-dependent failure budget is essential to this
argument; fixed failure probability would leave an unjustified
\(O(\tau)\) term.

\textbf{Lower proof.} Put \(V=r^{-1/2}\sum_{i\in S}E_i\) and compare
\(A_0\) with \(A_0+(\tau/2)V\). Both have mode ranks at most \(r\) and
norms at most \(M\). Their full Gaussian total variation is less than
\(1/4\) and their target difference is \((\tau/2)\sqrt{r/p}\). The
common-target interval argument at base \(A_0\) proves the lower bound.
This lower comparison uses the same outer class on which the algorithm
is honest.

\textbf{Computation.} Economy SVDs of the \(p\times p^{K-1}\) unfoldings
require \(O(Kp^{K+1})\) arithmetic and \(O(p^K+p^2)\) working storage
for fixed \(K\). Tensor projection has polynomial cost of the same
order. The new code uses economy SVDs; the previous prototype requested
full right singular matrices, which could allocate \(p^{2K-2}\) entries.
Thus the implementation now scales with the explicit tensor rather than
that avoidable allocation. It still reads the full tensor and does not
claim sublinear access cost.

\subsection{What these results do and do not
resolve}\label{what-these-results-do-and-do-not-resolve}

\begin{itemize}
\tightlist
\item
  Efficient high-dimensional minimax computation is resolved here for
  \(r\gtrsim\sqrt p\), where G1 gives the sharp all-budget rate.
  Original fixed-cap diffuse-loading weak signals remain computationally
  open.
\item
  All-budget loading rates are resolved for fixed polynomial decay and
  the specified nonorthogonal perturbations. General profiles remain
  open; the matching W2 statistical construction still uses ideal block
  fits.
\item
  General growing-cap adaptation is impossible in the G2 comparison,
  even between positive cap classes. P1 gives a positive, efficiently
  attained pointwise rate on a strong saturated Tucker family. It does
  not characterize pointwise lengths everywhere or optimal SNR
  thresholds.
\end{itemize}

No blanket tensor hardness theorem is inferred from low-degree results
for sparse tensor PCA. Those results concern a restricted computational
model; a testing reduction for arbitrary efficient intervals still needs
an explicit, justified computational hypothesis.

\subsection{Methodological references}\label{methodological-references}

\begin{itemize}
\tightlist
\item
  The sparse-functional precursor for G1 is Collier, Comminges and
  Tsybakov, \href{https://arxiv.org/abs/1502.00665}{\emph{Minimax
  estimation of linear and quadratic functionals on sparsity classes}}.
\item
  For the distinct estimation-adaptation problem, see Collier,
  Comminges, Tsybakov and Verzélen,
  \href{https://arxiv.org/abs/1611.09744}{\emph{Optimal adaptive
  estimation of linear functionals under sparsity}}.
\item
  The restricted computational perspective in W3 follows the methodology
  of Choo and d'Orsi, \href{https://arxiv.org/abs/2106.06308}{\emph{The
  Complexity of Sparse Tensor PCA}}.
\item
  For general tensor-functional tangent inference under signal
  conditions, see Xu, Chen and Han,
  \href{https://arxiv.org/abs/2501.16223}{\emph{Statistical Inference
  for Low-Rank Tensor Models}}.
\end{itemize}

The repository-specific dependency and novelty distinctions are recorded
in the accompanying audit and literature update.


\section{R1. Scalar companion and risk
proof}\label{r1.-scalar-companion-and-risk-proof}\label{sup:R1}

This is a new statement for the revision, with the complete argument
below. It does not alter the confidence interval or assert pointwise
minimax squared risk.

Let \(Y=\psi+\tau g\), \(g\sim N(0,1)\), and suppose the supplied
structural range is \(|\psi|\le h\). Let \(I\) be the intersection of
any candidate interval with \([Y-d,Y+d]\) and \([-h,h]\), where
\(d=z_{1-\alpha/4}\tau\). The candidate may depend arbitrarily on the
data. The reported confidence interval may retain its existing
empty-intersection convention \(\{0\}\).

Define a \textbf{separate scalar estimator} by \(\widehat\psi=0\) when
\(h\le\tau\). Otherwise, use the midpoint of the actual nonempty
intersection; on an empty intersection use \(\Pi_{[-h,h]}Y\). Replacing
an exactly zero singleton by the clipped direct estimate also satisfies
the bound. Then, for every admissible mean,

\[E(\widehat\psi-\psi)^2\le
\begin{cases}
h^2,&h\le\tau,\\
\tau^2\{1+2z_{1-\alpha/4}\sqrt{2/\pi}+z_{1-\alpha/4}^2\},&h>\tau.
\end{cases}\]

Thus the risk is at most \(C_\alpha\min(h^2,\tau^2)\), uniformly over
the norm budget and signal magnitude. No candidate coverage, separation,
rank recovery, or independence from \(Y\) is needed for this
\textbf{risk} bound. Honesty of the reported interval continues to
require its original assumptions and calibration.

\textbf{Proof.} On \(h\le\tau\), zero has squared error at most \(h^2\).
Otherwise, on a nonempty intersection its midpoint \(m\) is in
\([Y-d,Y+d]\), so \(|m-\psi|\le|Y-\psi|+d\). On an empty intersection,
projection onto an interval containing \(\psi\) is nonexpansive relative
to \(\psi\): \(|\Pi_{[-h,h]}Y-\psi|\le|Y-\psi|\). The same argument
covers a zero-singleton replacement. Consequently the error on every
observation is at most \(\tau|g|+d\). Taking expectations, using
\(Eg^2=1\) and \(E|g|=\sqrt{2/\pi}\), proves the second branch. Since
\(C_\alpha\ge1\), both branches imply the final inequality. No
conditioning on a successful pilot or deletion of rare events occurs.

For A4/P1 and the normalized diagonal loading, \(h=M\sqrt{r/p}\). The
cap-valid intervals already intersect the direct interval at error
\(\alpha/2\), so \(d=z_{1-\alpha/4}\tau\) is exactly their existing
direct radius. The scalar companion therefore removes the \(M^2\) loss
from the arbitrary zero sentinel without changing any coverage event or
interval length. For general fixed \(L\), replace \(\tau\) by
\(\tau\|L\|_F\) and use any valid supplied structural range. This is a
noise-scale risk bound, not a proof of the smaller P1 pointwise
squared-risk rate or of the ideal full-data minimax risk.

\bibliographystyle{plainnat}
\bibliography{references}
\end{document}
